% Les principes de la diplomatie camerounaise — ECM Tle Tech — Cours signé OnBuch+
% Programme officiel MINESEC. Compiler avec : tectonic N01-principes-diplomatie-camerounaise.tex
\newcommand{\DOCMATIERE}{ECM}
\newcommand{\DOCNIVEAU}{Tle Tech}
\newcommand{\DOCMODULE}{ECM – classe de Terminale technique}
\newcommand{\DOCLECON}{N01}
\newcommand{\DOCTITRE}{Les principes de la diplomatie camerounaise}
% OnBuch+ — préambule « cours signés » ENRICHI (compilé avec tectonic / XeLaTeX)
% Système visuel « Pop » d'OnBuch+ : fond crème (jamais blanc), bordures encre
% épaisses, ombres « dures » décalées, titres Archivo Black en capitales.
% Version MATHÉMATIQUES : graphes (pgfplots/tikz), boîtes pédagogiques
% (définition, propriété, méthode, attention, le savais-tu, exercice résolu,
% exercice, TP, bilan), page de garde et signature OnBuch+ ancrée en bas.
%
% Macros attendues dans main.tex AVANT \input{preamble} :
%   \DOCMATIERE  \DOCNIVEAU  \DOCMODULE  \DOCLECON (numéro)  \DOCTITRE
\documentclass[11pt]{article}

\usepackage{fontspec}
\usepackage[french]{babel}
\usepackage[a4paper,margin=2cm,top=3.4cm,bottom=2.8cm,headheight=1.4cm,footskip=1.3cm]{geometry}
\usepackage{amsmath,amssymb,amsfonts}
\usepackage{mathtools} % \xmapsto, \coloneqq... souvent utilisés par les modèles
\usepackage{cancel} % simplifications barrées (bilans de forces)
% \abs{z} (module, valeur absolue) et \norm{u} (norme), utilisés par les modèles
\providecommand{\abs}{}\let\abs\relax\DeclarePairedDelimiter{\abs}{\lvert}{\rvert}
\providecommand{\norm}{}\let\norm\relax\DeclarePairedDelimiter{\norm}{\lVert}{\rVert}
\usepackage[table]{xcolor} % [table] : fournit aussi \rowcolors (alternance de lignes)
\usepackage{pagecolor}
\usepackage{tikz}
\usetikzlibrary{arrows.meta,positioning,calc,shapes.geometric,shapes.misc,shapes.arrows,decorations.pathmorphing,decorations.pathreplacing,decorations.markings,patterns,shadows,mindmap,trees,fit,backgrounds,matrix,intersections,angles,quotes}
\usepackage{pgfplots}
\usepackage[version=4]{mhchem} % équations nucléaires \ce{^{235}_{92}U}
\usepackage[european,siunitx]{circuitikz} % schémas électriques
\pgfplotsset{compat=1.18}
\usepgfplotslibrary{fillbetween,groupplots} % aires entre courbes (calcul intégral)
\usepackage{enumitem}
\usepackage{fancyhdr}
% [most] charge toutes les bibliothèques tcolorbox (listings, minted, raster,
% documentation...) et gonfle inutilement la mémoire TeX sur un document long
% (~300 boîtes) : on ne charge que ce qui est réellement utilisé.
\usepackage[skins,breakable]{tcolorbox}
\usepackage{graphicx}
\usepackage{colortbl}
\usepackage{booktabs}
\usepackage{tabularx}
\usepackage{multirow}
\usepackage{diagbox} % case d'en-tête diagonale des tableaux à double entrée
% Colonnes extensibles centrée / gauche / droite (C, L, R), souvent utilisées par les modèles
\newcolumntype{C}{>{\centering\arraybackslash}X}
\newcolumntype{L}{>{\raggedright\arraybackslash}X}
\newcolumntype{R}{>{\raggedleft\arraybackslash}X}
\usepackage{array}
\usepackage{multicol}
\usepackage{siunitx}
% parse-numbers=false : accepte aussi \SI{2,5 \times 10^{-3}}{...}
\sisetup{locale=FR,output-decimal-marker={,},group-minimum-digits=4,parse-numbers=false}
\DeclareSIUnit\ohm{\ensuremath{\symup{\Omega}}} % U+2126 absent de la police
\DeclareSIUnit\division{div} % réglages d'oscilloscope (ms/div, V/div)
\DeclareSIUnit\tr{tr} % tours (vitesse de rotation, ex. \SI{1500}{\tr\per\minute})
\DeclareSIUnit\molar{M} % concentration molaire (ex. \SI{3}{\molar}, \SI{1}{\micro\molar})
\DeclareSIUnit\an{an} % année (le modèle confond parfois avec \year, un compteur TeX) — ex. \SI{5}{\kilo\meter\per\an}
\DeclareSIUnit\FCFA{FCFA} % franc CFA (ex. \SI{500}{\FCFA\per\kilo\gram})
\DeclareSIUnit\degreeBrix{\textdegree Brix} % teneur en sucre d'un sirop/jus (réfractométrie)
\DeclareSIUnit\month{mois} % le modèle utilise parfois \month, un compteur TeX, comme unité de durée
\DeclareSIUnit\microliter{\micro\liter} % le modèle écrit parfois \microliter en un seul token
\DeclareSIUnit\million{millions} % dénombrement (ex. \SI{15}{\million\per\mL} spermatozoïdes)
\usepackage{qrcode}
% \degree : commande fréquemment utilisée par les modèles pour un angle ou une
% température en dehors de \SI{}{} (non fournie par les paquets chargés).
\providecommand{\degree}{^{\circ}}
% \qed (fin de démonstration), utilisé par les modèles sans amsthm
\providecommand{\qed}{\hfill$\square$}
% \barre : double barre des valeurs interdites dans un tableau de variations (tkz-tab)
\providecommand{\barre}{\ensuremath{\|}}
% environnement proof (amsthm non chargé) utilisé par les modèles
\ifdefined\proof\else\newenvironment{proof}[1][Démonstration]{\par\noindent\textit{#1.}\ }{\qed\par}\fi
\usepackage{lastpage}
\usepackage{hyperref}
\hypersetup{hidelinks}

% ============================================================
% Palette « Pop » OnBuch+ — mêmes hex que la classe P (plus_ui.dart)
% ============================================================
\definecolor{poporange}{HTML}{F59321}
\definecolor{popdark}{HTML}{E07A0C}
\definecolor{popcream}{HTML}{FFF6E8}
\definecolor{popink}{HTML}{1C1714}
\definecolor{poppurple}{HTML}{7A5AE0}
\definecolor{popgreen}{HTML}{2E9E6A}
\definecolor{poppink}{HTML}{E0457B}
\definecolor{popblue}{HTML}{2F7DE1}
\definecolor{popgold}{HTML}{C98A12}
\definecolor{popmuted}{HTML}{8A7F76}
\definecolor{popline}{HTML}{EADFD2}
\definecolor{popgray}{HTML}{8A7F76}
\colorlet{popgrey}{popgray}
% Alias tolérants (le modèle nomme parfois les couleurs différemment)
\colorlet{popwhite}{white}
\colorlet{popblack}{popink}
\colorlet{popred}{poppink}
\colorlet{popyellow}{popgold}
% Teintes claires (fonds de tableaux/encadrés) — le modèle nomme parfois
% les couleurs avec un suffixe L (ex. popblueL) sans les avoir définies.
\colorlet{poporangeL}{poporange!20}
\colorlet{popdarkL}{popdark!20}
\colorlet{poppurpleL}{poppurple!20}
\colorlet{popgreenL}{popgreen!20}
\colorlet{poppinkL}{poppink!20}
\colorlet{popblueL}{popblue!20}
\colorlet{popgoldL}{popgold!20}
\colorlet{popredL}{popred!20}
\colorlet{popgrayL}{popgray!20}
% Teintes claires pour fonds de tableaux / zones
\colorlet{popblueL}{popblue!10!white}
\colorlet{poporangeL}{poporange!14!white}
\colorlet{popgreenL}{popgreen!12!white}
\colorlet{poppurpleL}{poppurple!12!white}
\colorlet{poppinkL}{poppink!12!white}
\colorlet{popgoldL}{popgold!14!white}

\pagecolor{popcream}

% ============================================================
% Typographie
% ============================================================
\defaultfontfeatures{Ligatures=TeX}
\setmainfont{PlusJakartaSans-Medium.ttf}[Path=../../../fonts/,BoldFont=PlusJakartaSans-Bold.ttf]
% Maths en sans-serif (Fira Math) avec chiffres et lettres droites de Plus
% Jakarta, pour que les formules se fondent dans le texte.
\usepackage[math-style=ISO,bold-style=ISO]{unicode-math}
\setmathfont{FiraMath-Regular.otf}[Path=../../../fonts/]
\setmathfont{PlusJakartaSans-Medium.ttf}[Path=../../../fonts/,range={up/num,up/{Latin,latin},\mathup}]
\usepackage{icomma} % virgule décimale sans espace en mode math
% Caractères mathématiques Unicode absents des polices de texte (Plus Jakarta,
% Archivo Black) : rendus par leur équivalent mathématique, en texte comme en
% formule — sinon ils s'affichent en carré vide (ex. « Arithmétique dans ℤ »).
\usepackage{newunicodechar}
\newunicodechar{ℝ}{\ensuremath{\mathbb{R}}}
\newunicodechar{ℤ}{\ensuremath{\mathbb{Z}}}
\newunicodechar{ℂ}{\ensuremath{\mathbb{C}}}
\newunicodechar{ℕ}{\ensuremath{\mathbb{N}}}
\newunicodechar{ℚ}{\ensuremath{\mathbb{Q}}}
\newunicodechar{𝔻}{\ensuremath{\mathbb{D}}}
\newunicodechar{α}{\ensuremath{\alpha}}
\newunicodechar{β}{\ensuremath{\beta}}
\newunicodechar{γ}{\ensuremath{\gamma}}
\newunicodechar{δ}{\ensuremath{\delta}}
\newunicodechar{ε}{\ensuremath{\varepsilon}}
\newunicodechar{θ}{\ensuremath{\theta}}
\newunicodechar{λ}{\ensuremath{\lambda}}
\newunicodechar{μ}{\ensuremath{\mu}}
\newunicodechar{σ}{\ensuremath{\sigma}}
\newunicodechar{τ}{\ensuremath{\tau}}
\newunicodechar{φ}{\ensuremath{\varphi}}
\newunicodechar{ψ}{\ensuremath{\psi}}
\newunicodechar{ω}{\ensuremath{\omega}}
\newunicodechar{Σ}{\ensuremath{\Sigma}}
% majuscules produites par \MakeUppercase dans les titres (α-aminés -> Α)
\newunicodechar{Α}{\ensuremath{\alpha}}
\newunicodechar{Β}{\ensuremath{\beta}}
\newunicodechar{Γ}{\ensuremath{\Gamma}}
\newunicodechar{Δ}{\ensuremath{\Delta}}
\newunicodechar{Ε}{\ensuremath{\varepsilon}}
\newunicodechar{Θ}{\ensuremath{\Theta}}
\newunicodechar{Λ}{\ensuremath{\Lambda}}
\newunicodechar{Μ}{\ensuremath{\mu}}
\newunicodechar{Τ}{\ensuremath{\tau}}
\newunicodechar{Φ}{\ensuremath{\Phi}}
\newunicodechar{Ψ}{\ensuremath{\Psi}}
\newunicodechar{Ω}{\ensuremath{\Omega}}
\newunicodechar{↦}{\ensuremath{\mapsto}}
\newunicodechar{∈}{\ensuremath{\in}}
\newunicodechar{∉}{\ensuremath{\notin}}
\newunicodechar{∧}{\ensuremath{\wedge}}
\newunicodechar{⇔}{\ensuremath{\Leftrightarrow}}
\newunicodechar{⟺}{\ensuremath{\Longleftrightarrow}}
\newunicodechar{⇒}{\ensuremath{\Rightarrow}}
\newunicodechar{⟹}{\ensuremath{\Longrightarrow}}
\newunicodechar{∘}{\ensuremath{\circ}}
\newunicodechar{∪}{\ensuremath{\cup}}
\newunicodechar{∩}{\ensuremath{\cap}}
\newunicodechar{≡}{\ensuremath{\equiv}}
\newunicodechar{⊂}{\ensuremath{\subset}}
\newunicodechar{⊥}{\ensuremath{\perp}}
\newunicodechar{∃}{\ensuremath{\exists}}
\newunicodechar{∀}{\ensuremath{\forall}}
\newunicodechar{∛}{\ensuremath{\sqrt[3]{\,}}}
\newunicodechar{∜}{\ensuremath{\sqrt[4]{\,}}}
\newunicodechar{⃗}{\ensuremath{{}^{\to}}}
\newunicodechar{ⁿ}{\ifmmode^{n}\else\textsuperscript{\ensuremath{n}}\fi}
\newunicodechar{ˣ}{\ifmmode^{x}\else\textsuperscript{\ensuremath{x}}\fi}
\newunicodechar{ᵇ}{\ifmmode^{b}\else\textsuperscript{\ensuremath{b}}\fi}
\newunicodechar{ʳ}{\ifmmode^{r}\else\textsuperscript{\ensuremath{r}}\fi}
\newunicodechar{ᵘ}{\ifmmode^{u}\else\textsuperscript{\ensuremath{u}}\fi}
\newunicodechar{ʸ}{\ifmmode^{y}\else\textsuperscript{\ensuremath{y}}\fi}
\newunicodechar{ᵃ}{\ifmmode^{a}\else\textsuperscript{\ensuremath{a}}\fi}
\newunicodechar{ᵉ}{\ifmmode^{e}\else\textsuperscript{\ensuremath{e}}\fi}
\newunicodechar{ᵏ}{\ifmmode^{k}\else\textsuperscript{\ensuremath{k}}\fi}
\newunicodechar{ᵖ}{\ifmmode^{p}\else\textsuperscript{\ensuremath{p}}\fi}
\newunicodechar{ᴺ}{\ifmmode^{N}\else\textsuperscript{\ensuremath{N}}\fi}
\newunicodechar{ᶜ}{\ifmmode^{c}\else\textsuperscript{\ensuremath{c}}\fi}
\newunicodechar{ʰ}{\ifmmode^{h}\else\textsuperscript{\ensuremath{h}}\fi}
\newunicodechar{ᵠ}{\ifmmode^{q}\else\textsuperscript{\ensuremath{q}}\fi}
\newunicodechar{ₙ}{\ifmmode_{n}\else\textsubscript{\ensuremath{n}}\fi}
\newunicodechar{ₐ}{\ifmmode_{a}\else\textsubscript{\ensuremath{a}}\fi}
\newunicodechar{ᵢ}{\ifmmode_{i}\else\textsubscript{\ensuremath{i}}\fi}
\newunicodechar{ᵣ}{\ifmmode_{r}\else\textsubscript{\ensuremath{r}}\fi}
\newunicodechar{ₖ}{\ifmmode_{k}\else\textsubscript{\ensuremath{k}}\fi}
\newunicodechar{ᵦ}{\ifmmode_{\beta}\else\textsubscript{\ensuremath{\beta}}\fi}
\newunicodechar{ₓ}{\ifmmode_{x}\else\textsubscript{\ensuremath{x}}\fi}
\newfontfamily\popdisplay{ArchivoBlack-Regular.ttf}[Path=../../../fonts/]
\newfontfamily\poplogo{SpaceGrotesk-Bold.ttf}[Path=../../../fonts/]
\newfontfamily\popsemi{PlusJakartaSans-SemiBold.ttf}[Path=../../../fonts/]
\newfontfamily\popxbold{PlusJakartaSans-ExtraBold.ttf}[Path=../../../fonts/]
\newfontfamily\popxboldls{PlusJakartaSans-ExtraBold.ttf}[Path=../../../fonts/,LetterSpace=3.0]

\setlist[enumerate]{leftmargin=1.7em,itemsep=5pt,topsep=4pt}
\setlist[itemize]{leftmargin=1.7em,itemsep=4pt,topsep=4pt,label={\color{poporange}\textbullet}}
\setlength{\parindent}{0pt}
\setlength{\parskip}{5pt}
\renewcommand{\arraystretch}{1.25}

% pgfplots : style Pop par défaut
\pgfplotsset{
  every axis/.append style={axis line style={popink,line width=1pt},tick style={popink},
    grid style={popline,line width=0.5pt},label style={font=\small},tick label style={font=\footnotesize}},
  popcurve/.style={poporange,line width=1.8pt,no markers,smooth},
}
% Même style disponible dans un \draw TikZ simple (hors pgfplots)
\tikzset{popcurve/.style={poporange,line width=1.8pt}}
\tikzset{popgrid/.style={popline,very thin}}
\colorlet{popcurve}{poporange} % le nom du style est parfois utilisé comme couleur

% ============================================================
% Pill / squiggle
% ============================================================
\newcommand{\poppill}[3]{%
  \tikz[baseline=-0.55ex]\node[draw=popink,line width=1.2pt,rounded corners=2.6pt,%
    fill=#2,inner xsep=6.5pt,inner ysep=3pt]{%
    \popxboldls\fontsize{7}{7}\selectfont\color{#3}\MakeUppercase{#1}};%
}
\newcommand{\popsquiggle}[1]{%
  \tikz[baseline=-0.4ex]{
    \draw[#1,line width=2.6pt,line cap=round]
      plot [smooth] coordinates {(0,0.09) (0.34,0.01) (0.68,0.21) (1.02,0.01) (1.36,0.10)};
  }%
}
% Mot-clé surligné (marqueur orange sous le texte)
\newcommand{\cle}[1]{{\popxbold\color{popdark}#1}}

% ============================================================
% En-tête / pied
% ============================================================
\pagestyle{fancy}
\fancyhf{}
\renewcommand{\headrulewidth}{0pt}
\renewcommand{\footrulewidth}{0pt}
\lhead{\raisebox{-0.15ex}{\poplogo\fontsize{13}{13}\selectfont\color{popink}OnBuch\color{poporange}+}}
\rhead{\poppill{\DOCMATIERE\ \textbullet\ \DOCNIVEAU\ \textbullet\ Leçon \DOCLECON}{white}{popink}}
\lfoot{\popxbold\fontsize{7.6}{7.6}\selectfont\color{popmuted}\MakeUppercase{Cours signé OnBuch+}\ \textbullet\ {\popsemi\DOCTITRE}}
\rfoot{\tikz[baseline=-0.6ex]\node[rounded corners=8.5pt,fill=popink,minimum height=17pt,minimum width=17pt,inner xsep=5pt,inner sep=0]{\popxbold\fontsize{8}{8}\selectfont\color{popcream}\thepage\,/\,\pageref*{LastPage}};}

% ============================================================
% Titres
% ============================================================
\newcommand{\chaptitre}[2]{%
  \begin{tcolorbox}[enhanced,colback=poporange,colframe=popink,boxrule=2.6pt,arc=11pt,
      drop shadow=popink,left=15pt,right=15pt,top=12pt,bottom=11pt,before skip=3pt,after skip=18pt]
    \poppill{#2}{popink}{popcream}\par\vspace{9pt}
    {\raggedright\hyphenpenalty=10000\exhyphenpenalty=10000\popdisplay\fontsize{22}{24}\selectfont\color{popcream}\MakeUppercase{#1}\par}
  \end{tcolorbox}
}
% Section numérotée : pastille orange avec le numéro + titre Archivo
\newcounter{popsec}
\newcommand{\coursec}[1]{%
  \stepcounter{popsec}\setcounter{popsub}{0}%
  \par\vspace{16pt}\noindent
  \begin{minipage}{\linewidth}
  \tikz[baseline=-0.7ex]\node[draw=popink,line width=1.6pt,fill=poporange,rounded corners=4pt,minimum size=22pt,inner sep=0]{\popdisplay\fontsize{12}{12}\selectfont\color{popcream}\thepopsec};%
  \hspace{8pt}{\popdisplay\fontsize{15}{17}\selectfont\color{popink}\MakeUppercase{#1}}\par
  \vspace{4pt}\noindent{\color{poporange}\rule{36pt}{3.6pt}}
  \end{minipage}\par\nopagebreak\vspace{8pt}
}
\newcounter{popsub}
\newcommand{\courssub}[1]{%
  \stepcounter{popsub}%
  \par\vspace{9pt}\noindent{\popxbold\fontsize{11.5}{13}\selectfont\color{popink}{\color{poporange}\thepopsec.\thepopsub}\hspace{6pt}#1}\par\nopagebreak\vspace{4pt}
}

% ============================================================
% Boîtes pédagogiques — même recette : fond blanc, bordure épaisse colorée,
% ombre dure encre, badge plein. Titre optionnel [#1].
% ============================================================
\tcbset{popbase/.style={breakable,enhanced,colback=white,boxrule=2.3pt,arc=9pt,
  shadow={3pt}{-3pt}{0pt}{popink},left=14pt,right=14pt,top=11pt,bottom=11pt,before skip=11pt,after skip=15pt,
  fontupper=\popsemi\fontsize{10.6}{14.5}\selectfont\color{popink},
  fontlower=\popsemi\fontsize{10.6}{14.5}\selectfont\color{popink}}}
% Coche dessinée (le glyphe ✓ n'existe pas dans les polices du document)
\newcommand{\popcheck}[1]{\tikz[baseline=-0.5ex]\draw[#1,line width=1.6pt,line cap=round,line join=round](-0.13,0.0)--(-0.04,-0.09)--(0.13,0.1);}
\providecommand{\coche}{\popcheck{popgreen}}
\providecommand{\resultat}[1]{\par\smallskip\noindent\fcolorbox{popgreen}{popgreenL}{\parbox{\dimexpr\linewidth-2\fboxsep-2\fboxrule\relax}{\textbf{Résultat.} #1}}\par\smallskip}
\newcommand{\popboxhead}[3]{\poppill{#1}{#2}{white}\ifx\relax#3\relax\else\hspace{6pt}{\popxbold\fontsize{10.6}{12}\selectfont\color{popink}#3}\fi\par\vspace{7pt}}

\newtcolorbox{aretenir}[1][]{popbase,colframe=popgreen,before upper={\popboxhead{À retenir}{popgreen}{#1}}}
\newtcolorbox{exemplebox}[1][]{popbase,colframe=poporange,before upper={\popboxhead{Exemple}{poporange}{#1}}}
\newtcolorbox{definition}[1][]{popbase,colframe=popblue,before upper={\popboxhead{Définition}{popblue}{#1}}}
\newtcolorbox{propriete}[1][]{popbase,colframe=poppurple,before upper={\popboxhead{Propriété}{poppurple}{#1}}}
\newtcolorbox{methode}[1][]{popbase,colframe=popgold,before upper={\popboxhead{Méthode}{popgold}{#1}}}
\newtcolorbox{attention}[1][]{popbase,colframe=poppink,before upper={\popboxhead{Attention}{poppink}{#1}}}
\newtcolorbox{experience}[1][]{popbase,colframe=popblue,colback=popblueL,before upper={\popboxhead{Expérience}{popblue}{#1}}}
% « Le savais-tu ? » : carte violette pleine (contexte, culture, Cameroun)
\newtcolorbox{savaistu}[1][]{popbase,colback=poppurple,colframe=popink,
  fontupper=\popsemi\fontsize{10.4}{14.2}\selectfont\color{white},
  before upper={\poppill{Le savais-tu\,?}{white}{poppurple}\ifx\relax#1\relax\else\hspace{6pt}{\popxbold\color{white}#1}\fi\par\vspace{7pt}}}
% Exercice résolu : énoncé en haut, \tcblower puis la solution sur fond crème
\newcounter{popexo}\newcounter{popexr}
\newtcolorbox{exoresolu}[1][]{popbase,colframe=poporange,colbacklower=popcream,
  segmentation style={popink,line width=1.2pt,dashed},
  before upper={\stepcounter{popexr}\popboxhead{Exercice résolu \thepopexr}{poporange}{#1}},
  before lower={\poppill{Solution}{popink}{popcream}\par\vspace{6pt}}}
% Exercice (non résolu en place) — la correction est dans la partie Corrigés
\newtcolorbox{exercice}[1][]{popbase,colframe=popink,
  before upper={\stepcounter{popexo}\popboxhead{Exercice \thepopexo}{popink}{#1}}}
% Corrigé
\newtcolorbox{corrige}[1][]{popbase,colframe=popgreen,colback=popgreenL,
  before upper={\popboxhead{Corrigé}{popgreen}{#1}}}
% Objectifs / prérequis / bilan
\newtcolorbox{objectifs}{popbase,colframe=popink,colback=poporangeL,
  before upper={\popboxhead{Objectifs de la leçon}{popink}{}}}
\newtcolorbox{prerequis}{popbase,colframe=popink,colback=popblueL,
  before upper={\popboxhead{Prérequis}{popblue}{}}}
\newtcolorbox{bilan}{popbase,colframe=popink,colback=white,boxrule=2.8pt,
  before upper={\popboxhead{Fiche bilan}{poporange}{L'essentiel à connaître pour l'examen}}}
% Figure encadrée avec légende
\newtcolorbox{popfigure}[1][]{enhanced,breakable=false,colback=white,colframe=popline,boxrule=1.4pt,arc=8pt,
  left=10pt,right=10pt,top=10pt,bottom=8pt,before skip=10pt,after skip=12pt,halign=center,
  lower separated=false,halign lower=center,
  fontlower=\popsemi\fontsize{9}{11.5}\selectfont\color{popmuted},
  #1}
\newcommand{\legende}[1]{\par\vspace{6pt}{\popsemi\fontsize{9}{11.5}\selectfont\color{popmuted}#1}}

% ============================================================
% Signature OnBuch+
% ============================================================
\newcommand{\poplogoinline}[1]{{\poplogo\fontsize{#1}{#1}\selectfont\color{popink}OnBuch\color{poporange}+}}
\newcommand{\signatureonbuch}{%
  \begin{tcolorbox}[enhanced,colback=popink,colframe=popink,boxrule=2.2pt,arc=10pt,
      drop shadow=poporange,left=14pt,right=14pt,top=10pt,bottom=10pt,before skip=0pt,after skip=0pt]
    \tikz[baseline=-0.7ex]\node[circle,fill=popgreen,draw=popcream,line width=1.3pt,minimum size=22pt,inner sep=0]{\popcheck{white}};%
    \hspace{9pt}\begin{minipage}[c]{0.62\linewidth}
      {\popxbold\fontsize{12}{13}\selectfont\color{popcream}Signé\ \textbullet\ {\poplogo OnBuch\color{poporange}+}}\par\vspace{2pt}
      {\raggedright\popsemi\fontsize{8}{10.5}\selectfont\color{popline}\MakeUppercase{Équipe pédagogique OnBuch+}\\\MakeUppercase{Conforme au programme officiel MINESEC}\par}
    \end{minipage}\hfill
    {\poplogo\fontsize{22}{22}\selectfont\color{popcream}OnBuch\color{poporange}+}
  \end{tcolorbox}%
}

% ============================================================
% Page de garde
% #1 titre · #2 matière · #3 niveau · #4 module · #5 n° leçon · #6 durée ·
% #7 description / objectifs (texte ou itemize)
% ============================================================
\newcommand{\pagegarde}[7]{%
  \thispagestyle{empty}
  \newgeometry{a4paper,margin=1.7cm,top=1.6cm,bottom=1.4cm}%
  \begin{tikzpicture}[remember picture,overlay]
    \fill[poporange!18] ([xshift=-3.2cm,yshift=-3.6cm]current page.north east) circle (4.6cm);
    \fill[poppurple!14] ([xshift=2.4cm,yshift=7.6cm]current page.south west) circle (3.8cm);
  \end{tikzpicture}%
  {\poplogo\fontsize{30}{30}\selectfont\color{popink}OnBuch\color{poporange}+}\hfill
  \raisebox{0.6ex}{\poppill{Cours signé \textbullet\ Premium}{popink}{popcream}}\par
  \vspace{1pt}\popsquiggle{poppurple}\par
  \vspace{8pt}
  \begin{tcolorbox}[enhanced,colback=poporange,colframe=popink,boxrule=2.8pt,arc=13pt,
      drop shadow=popink,left=18pt,right=18pt,top=12pt,bottom=13pt,after skip=10pt]
    \poppill{#2 \textbullet\ #3}{popink}{popcream}\hspace{5pt}\poppill{#4}{white}{popink}\par\vspace{8pt}
    {\popdisplay\fontsize{14}{14}\selectfont\color{popink}LEÇON #5}\par\vspace{4pt}
    {\raggedright\hyphenpenalty=10000\exhyphenpenalty=10000\popdisplay\fontsize{25}{27}\selectfont\color{popcream}\MakeUppercase{#1}\par}
    \vspace{8pt}
    {\popxbold\fontsize{9.5}{9.5}\selectfont\color{popink}\MakeUppercase{Durée conseillée : #6}}
  \end{tcolorbox}
  \begin{tcolorbox}[enhanced,colback=white,colframe=popink,boxrule=2.2pt,arc=10pt,
      drop shadow=popink,left=16pt,right=16pt,top=10pt,bottom=10pt,after skip=8pt]
    \poppill{Dans cette leçon}{poppurple}{white}\par\vspace{5pt}
    \popsemi\fontsize{9.3}{12.6}\selectfont\color{popink}#7
  \end{tcolorbox}
  \noindent\begin{minipage}[t]{0.487\linewidth}
    \begin{tcolorbox}[enhanced,colback=popgreen,colframe=popink,boxrule=2.2pt,arc=10pt,
        drop shadow=popink,left=11pt,right=11pt,top=10pt,bottom=10pt,height=3.1cm,valign=center]
      \tcbox[colback=white,colframe=popink,boxrule=1.3pt,arc=5pt,boxsep=0pt,nobeforeafter,
        left=3pt,right=3pt,top=3pt,bottom=3pt,valign=center]{\qrcode[height=1.9cm]{https://play.google.com/store/apps/details?id=com.luvvix.onbuch&hl=fr}}%
      \hspace{8pt}\begin{minipage}[c]{0.5\linewidth}
        {\popdisplay\fontsize{11}{12.5}\selectfont\color{white}\MakeUppercase{Télécharge OnBuch}}\par\vspace{4pt}
        {\raggedright\popsemi\fontsize{8.4}{11}\selectfont\color{white}Gratuit \textbullet\ Google Play\\Tuteur IA, annales, concours, résultats\par}
      \end{minipage}
    \end{tcolorbox}
  \end{minipage}\hfill
  \begin{minipage}[t]{0.487\linewidth}
    \begin{tcolorbox}[enhanced,colback=poppurple,colframe=popink,boxrule=2.2pt,arc=10pt,
        drop shadow=popink,left=11pt,right=11pt,top=10pt,bottom=10pt,height=3.1cm,valign=center]
      \tikz[baseline=-0.7ex]\node[circle,fill=white,draw=popink,line width=1.3pt,minimum size=1.6cm,inner sep=0]{\popdisplay\fontsize{24}{24}\selectfont\color{poppurple}+};%
      \hspace{8pt}\begin{minipage}[c]{0.6\linewidth}
        {\popdisplay\fontsize{11}{12.5}\selectfont\color{white}\MakeUppercase{Passe à OnBuch+}}\par\vspace{4pt}
        {\raggedright\popsemi\fontsize{8.4}{11}\selectfont\color{white}Tous les cours signés, Tuteur IA illimité, fiches PDF — onglet OnBuch+ de l'app\par}
      \end{minipage}
    \end{tcolorbox}
  \end{minipage}
  \vspace{8pt}
  \popcontenu
  \vfill
  \signatureonbuch
  \restoregeometry
  \newpage
}

% Rangée « Ce cours contient » (page de garde)
\newcommand{\poptile}[3]{% #1 couleur #2 gros texte #3 légende
  \begin{tcolorbox}[enhanced,colback=#1,colframe=popink,boxrule=2pt,arc=9pt,shadow={2.5pt}{-2.5pt}{0pt}{popink},
      left=6pt,right=6pt,top=9pt,bottom=9pt,halign=center,valign=center,height=1.75cm,width=\linewidth,nobeforeafter]
    {\popdisplay\fontsize{12.5}{13}\selectfont\color{white}\MakeUppercase{#2}}\par\vspace{3pt}
    {\centering\popsemi\fontsize{7.8}{9.5}\selectfont\color{white}#3\par}
  \end{tcolorbox}}
\newcommand{\popcontenu}{%
  \par\noindent{\popxboldls\fontsize{7.5}{7.5}\selectfont\color{popmuted}CE COURS SIGNÉ CONTIENT}\par\vspace{7pt}
  \noindent\begin{minipage}[t]{0.235\linewidth}\poptile{popblue}{Cours}{complet, illustré, pas à pas}\end{minipage}\hfill
  \begin{minipage}[t]{0.235\linewidth}\poptile{poporange}{Méthodes}{et exercices résolus}\end{minipage}\hfill
  \begin{minipage}[t]{0.235\linewidth}\poptile{poppink}{Exercices}{type examen + corrigés}\end{minipage}\hfill
  \begin{minipage}[t]{0.235\linewidth}\poptile{popgreen}{Fiche bilan}{l'essentiel à retenir}\end{minipage}\par}

% Sommaire
\newtcolorbox{sommaire}{enhanced,colback=white,colframe=popink,boxrule=2.2pt,arc=10pt,
  drop shadow=popink,left=18pt,right=18pt,top=13pt,bottom=13pt,before skip=6pt,after skip=20pt,
  before upper={\poppill{Sommaire}{poppurple}{white}\par\vspace{9pt}\popsemi\fontsize{10.6}{14.5}\selectfont\color{popink}}}

% Signature de fin de cours — poussée tout en bas de la dernière page
\newcommand{\findecours}{%
  \par\vfill\nopagebreak
  \begin{center}\popsquiggle{poporange}\end{center}\vspace{4pt}
  \signatureonbuch
}
\AtBeginDocument{\def\llbracket{\mathopen{[\mkern-3.5mu[}}\def\rrbracket{\mathclose{]\mkern-3.5mu]}}} % intervalles d'entiers (glyphe ⟦ absent de la police)
% Glyphes absents de Fira Math : redéfinis à partir de glyphes disponibles
\AtBeginDocument{%
  \def\setminus{\mathbin{\backslash}}\let\smallsetminus\setminus
  \def\vdots{\mathord{\vbox{\baselineskip3.2pt\lineskiplimit0pt\kern4pt\hbox{.}\hbox{.}\hbox{.}}}}%
  \def\ddots{\mathinner{\mkern1mu\raise6.5pt\hbox{.}\mkern2mu\raise3.6pt\hbox{.}\mkern2mu\raise0.7pt\hbox{.}\mkern1mu}}%
  \def\triangle{\mathord{\scalebox{0.9}{$\symup{\Delta}$}}}%
  \def\nabla{\mathord{\raisebox{\depth}{\scalebox{1}[-1]{$\symup{\Delta}$}}}}%
  \def\checkmark{\popcheck{popdark}}%
  \def\star{\mathbin{\ast}}\def\bigstar{\mathord{\ast}}%
  \def\diamond{\mathbin{\scalebox{0.6}{$\Diamond$}}}%
  \def\bigcup{\mathop{\mathchoice{\raisebox{-0.3em}{\scalebox{1.7}{$\cup$}}}{\scalebox{1.25}{$\cup$}}{\cup}{\cup}}\displaylimits}%
  \def\bigcap{\mathop{\mathchoice{\raisebox{-0.3em}{\scalebox{1.7}{$\cap$}}}{\scalebox{1.25}{$\cap$}}{\cap}{\cap}}\displaylimits}%
  \def\lesseqqgtr{\mathrel{\lessgtr}}\def\bigoplus{\mathop{\oplus}\displaylimits}%
}
\providecommand{\ssi}{si et seulement si} % abréviation fréquente
\AtBeginDocument{\providecommand{\wideparen}{\overparen}} % arc de cercle

%% FIGFIX-BEGIN v5 — correctifs globaux des figures (docs/onbuch-generation/tools/figfix)
\usepackage{adjustbox}\usepackage{etoolbox}\tcbuselibrary{raster}
% toute figure TikZ/pgfplots plus large que la ligne est réduite à la largeur disponible
\BeforeBeginEnvironment{tikzpicture}{\begin{adjustbox}{max width=\linewidth}}
\AfterEndEnvironment{tikzpicture}{\end{adjustbox}}
% légendes pgfplots lisibles par-dessus les courbes
\pgfplotsset{every axis/.append style={legend style={fill=white,fill opacity=0.92,text opacity=1,draw=black!15,inner sep=3pt}}}
% étiquettes d'arêtes (syntaxe edge["..."]) avec fond blanc
\tikzset{every edge quotes/.append style={fill=white,inner sep=1.2pt}}
%% FIGFIX-END
\begin{document}

\pagegarde{Les principes de la diplomatie camerounaise}{ECM}{Tle Tech}{ECM – classe de Terminale technique}{N01}{4 séances (fiche de progression harmonisée)}{Cette leçon présente les relations internationales, les formes de la coopération internationale (dont l'assistance humanitaire) et les principes fondamentaux qui guident la diplomatie du Cameroun sur la scène mondiale. Elle s'appuie sur des documents authentiques (cartes, photographies, données chiffrées) et sur des situations concrètes vécues au Cameroun.
\vspace{4pt}
\begin{multicols}{2}\small
\begin{itemize}
\item À la fin de cette leçon, je sais définir les relations internationales et identifier leurs acteurs
\item Je sais distinguer les aspects (politique, économique, culturel, militaire) et les formes (bilatérale, multilatérale) de la coopération internationale
\item Je sais expliquer ce qu'est l'assistance humanitaire et illustrer son action au Cameroun (réfugiés, catastrophes)
\item Je sais énoncer et expliquer les principes de la diplomatie camerounaise (non-ingérence, coexistence pacifique, non-alignement, règlement pacifique des conflits)
\item Je sais analyser des documents (cartes, photographies, données chiffrées) relatifs à l'action internationale du Cameroun
\item Je sais appliquer ces notions à des situations concrètes de la vie courante (accueil de réfugiés, coopération avec la Chine, maintien de la paix)
\end{itemize}\end{multicols}}

\begin{sommaire}
\begin{multicols}{2}
\begin{enumerate}
\item Le concept de relations internationales
\item Les aspects de la coopération internationale
\item Les formes de la coopération internationale
\item Dossier 1 : L'assistance humanitaire
\item Les principes de la diplomatie camerounaise
\item TP : cartographier et argumenter l'action internationale du Cameroun
\item Activité d'intégration
\item Méthodes et exercices résolus
\item Exercices
\item Corrigés des exercices
\item Fiche bilan
\end{enumerate}
\end{multicols}
\end{sommaire}

\begin{objectifs}
\begin{itemize}
\item À la fin de cette leçon, je sais définir les relations internationales et identifier leurs acteurs
\item Je sais distinguer les aspects (politique, économique, culturel, militaire) et les formes (bilatérale, multilatérale) de la coopération internationale
\item Je sais expliquer ce qu'est l'assistance humanitaire et illustrer son action au Cameroun (réfugiés, catastrophes)
\item Je sais énoncer et expliquer les principes de la diplomatie camerounaise (non-ingérence, coexistence pacifique, non-alignement, règlement pacifique des conflits)
\item Je sais analyser des documents (cartes, photographies, données chiffrées) relatifs à l'action internationale du Cameroun
\item Je sais appliquer ces notions à des situations concrètes de la vie courante (accueil de réfugiés, coopération avec la Chine, maintien de la paix)
\item Je sais construire un croquis ou un diagramme simple présentant les partenaires du Cameroun
\item Je sais rédiger et justifier une démarche, une réponse ou une conclusion argumentée sur un sujet de diplomatie camerounaise
\end{itemize}
\end{objectifs}

\begin{prerequis}
\begin{itemize}
\item La notion d'État et de souveraineté (cours de Première sur l'État camerounais)
\item Les grandes organisations internationales (ONU, Union africaine) vues en Histoire-Géographie
\item La localisation du Cameroun et de ses voisins en Afrique centrale
\item La notion de citoyenneté et de solidarité (ECM, classes antérieures)
\end{itemize}
\end{prerequis}

\newpage

\coursec{Le concept de relations internationales}

\textbf{Situation d'accroche.} En août 2023, un même mois, le Cameroun envoie un contingent de casques bleus en République centrafricaine (MINUSCA) pour le maintien de la paix, accueille des réfugiés nigérians dans le camp de Minawao (département du Mayo-Tsanaga, Extrême-Nord), signe des accords de financement d'infrastructures routières avec la Chine et accueille à Yaoundé un sommet ordinaire de la CEMAC. À première vue, ces actions semblent dispersées : sécurité, humanitaire, économie, politique régionale. Pourtant, elles obéissent toutes à une même logique : la diplomatie camerounaise, c'est-à-dire la manière dont le Cameroun dialogue avec le reste du monde pour défendre ses intérêts et promouvoir ses valeurs.

\textbf{Problématique.} Comment un État comme le Cameroun dialogue-t-il avec le reste du monde ? Sur quels principes repose son action internationale, et comment se manifeste concrètement sa coopération avec les autres États et organisations ? Pour répondre à ces questions, il faut d'abord comprendre ce que sont les relations internationales : c'est l'objet de cette première partie.

\courssub{Définition des relations internationales}

\textbf{Intuition.} Dans la vie quotidienne, un élève de Terminale technique entretient des relations avec ses camarades, ses professeurs, sa famille, son maître de stage. Ces relations peuvent être amicales, commerciales, conflictuelles. De la même manière, un État n'est jamais seul : il échange, négocie, coopère ou s'oppose avec d'autres États et d'autres organisations. Aucun pays, même le plus puissant, ne peut vivre en autarcie totale : le pétrole, les médicaments, les téléphones, les capitaux, les idées circulent au-delà des frontières.

\begin{definition}[Les relations internationales]
Les \cle{relations internationales} sont l'ensemble des liens \cle{politiques}, \cle{économiques}, \cle{culturels} et \cle{militaires} qui unissent les \cle{États} et les autres \cle{acteurs de la scène mondiale} (organisations internationales, ONG, entreprises multinationales, diasporas).
\end{definition}

Cette définition appelle trois remarques importantes.

\begin{itemize}
\item Les relations internationales sont \textbf{plurielles} : elles ne se limitent pas à la politique. Un contrat pétrolier, un match de football, un échange d'étudiants ou une alliance militaire sont aussi des relations internationales.
\item Elles dépassent les seuls \textbf{États} : une ONG comme la Croix-Rouge ou une entreprise comme MTN agissent aussi sur la scène mondiale.
\item Elles peuvent être \textbf{pacifiques} (coopération, échanges, traités) ou \textbf{conflictuelles} (sanctions, ruptures diplomatiques, rivalités, guerres).
\end{itemize}

\begin{attention}[Erreur fréquente : confusion intérieur / extérieur]
Ne pas confondre les \cle{relations internationales} (entre États et acteurs mondiaux, \emph{au-delà des frontières}) et la \cle{politique intérieure} (relations entre les citoyens et les institutions d'un même pays : élections, impôts, justice, décentralisation). Quand le Cameroun signe un accord avec la Chine, c'est une relation internationale ; quand un maire camerounais construit un marché communal, c'est une affaire intérieure. À l'examen, cette confusion fait perdre des points.
\end{attention}

\courssub{Les acteurs des relations internationales}

Qui parle, qui négocie, qui agit sur la scène mondiale ? On distingue cinq grandes catégories d'acteurs.

\textbf{1. Les États.} Ce sont les acteurs principaux et historiques. Un État possède un territoire, une population, un gouvernement et la \cle{souveraineté}, c'est-à-dire le pouvoir de décider librement de sa politique intérieure et extérieure. Le Cameroun, la France, la Chine, le Nigeria sont des États.

\textbf{2. Les organisations internationales.} Ce sont des associations d'États créées par un traité pour coopérer dans un domaine précis. On cite :
\begin{itemize}
\item l'\cle{ONU} (Organisation des Nations Unies, 1945) : maintien de la paix et de la sécurité internationale, développement, droits de l'homme ;
\item l'\cle{UA} (Union Africaine, 2002, héritière de l'OUA) : unité, solidarité et développement de l'Afrique ;
\item la \cle{CEMAC} (Communauté Économique et Monétaire de l'Afrique Centrale) : intégration économique de six pays d'Afrique centrale (Cameroun, RCA, Tchad, Congo, Gabon, Guinée équatoriale), monnaie commune (FCFA).
\end{itemize}

\textbf{3. Les ONG} (Organisations Non Gouvernementales) : associations indépendantes des gouvernements qui interviennent dans l'humanitaire, l'environnement ou les droits humains (Croix-Rouge internationale, Médecins Sans Frontières, Plan International, Amnesty International).

\textbf{4. Les entreprises multinationales} : grandes sociétés présentes dans plusieurs pays (TotalEnergies, MTN, Orange, Bolloré Logistics, Société Générale). Leur poids économique leur donne une influence réelle sur les décisions des États.

\textbf{5. Les diasporas} : communautés de ressortissants d'un pays vivant à l'étranger. La diaspora camerounaise, par ses envois d'argent (transferts de fonds) et ses investissements, est un acteur économique et culturel de premier plan.

\begin{popfigure}
\begin{center}
\begin{tikzpicture}[scale=1.0, every node/.style={font=\small}]
\node[circle,draw=popdark,fill=poporangeL,minimum size=2.6cm,align=center,font=\bfseries\small] (cam) at (0,0) {CAMEROUN};
\node[rectangle,rounded corners,draw=popblue,fill=popblueL,align=center,minimum width=2.8cm] (etats) at (90:3.8) {\textbf{États souverains}\\ France, Chine, Nigeria, USA\ldots};
\node[rectangle,rounded corners,draw=popgreen,fill=popgreenL,align=center,minimum width=2.8cm] (onu) at (40:4.1) {\textbf{Organisations int.}\\ ONU, UA, CEMAC, CEEAC};
\node[rectangle,rounded corners,draw=poppurple,fill=poppurpleL,align=center,minimum width=2.8cm] (ong) at (-40:4.1) {\textbf{ONG}\\ Croix-Rouge, MSF, Plan Int.};
\node[rectangle,rounded corners,draw=popgold,fill=popgoldL,align=center,minimum width=2.8cm] (mult) at (-140:3.8) {\textbf{Multinationales}\\ TotalEnergies, MTN, Orange};
\node[rectangle,rounded corners,draw=poppink,fill=poppinkL,align=center,minimum width=2.8cm] (dias) at (140:3.8) {\textbf{Diasporas}\\ Transferts, investissements, lobbying};
\draw[<->,thick,popblue] (cam) -- (etats);
\draw[<->,thick,popgreen] (cam) -- (onu);
\draw[<->,thick,poppurple] (cam) -- (ong);
\draw[<->,thick,popgold] (cam) -- (mult);
\draw[<->,thick,poppink] (cam) -- (dias);
\end{tikzpicture}
\end{center}
\legende{Schéma en étoile des acteurs des relations internationales du Cameroun. Chaque flèche double symbolise un échange : le Cameroun reçoit (aides, investissements, technologies, appui politique) et donne (casques bleus, matières premières, voix diplomatiques, marché de consommation).}
\end{popfigure}

\begin{savaistu}[Le Cameroun, un carrefour diplomatique original]
Le Cameroun entretient des relations diplomatiques avec plus de \textbf{180 États} (sur 193 membres de l'ONU) et abrite à Yaoundé une \textbf{quarantaine d'ambassades étrangères} résidentes. Surnommé « l'Afrique en miniature » par sa diversité, il possède une originalité historique : il est membre à la fois du \textbf{Commonwealth} (depuis 1995) et de la \textbf{Francophonie} (institutionnelle depuis 1970). Cette double appartenance, héritée de la colonisation allemande (1884-1916), puis du mandat français et britannique, lui offre une capacité de dialogue unique sur la scène internationale.
\end{savaistu}

\courssub{Les domaines des relations internationales}

Les relations internationales se déploient dans plusieurs domaines, souvent imbriqués les uns dans les autres.

\begin{center}
\begin{tabularx}{\linewidth}{|l|X|X|}\hline
\rowcolor{popblueL} \textbf{Domaine} & \textbf{Contenu principal} & \textbf{Exemple camerounais récent} \\ \hline
Diplomatique / Politique & Dialogue entre États, sommets, visites officielles, médiation, vote à l'ONU & Médiation du Cameroun dans la crise centrafricaine ; vote résolutions ONU \\ \hline
Économique & Commerce international, investissements directs (IDE), aide au développement, dette & Accords routes/barrages avec la Chine ; Partenariat UE-Cameroun (APE) \\ \hline
Culturel & Échanges artistiques, linguistiques, universitaires, scientifiques, patrimoine & Bourses d'études (France, Chine, Maroc, Russie) ; Institut Goethe, Institut français \\ \hline
Sportif & Compétitions internationales, rayonnement du pays, diplomatie par le sport & Lions Indomptables (CAN, Coupe du monde) ; organisation CHAN 2020 / CAN 2021 \\ \hline
Sécuritaire / Militaire & Alliances, formation, maintien de la paix (OMP), lutte antiterroriste, renseignement & Casques bleus MINUSCA (RCA) ; Force multinationale mixte (FMM) Lac Tchad \\ \hline
\end{tabularx}
\end{center}

\begin{exemplebox}[Le football, un outil de diplomatie publique]
Un match international des Lions Indomptables n'est pas qu'une fête sportive : c'est un vecteur de relations internationales. Il projette l'image du Cameroun dans le monde entier (soft power), crée des échanges entre fédérations (FECAFOOT / FIFA / CAF), supporteurs et médias, et peut même apaiser des tensions politiques entre pays. On parle de \emph{diplomatie sportive}. De même, l'organisation de la CAN 2021 (jouée en 2022) a nécessité une coopération intense avec la CAF, les 23 pays participants, des entreprises de construction (Chine, Turquie, France) et des médias internationaux.
\end{exemplebox}

\courssub{Les instruments des relations internationales}

Pour dialoguer avec le monde, les États disposent d'instruments juridiques et institutionnels précis.

\begin{definition}[Traité, Convention, Accord international]
Un \cle{traité} (ou accord international) est un engagement écrit, conclu entre sujets de droit international (États, organisations internationales), régi par le droit international, qu'il soit contenu dans un instrument unique ou dans plusieurs instruments connexes, quelle que soit sa dénomination (Convention, Accord, Protocole, Charte, Pacte, Statut). Une \cle{convention} désigne souvent un traité multilatéral à portée normative (ex. Convention de Vienne sur le droit des traités, Convention sur les droits de l'enfant).
\end{definition}

Les autres instruments institutionnels sont :
\begin{itemize}
\item les \cle{ambassades} (ou missions diplomatiques) : représentations permanentes d'un État auprès d'un autre État, installées dans la capitale, dirigées par un \cle{ambassadeur} (chef de mission). L'ambassade du Cameroun à Paris défend les intérêts camerounais en France ; l'ambassade de France à Yaoundé représente la France au Cameroun ;
\item les \cle{consulats} : antennes de l'ambassade situées dans d'autres villes que la capitale (ex. Consulat du Cameroun à Lyon, à Marseille, à Douala pour les pays étrangers), chargées surtout de la \cle{protection consulaire} : services aux citoyens (passeports, actes d'état civil, visas, assistance aux détenus, rapatriements) ;
\item les \textbf{missions permanentes} : représentations auprès des organisations internationales (Mission permanente du Cameroun auprès de l'ONU à New York et Genève, auprès de l'UA à Addis-Abéba) ;
\item les \textbf{sommets et conférences internationales} : réunions au plus haut niveau (Chefs d'État, Ministres) pour décider de l'orientation politique (Sommet CEMAC, Assemblée générale ONU, COP Climat, Sommet Afrique-France) ;
\item les \textbf{missions diplomatiques spéciales} : envoyés spéciaux, médiateurs, observateurs électoraux, missions de bons offices.
\end{itemize}

\begin{popfigure}
\begin{center}
\begin{tikzpicture}[scale=0.9]
% Cadre monde
\draw[fill=popcream,draw=popline,rounded corners=4pt] (-7.5,-3) rectangle (7.5,3);
% Zones textuelles
\node[align=center,font=\footnotesize] at (-5.5,2.2) {\textbf{AMÉRIQUES}\\ Washington, Ottawa\\ Brasília, Mexico};
\node[align=center,font=\footnotesize] at (-0.5,2.3) {\textbf{EUROPE}\\ Paris, Londres, Berlin\\ Bruxelles, Rome, Moscou};
\node[align=center,font=\footnotesize] at (5,2.3) {\textbf{ASIE / MOYEN-ORIENT}\\ Pékin, Tokyo, New Delhi\\ Riyad, Ankara, Doha};
\node[align=center,font=\footnotesize] at (0.5,0.0) {\textbf{AFRIQUE}\\ Abuja, Pretoria, Alger\\ Le Caire, Addis-Abéba (UA)\\ Libreville, N'Djamena, Brazzaville (CEMAC)};
% Points ambassades (orange)
\foreach \p in {(-5.5,1.5),(-5.0,1.5),(-6.0,1.5), (-0.5,1.6),(-1.0,1.6),(0.0,1.6), (5,1.6),(5.5,1.6),(4.5,1.6), (0.5,-0.7),(1.0,-0.7),(0.0,-0.7)}
  \fill[poporange] \p circle (0.07);
% Yaoundé
\fill[popdark] (0.5,-2.2) circle (0.12);
\node[font=\footnotesize\bfseries,popdark,align=center] at (0.5,-2.7) {Yaoundé\\ Siège de la diplomatie\\ camerounaise};
% Flèches
\draw[->,thick,popdark] (0.5,-2.1) -- (-4.5,1.6);
\draw[->,thick,popdark] (0.5,-2.1) -- (-0.4,1.7);
\draw[->,thick,popdark] (0.5,-2.1) -- (4.5,1.7);
\draw[->,thick,popdark] (0.5,-2.1) -- (0.5,-0.8);
% Flèche Nord
\draw[->,thick,popdark] (6.5,2.5) -- (6.5,2.8) node[above,font=\tiny\bfseries] {N};
\end{tikzpicture}
\end{center}
\legende{Carte schématique du monde (projection très simplifiée) localisant les principaux pôles de partenaires diplomatiques du Cameroun. Points orange : zones de concentration d'ambassades camerounaises à l'étranger. Flèches : flux diplomatiques bilatéraux. Yaoundé (point noir) accueille en retour une quarantaine d'ambassades étrangères et les missions permanentes auprès des organisations sous-régionales.}
\end{popfigure}

\courssub{Étude de documents et situations concrètes}

\begin{methode}[Analyser un document diplomatique (photo, communiqué, carte, discours)]
Pour analyser un document relatif aux relations internationales, suivre quatre étapes rigoureuses :
\begin{enumerate}
\item \textbf{Identifier les acteurs} : quels États (Chefs d'État, Ministres), quelles organisations (ONU, UA, CEMAC), quelles personnalités, quelles entreprises ou ONG sont présents ou signataires ?
\item \textbf{Situer le lieu et la date} : où et quand l'événement a-t-il lieu ? Le choix du lieu (capital, siège d'organisation, zone de conflit) est souvent un message politique en soi.
\item \textbf{Déterminer l'objet et le domaine} : de quoi traite concrètement l'acte (signature d'accord, déclaration commune, résolution, aide d'urgence) ? À quel domaine rattache-t-on l'action (politique, économique, sécuritaire, culturel, humanitaire) ?
\item \textbf{Dégager les enjeux et intérêts} : que cherche chaque acteur ? (Sécurité, accès aux ressources/marchés, image/soft power, influence régionale, appui budgétaire, respect du droit international).
\end{enumerate}
Conclure en rattachant l'analyse à la définition des relations internationales (liens politiques, économiques, culturels, militaires entre acteurs).
\end{methode}

\begin{exoresolu}[Analyse d'une situation : Visite d'État d'un partenaire stratégique au Cameroun]
Énoncé. Le Président de la République d'un pays partenaire effectue une visite d'État de trois jours à Yaoundé. Au programme : entretien en tête-à-tête avec le Président de la République du Cameroun, signature de cinq accords de coopération (santé, routes, formation professionnelle, agriculture, défense), visite d'une usine de transformation de bois financée par ce pays, dépôt de gerbe au monument de la Réunification. À l'aide de la méthode d'analyse, montre que cette visite illustre les relations internationales.
\tcblower
Solution détaillée.
\begin{enumerate}
\item \textbf{Acteurs} : Deux États souverains représentés par leurs Chefs d'État, leurs Ministres des Affaires étrangères et sectoriels (Santé, Travaux publics, Éducation, Défense), ainsi que des opérateurs économiques des deux pays (usine de bois).
\item \textbf{Lieu et date} : Yaoundé, capitale politique du Cameroun, siège des institutions. Le déplacement du Chef d'État étranger honore le pays hôte et marque l'importance du partenariat.
\item \textbf{Objet et domaines} : Renforcement de la coopération bilatérale. Domaines mobilisés : \emph{Diplomatique} (entretiens, honneurs), \emph{Économique} (accords routes, usine, investissements), \emph{Social / Culturel} (santé, formation professionnelle), \emph{Sécuritaire} (accord défense).
\item \textbf{Enjeux} : Pour le Cameroun : diversification des financements, transfert de technologies, formation de la jeunesse, équipement militaire, appui à l'industrialisation (transformation locale du bois). Pour le pays partenaire : sécurisation d'approvisionnements (bois, matières premières), accès au marché camerounais et sous-régional (CEMAC), renforcement de son influence géopolitique en Afrique centrale.
\end{enumerate}
\textbf{Conclusion} : Cette visite d'État illustre parfaitement la définition des relations internationales, car elle crée et renforce des liens \cle{politiques} (dialogue au sommet), \cle{économiques} (investissements, commerce), \cle{culturels/sociaux} (formation, santé) et \cle{militaires} (défense) entre deux États, par l'instrument juridique des \cle{traités/accords} signés et le relais institutionnel des \cle{ambassades} qui en assureront le suivi.
\end{exoresolu}

\begin{exercice}[Entraînement : Appliquer les notions]
1. Définis les relations internationales en citant leurs quatre grands domaines constitutifs.\\
2. Classe les acteurs suivants dans les cinq catégories vues en cours : \textbf{ONU}, \textbf{Nigeria}, \textbf{Croix-Rouge internationale}, \textbf{MTN}, \textbf{Diaspora camerounaise d'Allemagne}.\\
3. Un élève affirme : « Le match Cameroun--Brésil lors de la Coupe du monde n'a rien à voir avec les relations internationales, c'est juste du sport. » Réfute cette affirmation en développant \textbf{deux arguments} distincts fondés sur le cours.\\
4. Différencie clairement le rôle d'une \textbf{ambassade} et celui d'un \textbf{consulat}.
\end{exercice}

\begin{corrige}[Exercice 1]
1. Les relations internationales sont l'ensemble des liens \cle{politiques}, \cle{économiques}, \cle{culturels} et \cle{militaires} qui unissent les \cle{États} et les autres \cle{acteurs de la scène mondiale} (organisations internationales, ONG, entreprises multinationales, diasporas).\\
2. \textbf{ONU} : Organisation internationale ; \textbf{Nigeria} : État ; \textbf{Croix-Rouge internationale} : ONG ; \textbf{MTN} : Entreprise multinationale ; \textbf{Diaspora camerounaise d'Allemagne} : Diaspora.\\
3. \textbf{Argument 1 (Rayonnement / Soft power) :} Le match oppose deux nations souveraines et projette l'image, la culture et l'identité du Cameroun auprès d'une audience mondiale de plusieurs milliards de téléspectateurs : c'est de la diplomatie publique.\\
   \textbf{Argument 2 (Cadre institutionnel / Échanges) :} La rencontre suppose des accords entre fédérations (FECAFOOT / CBF / FIFA / CAF), l'obtention de visas, le déplacement d'une délégation officielle (ministre, ambassadeur), la présence de supporters et de médias des deux pays : ce sont des flux humains, juridiques et médiatiques transfrontaliers.\\
4. L'\textbf{ambassade}, située dans la \cle{capitale} du pays d'accueil et dirigée par un \cle{ambassadeur}, est la représentation politique permanente : elle conduit la diplomatie, négocie les traités, suit la politique locale et rend compte à son gouvernement. Le \textbf{consulat}, situé dans d'autres \cle{grandes villes} (pas nécessairement la capitale), est une administration de proximité : il assure la \cle{protection consulaire} (délivrance passeports, visas, actes d'état civil, aide aux ressortissants en difficulté, organisation des votes à l'étranger).
\end{corrige}

\begin{aretenir}[L'essentiel de la section : Le concept de relations internationales]
\begin{itemize}
\item \textbf{Définition} : Les \cle{relations internationales} = ensemble des liens \cle{politiques}, \cle{économiques}, \cle{culturels}, \cle{militaires} entre \cle{États} et \cle{acteurs non étatiques} (OI, ONG, Multinationales, Diasporas).
\item \textbf{Acteurs} : 5 catégories : États (souveraineté), Organisations internationales (ONU, UA, CEMAC), ONG, Multinationales, Diasporas.
\item \textbf{Domaines} : Diplomatique, Économique, Culturel, Sportif, Sécuritaire (souvent imbriqués).
\item \textbf{Instruments} : Juridiques (\cle{traité}, \cle{convention}) et Institutionnels (\cle{ambassade}, \cle{consulat}, mission permanente, sommets, missions spéciales).
\item \textbf{Cameroun} : Acteur majeur (relations avec 180+ pays, 40+ ambassades à Yaoundé, double appartenance Commonwealth/Francophonie, contributeur OMP, hub CEMAC).
\end{itemize}
\end{aretenir}

\coursec{Les principes de la diplomatie camerounaise}

\courssub{Le concept de relations internationales}

\begin{definition}[Relations Internationales (RI)]
Les \textbf{Relations Internationales} désignent l'ensemble des interactions, formelles ou informelles, conflictuelles ou coopératives, qui s'établissent entre les \textbf{acteurs} de la scène mondiale (États, organisations internationales, organisations non gouvernementales, entreprises multinationales, collectivités territoriales, individus) pour défendre leurs intérêts, résoudre des différends ou construire des biens communs.
\end{definition}

\begin{propriete}[Les acteurs majeurs des RI]
\begin{itemize}
    \item \textbf{Les États :} Acteurs souverains principaux, détenteurs du monopole de la violence légitime et de la personnalité juridique internationale (siège à l'ONU).
    \item \textbf{Les Organisations Internationales (OI) :} ONU, UA, CEMAC, CEEAC, OIF, OMC, FMI, Banque Mondiale. Elles offrent des cadres normatifs et des forums de négociation.
    \item \textbf{Les ONG internationales :} MSF, Amnesty International, Croix-Rouge, WWF. Acteurs de plaidoyer, d'urgence humanitaire et de surveillance des droits humains.
    \item \textbf{Les Entreprises Multinationales (EMN) :} Acteurs économiques puissants (ex. : Bolloré, CHEC, TotalEnergies, Dangote au Cameroun) influençant les politiques publiques par l'investissement.
    \item \textbf{Les collectivités territoriales :} Coopération décentralisée (villes jumelées, régions partenaires).
\end{itemize}
\end{propriete}

\begin{exemplebox}[Distinction Relations Internationales vs Coopération Internationale]
\begin{center}
\begin{tabularx}{\linewidth}{|l|X|X|}\hline
\textbf{Critère} & \textbf{Relations Internationales (Cadre global)} & \textbf{Coopération Internationale (Volet constructif)} \\ \hline
\textbf{Nature} & Inclut conflits, rivalités, sanctions, guerre froide, diplomatie coercitive. & Action volontaire, solidaire, mutuellement avantageuse (gain-gagnant). \\ \hline
\textbf{Objectif} & Défense de l'intérêt national, puissance, sécurité, influence. & Développement partagé, résolution problèmes communs, paix durable. \\ \hline
\textbf{Outils} & Traité, alliance, sanction, embargo, médiation, force armée. & Accord-cadre, don, prêt concessionnel, expertise, formation, transfert tech. \\ \hline
\end{tabularx}
\end{center}
\legende{La coopération est un \emph{sous-ensemble} des relations internationales : toute coopération est une relation internationale, mais toute relation internationale n'est pas de la coopération.}
\end{exemplebox}

\courssub{Les aspects de la coopération internationale}

Après avoir défini les relations internationales comme l'ensemble des interactions entre acteurs étatiques et non étatiques sur la scène mondiale, nous comprenons que la \cle{coopération internationale} en constitue le volet le plus constructif et le plus opérationnel. Là où les relations internationales englobent aussi les conflits et les rivalités, la coopération désigne spécifiquement la volonté d'agir \emph{ensemble} pour transformer une situation initiale en une situation meilleure pour toutes les parties. C'est le passage de la coexistence à la co-construction.

\courssub{Définition et fondements de la coopération internationale}

\begin{definition}[La coopération internationale]
La \textbf{coopération internationale} est l'action solidaire menée en commun par plusieurs États ou organisations (internationales, régionales, non gouvernementales) pour résoudre des problèmes partagés, atteindre des objectifs communs et promouvoir le développement mutuel, dans le respect de la souveraineté de chacun et du droit international.
\end{definition}

Cette définition repose sur trois piliers indissociables :
\begin{itemize}
    \item \textbf{La solidarité :} Reconnaissance d'une interdépendance (aucun pays ne peut seul faire face au changement climatique, aux pandémies ou au terrorisme transfrontalier).
    \item \textbf{La réciprocité :} Chaque partenaire apporte une contribution (financière, technique, humaine, diplomatique) et en retire un avantage, même asymétrique.
    \item \textbf{La contractualisation :} La coopération se formalise par des accords, traités, conventions ou protocoles d'entente qui fixent les objectifs, les moyens, le calendrier et les instances de suivi-évaluation.
\end{itemize}

\begin{methode}[Comment identifier et classifier un acte de coopération]
Pour analyser une situation concrète et en déterminer l'aspect dominant, suivez cette démarche en quatre étapes :
\begin{enumerate}
    \item \textbf{Identifier les acteurs :} Quels États, organisations internationales (ONU, UA, CEMAC), ONG, collectivités territoriales ou entreprises publiques/privées sont parties prenantes ?
    \item \textbf{Repérer le besoin ou le problème visé :} S'agit-il d'un déficit infrastructurel, d'une menace sécuritaire, d'une urgence sanitaire, d'un besoin de formation, de la préservation d'un écosystème partagé ?
    \item \textbf{Qualifier la nature des flux :} Argent (dons, prêts, investissements), biens (équipements, médicaments, vivres), savoir-faire (experts, formation, transfert de technologie), personnel (militaires, médecins, enseignants), cadres juridiques (traité, accord).
    \item \textbf{Classer l'aspect principal (sans exclure les autres) :} Politique/diplomatique, Économique/commercial, Culturel/scientifique, Militaire/sécuritaire, Sanitaire/environnemental.
\end{enumerate}
\end{methode}

\begin{exemplebox}[Application de la méthode : le projet de modernisation de l'École de Guerre de Yaoundé]
\begin{itemize}
    \item \textbf{Acteurs :} Cameroun (État hôte), France (partenaire technique/financier), partenaires africains (stagiaires de la sous-région).
    \item \textbf{Besoin :} Renforcement des capacités de commandement et d'état-major pour la sécurité régionale.
    \item \textbf{Flux :} Infrastructures (bâtiments), expertise (officiers instructeurs français), financement partagé, programmes pédagogiques.
    \item \textbf{Classification :} \textbf{Principalement militaire et sécuritaire} (formation, interopérabilité), avec une dimension \textbf{diplomatique} (renforcement des liens d'amitié, influence stratégique) et \textbf{économique} (marchés de construction, logistique).
\end{itemize}
\end{exemplebox}

\courssub{L'aspect politique et diplomatique : le dialogue au service de la paix}

C'est la matrice historique de la coopération. Elle vise à prévenir les conflits, à les résoudre par la médiation et à construire des cadres juridiques stables.

\begin{itemize}
    \item \textbf{Les alliances et organisations régionales :} Le Cameroun est membre fondateur de l'UA, de la CEMAC, de la CEEAC, de la CIRGL. Ces cadres offrent des mécanismes de concertation permanente (sommet des chefs d'État, conseils des ministres).
    \item \textbf{La médiation et la « diplomatie préventive » :} Le Cameroun a joué un rôle clé dans la résolution de la crise centrafricaine (Accords de Libreville 2008, puis facilitation du dialogue politique 2019-2020 sous l'égide de l'UA et de la CEEAC). Le Président Paul Biya a souvent été mandaté comme facilitateur ou médiateur (ex. : crise post-électorale en Guinée-Conakry, tensions Tchad-Soudan).
    \item \textbf{La coopération juridique :} Extradition, entraide judiciaire, harmonisation des législations (droit des affaires OHADA, droit pénal international via la CPI).
\end{itemize}

\begin{exoresolu}[Analyse d'un sommet : le rôle du Cameroun dans le processus de paix en Centrafrique]
\textbf{Document :} Extrait du communiqué final du Sommet de la CEEAC sur la RCA (N'Djaména, janvier 2019) — « Les Chefs d'État mandatent le Président Paul Biya pour poursuivre la facilitation du dialogue politique inclusif... »
\textbf{Questions :}
\begin{enumerate}
    \item Quel aspect de la coopération illustre ce mandat ?
    \item Pourquoi le choix du Cameroun est-il pertinent ?
\end{enumerate}
\tcblower
\textbf{Solution :}
\begin{enumerate}
    \item C'est l'\textbf{aspect politique et diplomatique} (médiation, facilitation, diplomatie préventive). Il s'agit d'une coopération \emph{multilatérale} (CEEAC) mise en œuvre par un acteur \emph{étatique} (Cameroun) au profit d'un tiers (RCA).
    \item \textbf{Pertinence :} (1) Voisin direct (frontière de 900 km), impact direct des flux de réfugiés et de l'insécurité transfrontalière. (2) Histoire partagée (ancien territoire de l'AEF, liens humains forts). (3) Crédibilité diplomatique (membre permanent du Conseil de paix et de sécurité de l'UA, expérience de médiation). (4) Neutralité perçue (pas de troupes au sol en RCA dans le cadre de la MINUSCA, contrairement au Tchad ou au Rwanda).
\end{enumerate}
\end{exoresolu}

\courssub{L'aspect économique et commercial : le moteur du développement}

C'est aujourd'hui l'aspect le plus visible et le plus quantifié. Il recouvre l'aide publique au développement (APD), les investissements directs étrangers (IDE), le commerce, le transfert de technologies et la coopération monétaire.

\begin{itemize}
    \item \textbf{L'Aide Publique au Développement (APD) :} Dons et prêts à conditions concessionnelles (taux d'intérêt bas, longues maturités). Principaux bailleurs : France (via AFD), Banque mondiale (IDA), Banque africaine de développement (BAD), Union européenne, Allemagne (KfW), Japon (JICA), Chine (prêts concessionnels et dons).
    \item \textbf{Les Investissements Directs Étrangers (IDE) :} Création d'usines, acquisition de participations, concessions (ports, barrages, mines, télécoms). Le Cameroun attire dans l'énergie (barrages de Nachtigal, Memve'ele), les infrastructures (routes, ports), l'agro-industrie, les télécoms.
    \item \textbf{Le commerce international :} Exportations (bois, cacao, coton, pétrole, aluminium, bananes) et importations (produits pétroliers raffinés, machines, véhicules, céréales, médicaments). La Zone de Libre-Échange Continentale Africaine (ZLECAf) est le nouveau cadre de coopération commerciale majeure.
\end{itemize}

\begin{exemplebox}[Focus : La coopération sino-camerounaise et le Port en eau profonde de Kribi]
Le \textbf{Port autonome de Kribi (PAK)} et sa zone industrielle adjacente constituent l'emblème de la coopération économique Cameroun-Chine.
\begin{itemize}
    \item \textbf{Acteurs :} État du Cameroun (maître d'ouvrage), China Harbour Engineering Company (CHEC) — constructeur/exploitant (phase 1), Exim Bank of China — financeur (prêt concessionnel \approx 80\% du coût), entreprises camerounaises (sous-traitance).
    \item \textbf{Coût :} Phase 1 : \num{600} milliards FCFA ; Phase 2 (en cours) : \num{800} milliards FCFA.
    \item \textbf{Performances :} Capacité d'accueil de navires de \num{300} mètres (post-Panamax), tirant d'eau \num{18} m. Plateforme logistique pour la sous-région (Tchad, RCA, Guinée équatoriale).
    \item \textbf{Externalités :} Création de \num{5000} emplois directs/indirects, formation de cadres portuaires, développement de la ville de Kribi, attraction d'industries (usine d'alumine, aciérie, engrais) dans la zone économique spéciale.
\end{itemize}
\end{exemplebox}

\begin{popfigure}
\begin{tikzpicture}
\begin{scope}[scale=0.9]
% Fond suggérant une photo satellite / vue aérienne
\fill[popblue!10] (-5,-3) rectangle (5,3);
% Représentation schématique du port
\draw[popdark, line width=2pt, fill=popblue!20] (-3,-1.5) -- (-3,1.5) -- (-1,1.5) -- (-1,-1.5) -- cycle; % Quai principal
\draw[popdark, line width=1.5pt, fill=popmuted!30] (-1,-1.5) -- (3,-1) -- (3,1) -- (-1,1.5) -- cycle; % Bassin
\draw[poporange, line width=1pt, fill=poporange!30] (-3.5,-1) rectangle (-2.5,1); % Grue 1
\draw[poporange, line width=1pt, fill=poporange!30] (-3.5,-0.2) rectangle (-2.5,0.2); % Grue 2
\draw[poporange, line width=1pt, fill=poporange!30] (-3.5,0.6) rectangle (-2.5,1); % Grue 3
% Navire
\draw[popdark, line width=1.5pt, fill=popmuted!50] (0.5,-0.6) .. controls (1.5,-0.6) and (2.5,0.6) .. (2.5,0.6) -- (2.5,0.8) .. controls (1.5,0.8) and (0.5,0.8) .. (0.5,0.8) -- cycle;
\draw[popdark, fill=popdark] (1.5,0.8) rectangle (1.8,1.3); % Superstructure
% Légendes intégrées
\node[popdark, font=\small\bfseries, align=center] at (-2,2.2) {Quai principal (600 m)};
\node[popdark, font=\small, align=center] at (-3, -2.2) {Grues de quai (STS)};
\node[popdark, font=\small, align=center] at (2.5, -1.5) {Navire Post-Panamax};
\node[popdark, font=\small, align=center] at (3.5, 0) {Bassin \SI{18}{m}};
\node[popdark, font=\small, align=center] at (-4.5, 2.5) {Zone industrielle / ZES};
\draw[->, popdark, thick] (-4, 2.3) -- (-2.5, 1.8);
\end{scope}
\end{tikzpicture}
\legende{Schéma simplifié du Port Autonome de Kribi (Phase 1) : quai principal, bassin en eau profonde (\SI{18}{m}), grues de quai (STS) et navire de grande capacité. Infrastructure réalisée dans le cadre de la coopération économique Cameroun-Chine (financement Exim Bank, réalisation CHEC).}
\end{popfigure}

\begin{aretenir}[Chiffres clés : Coopération économique Cameroun — Principaux partenaires (Données récentes, estimations)]
\begin{center}
\begin{tabularx}{\linewidth}{|l|X|X|X|}\hline
\textbf{Partenaire} & \textbf{Nature dominante} & \textbf{Projets emblématiques} & \textbf{Volume échanges / Engagements (ordre de grandeur)} \\ \hline
\textbf{Chine} & IDE, Prêts concessionnels, Commerce & Port de Kribi, Barrages (Memve'ele, Nachtigal), Routes, Stade Olembé & \textbf{1er partenaire commercial} (\textgreater 1000 Mds FCFA/an), \textasciitilde 3000 Mds FCFA engagements infrastructure \\ \hline
\textbf{France / AFD} & APD (Dons/Prêts), IDE, Expertise & Eau/Assainissement (Yaoundé, Douala), Énergie (réseau), Santé, Éducation, Appui budgétaire & \textasciitilde 500-800 Mds FCFA/an (APD + IDE), 1er investisseur historique (stock) \\ \hline
\textbf{Union Européenne} & APD (FED), Appui budgétaire, Commerce & Gouvernance, Développement rural, Transport, ZLECAf & \textasciitilde 300-400 Mds FCFA / période pluriannuelle \\ \hline
\textbf{Banque Mondiale (IDA) / BAD} & Prêts concessionnels, Dons, Expertise technique & Transport (routes), Énergie, Filets sociaux, Climat, Compétences & \textasciitilde 1000-1500 Mds FCFA portefeuille actif \\ \hline
\textbf{Autres (Japon, Allemagne, Corée, Turquie, Pays du Golfe, Inde...)} & Projets sectoriels ciblés & Riziculture (Japon), Formation pro (Allemagne), Santé (Corée), Infrastructures (Turquie) & Variables, en croissance \\ \hline
\end{tabularx}
\end{center}
\legende{Sources : MINFI, Banque Mondiale, OCDE/DAC, Douanes Cameroun. Les montants sont des ordres de grandeur annuels ou de portefeuille, sujets à variation.}
\end{aretenir}

\begin{attention}[Piège : Confondre APD et IDE]
\begin{itemize}
    \item \textbf{APD (Aide Publique au Développement) :} Ressource \emph{non commerciale}, octroyée par des États/agences publiques à des pays en développement. Contient un \textbf{don} (\ge 25\% d'élément de don pour l'OCDE). Objectif : bien-être, réduction pauvreté. Ex. : construction d'un hôpital financé par don AFD.
    \item \textbf{IDE (Investissement Direct Étranger) :} Flux \emph{commercial/privé}. Un investisseur (entreprise, fonds souverain) acquiert un contrôle durable (\ge 10\% capital) dans une entreprise locale. Objectif : profit, parts de marché. Ex. : CHEC exploite le terminal à conteneurs de Kribi via une concession BOT (Build-Operate-Transfer).
    \item \textbf{Risque :} Un même projet (ex. barrage de Nachtigal) peut mêler les deux : prêt concessionnel BAD/AFD (part APD) + capital privé EDF/STOA (part IDE). Ne les additionnez pas bêtement.
\end{itemize}
\end{attention}

\courssub{L'aspect culturel et scientifique : le rayonnement et le capital humain}

Souvent sous-estimé, cet aspect est le ciment de la compréhension mutuelle et le vivier des compétences futures.

\begin{itemize}
    \item \textbf{Les bourses d'études et la formation :} Programme \textbf{Eiffel} (France), \textbf{CSC} (Chine), \textbf{DAAD} (Allemagne), \textbf{Fulbright} (USA), \textbf{MEXT} (Japon), \textbf{Stipendium Hungaricum} (Hongrie), bourses russes, marocaines, turques, indiennes. Des milliers de cadres camerounais formés à l'étranger (médecins, ingénieurs, administrateurs, chercheurs).
    \item \textbf{Les instituts culturels et centres de langue :} \textbf{Institut français du Cameroun} (IFC - 5 antennes), \textbf{Institut Confucius} (Yaoundé, Douala - langue chinoise, culture), \textbf{Goethe-Institut} (Yaoundé - allemand), \textbf{Centro Cultural Español}, \textbf{British Council}. Ils offrent cours de langue, certifications (DELF/DALF, HSK, TestDaF), médiathèques, programmation artistique.
    \item \textbf{La coopération universitaire et scientifique :} Cotutelles de thèse, laboratoires mixtes internationaux (LMI), projets de recherche conjoints (IRD, CIRAD, CNRS avec universités camerounaises : Yaoundé I, II, Douala, Ngaoundéré, Buea, Dschang). Ex. : recherche sur le paludisme, le VIH, les plantes médicinales, l'agroforesterie, les énergies renouvelables.
    \item \textbf{Le patrimoine et la francophonie :} Classement UNESCO (Réserve de la Dja, Paysage des chutes de la Lobé), participation active à l'OIF (Sommet de Yaoundé 1999, Jeux de la Francophonie).
\end{itemize}

\begin{savaistu}[Le « soft power » camerounais par la culture]
Le Cameroun exporte sa culture comme vecteur d'influence : la \textbf{makossa}, le \textbf{bikutsi}, la littérature (Mongo Beti, Ferdinand Oyono, Calixthe Beyala, Patrice Nganang), le cinéma (Jean-Pierre Dikongué Pipa, Rosine Mbakam), la mode (Imane Ayissi). La participation aux Jeux de la Francophonie, à la CAN, aux JO crée une image positive qui facilite les autres coopérations. L'organisation de la \textbf{CHAN 2020} et de la \textbf{CAN 2021} (reportée à 2022) a été un exercice de coopération logistique, sécuritaire et diplomatique majeur.
\end{savaistu}

\courssub{L'aspect militaire et sécuritaire : la défense collective face aux menaces transnationales}

La sécurité est un « bien public » qui ne s'arrête pas aux frontières. La coopération militaire est encadrée par le droit international (Charte ONU Art. 51, Traité de non-agression CEMAC, Protocole de la CEEAC).

\begin{itemize}
    \item \textbf{La Force Multinationale Mixte (FMM) :} Créée en 1994 (Nigéria, Niger, Tchad, Cameroun), revitalisée en 2015 sous l'égide de l'UA et de la CEEAC pour lutter contre \textbf{Boko Haram / Etat islamique en Afrique de l'Ouest (ISWAP)} dans le bassin du Lac Tchad. Quartier général : N'Djaména. Le Cameroun fournit le secteur 1 (Extrême-Nord) et le commandement adjoint.
    \item \textbf{Opérations conjointes bilatérales :} Patrouilles mixtes Cameroun-Nigéria (frontière), Cameroun-Tchad, Cameroun-RCA.
    \item \textbf{Formation et renforcement des capacités :} Écoles militaires (EMIA, École de Guerre, École de maintien de la paix de Douala - \textbf{EMPABB}) formées avec appui France, USA, Allemagne, ONU. Formation aux droits de l'homme, DIH, protection des civils.
    \item \textbf{Contributions aux Opérations de Maintien de la Paix (OMP) :} Le Cameroun est un contributeur majeur de troupes (Casques bleus) : MINUSCA (RCA), MINUSMA (Mali - fin 2023), MONUSCO (RDC), MINURSO (Sahara occidental). Cela génère des revenus, de l'expérience et du prestige diplomatique.
\end{itemize}

\begin{experience}[Simulation : Conseil de sécurité régional (CEEAC) — Crise frontalière]
\textbf{Scénario :} Incursion de groupes armés non identifiés depuis le pays voisin X dans la région de l'Est du Cameroun. 5 villages attaqués, 2000 déplacés.
\textbf{Rôles :} Ministre Défense Cameroun, Ministre Affaires Étrangères, Représentant ONU, Représentant UA, Chef d'État-Major, ONG humanitaire.
\textbf{Tâche :} En 20 min, proposer une réponse coordonnée mobilisant \textbf{au moins 3 aspects de la coopération} (ex. : Militaire \rightarrow FMM/patrouilles mixtes ; Humanitaire \rightarrow HCR/OCHA ; Diplomatique \rightarrow Note verbale, médiation CEEAC ; Sanitaire \rightarrow OMS/MSF).
\textbf{Débriefing :} Montre l'interdépendance des aspects. Une réponse purement militaire sans volet humanitaire et diplomatique échoue durablement.
\end{experience}

\courssub{L'aspect sanitaire et environnemental : la santé planétaire et les biens communs mondiaux}

La pandémie de COVID-19 a brutalement rappelé qu'un virus n'a pas de passeport. La coopération sanitaire est devenue une priorité absolue.

\begin{itemize}
    \item \textbf{Lutte contre les pandémies :} \textbf{COVID-19} : Mécanisme COVAX (Gavi, OMS, CEPI) \rightarrow livraison de vaccins (AstraZeneca, Pfizer, Sinopharm, Johnson\&Johnson). Appui technique (OMS, CDC Africa, ONG). \textbf{Ébola, Choléra, Paludisme, VIH/Sida} : Fonds mondial, PEPFAR (USA), Initiative Clinton, GAVI. Renforcement du système de santé (labos, chaîne de froid, santé communautaire).
    \item \textbf{Protection de l'environnement et du climat :} Le Cameroun, « Afrique en miniature », abrite une partie du \textbf{Bassin du Congo} (2e poumon vert mondial après l'Amazonie).
    \begin{itemize}
        \item \textbf{COMIFAC / CBFP :} Commission des Forêts d'Afrique Centrale / Partenariat pour les Forêts du Bassin du Congo. Gestion durable, lutte anti-braconnage, certification FSC.
        \item \textbf{Initiative pour la Résilience du Bassin du Lac Tchad :} Reboisement, gestion de l'eau, énergie solaire, lutte contre la désertification (Grande Muraille Verte).
        \item \textbf{Accords climat :} Ratification Accord de Paris. Contribution Déterminée au Niveau National (CDN) : réduction 35\% émissions à l'horizon 2030 (conditionnel à l'appui international). Mécanismes REDD+ (Réduction des émissions liées à la déforestation).
    \end{itemize}
    \item \textbf{Coopération nucléaire civile :} AIEA — formation d'experts, radiothérapie (cancer), agriculture (variétés mutées), gestion de l'eau (isotopes).
\end{itemize}

\begin{exoresolu}[Classer des actions de coopération sanitaire pendant la COVID-19]
\textbf{Documents :}
\begin{enumerate}
    \item Arrivée de 391 200 doses vaccin AstraZeneca via COVAX (avril 2021).
    \item Don de 200 000 doses Sinopharm par la Chine (avril 2021).
    \item Appui technique OMS : déploiement 50 experts épidémiologistes dans les 10 régions.
    \item Prêt BAD \num{88} milliards FCFA « Appui à la riposte COVID-19 » (budget, filets sociaux, intrants).
    \item Formation biomédicale (IRD/INSERM) : séquençage variants au Centre Pasteur Yaoundé.
\end{enumerate}
\textbf{Consigne :} Pour chaque action, précisez l'aspect dominant (Sanitaire / Économique / Diplomatique / Scientifique) et le type de coopération (Multilatérale / Bilatérale).
\tcblower
\textbf{Solution :}
\begin{center}
\begin{tabularx}{\linewidth}{|X|l|l|}\hline
\textbf{Action} & \textbf{Aspect dominant} & \textbf{Type} \\ \hline
1. Vaccins COVAX & Sanitaire (santé publique mondiale) & Multilatérale (ONU/Gavi) \\ \hline
2. Don Sinopharm Chine & Sanitaire + \textbf{Diplomatique} (soft power) & Bilatérale (Cameroun-Chine) \\ \hline
3. Experts OMS & Sanitaire + \textbf{Scientifique/Technique} & Multilatérale (ONU) \\ \hline
4. Prêt BAD & \textbf{Économique/Budgétaire} (soutien macroéco) + Sanitaire & Multilatérale (Banque régionale) \\ \hline
5. Séquençage variants & \textbf{Scientifique} (recherche, surveillance génomique) & Bilatérale/Multilatérale (Partenariat recherche) \\ \hline
\end{tabularx}
\end{center}
\emph{Remarquez la porosité : l'action 2 est sanitaire dans le fait, diplomatique dans l'intention politique. L'action 4 est économique dans l'instrument, sanitaire dans l'objectif.}
\end{exoresolu}

\courssub{Synthèse visuelle : la coopération camerounaise par aspects}

\begin{popfigure}
\begin{tikzpicture}
\begin{scope}[scale=0.9]
% Données approximatives basées sur l'intensité/visibilité des flux (non monétaires strictes)
% Politique: 15, Economique: 35, Culturel: 10, Militaire: 15, Sanitaire/Env: 25
\def\angleStart{90}
\foreach \label/\value/\color in {
    Économique \& Commercial/35/popblue,
    Sanitaire \& Environnemental/25/popgreen,
    Militaire \& Sécuritaire/15/poporange,
    Politique \& Diplomatique/15/poppurple,
    Culturel \& Scientifique/10/popgold}{
    \pgfmathsetmacro{\angleEnd}{\angleStart - \value*3.6}
    \draw[fill=\color!70, draw=white, thick] (0,0) -- (\angleStart:3) arc (\angleStart:\angleEnd:3) -- cycle;
    \pgfmathsetmacro{\midAngle}{(\angleStart + \angleEnd)/2}
    \pgfmathsetmacro{\labelRadius}{3.5}
    \node[font=\small\bfseries, text=\color!80!black] at (\midAngle:\labelRadius) {\label};
    \pgfmathsetmacro{\valRadius}{1.5}
    \node[font=\small, text=popdark, align=center, text width=2.5cm] at (\midAngle:\valRadius) {\value\%};
    \xdef\angleStart{\angleEnd}
}
% Légende centrale
\node[font=\footnotesize, text=popdark, align=center, text width=4cm] at (0,0) {Répartition qualitative\\des engagements\\de coopération\\du Cameroun};
\end{scope}
% Légende bas
\begin{scope}[yshift=-5.5cm]
\foreach \i/\color/\txt in {
    1/popblue/Économique \& Commercial,
    2/popgreen/Sanitaire \& Environnemental,
    3/poporange/Militaire \& Sécuritaire,
    4/poppurple/Politique \& Diplomatique,
    5/popgold/Culturel \& Scientifique}{
    \pgfmathsetmacro{\y}{-\i*0.6}
    \fill[\color!70] (0,\y) rectangle (0.5,\y+0.4);
    \node[anchor=west, font=\small] at (0.7,\y+0.2) {\txt};
}
\end{scope}
\end{tikzpicture}
\legende{Diagramme circulaire : répartition qualitative estimée des efforts de coopération internationale du Cameroun par aspect (basé sur volume financier, effectifs engagés, visibilité diplomatique et nombre d'accords actifs). L'aspect économique domine par les flux financiers (IDE, APD, Dette), suivi de près par le sanitaire/environnemental (urgence climatique et pandémies).}
\end{popfigure}

\courssub{Les formes de la coopération internationale}

Au-delà des aspects sectoriels, la coopération se distingue par sa \cle{forme juridique et institutionnelle}, qui détermine les modalités de décision, de financement et de mise en œuvre.

\begin{definition}[Formes de la coopération internationale]
On distingue traditionnellement cinq formes principales, souvent combinées dans la pratique camerounaise :
\begin{enumerate}
    \item \textbf{La coopération bilatérale :} Relation directe entre deux États (ex. : Cameroun-France, Cameroun-Chine, Cameroun-USA). Cadre : Commission mixte, Accord-cadre. Avantages : souplesse, adaptation aux besoins précis, visibilité politique.
    \item \textbf{La coopération multilatérale :} Relation via une organisation internationale (ONU, UA, UE, BAD, BM, FMI, OMS, UNESCO). Cadre : Programmes pays, Stratégies d'assistance. Avantages : mutualisation des ressources, normes universelles, légitimité, effet de levier.
    \item \textbf{La coopération décentralisée :} Relation entre collectivités territoriales (Régions, Communes). Ex. : Jumelage Yaoundé-Montreuil, Région du Littoral - Région française. Avantages : proximité, appropriation locale, développement territorial intégré.
    \item \textbf{La coopération non gouvernementale / humanitaire :} Mise en œuvre par des ONG internationales (MSF, CARE, Plan International, Croix-Rouge) ou nationales. Cadre : Convention de siège, accord de partenariat avec l'État hôte. Avantages : réactivité, accès zones difficiles, innovation sociale.
    \item \textbf{La coopération Sud-Sud et triangulaire :} Échanges entre pays en développement (ex. : Cameroun-Brésil agriculture, Cameroun-Maroc formation, Cameroun-Rwanda expérience post-conflit). \textbf{Triangulaire :} Un pays du Nord/émergent finance, un pays du Sud apporte l'expertise, un troisième pays bénéficie (ex. : Appui Japon/FAO pour la riziculture au Cameroun).
\end{enumerate}
\end{definition}

\begin{exemplebox}[La coopération triangulaire en action : Projet rizicole Ndop (Nord-Ouest)]
\begin{itemize}
    \item \textbf{Acteurs :} Japon (financement via JICA), FAO (expertise technique multilatérale), Cameroun (bénéficiaire, mise en œuvre via MINADER/UNVDA).
    \item \textbf{Mécanisme :} Transfert de variétés améliorées (NERICA),'aménagement bas-fonds, formation riziculteurs.
    \item \textbf{Intérêt :} Combinaison financement Nord + expertise technique adaptée Sud + appropriation nationale.
\end{itemize}
\end{exemplebox}

\begin{methode}[Rédiger une fiche de synthèse sur un partenariat bilatéral (ex. Bac / Probatoire)]
Pour structurer votre analyse d'un partenariat (ex. Cameroun-France, Cameroun-Chine, Cameroun-USA) :
\begin{enumerate}
    \item \textbf{En-tête :} Nom du partenaire, date établissement relations, cadre juridique (accord-cadre, commission mixte).
    \item \textbf{Tableau par aspects :} 5 colonnes (Politique, Économique, Culturel, Militaire, Sanitaire/Environnemental). Remplir avec 1 à 2 actions concrètes, chiffrées, datées par case.
    \item \textbf{Analyse des formes :} Identifier la part du bilatéral, multilatéral, décentralisé, ONG, Sud-Sud dans le portefeuille global.
    \item \textbf{Bilan critique :} Équilibre ? Réciprocité ? Alignement sur la SND30 (Stratégie Nationale de Développement 2020-2030) ? Points de friction (droits humains, gouvernance, dette, conditions politiques) ?
    \item \textbf{Conclusion :} Intérêt stratégique pour le Cameroun (diversification, levier de développement, sécurité).
\end{enumerate}
\end{methode}

\courssub{Dossier 1 : L'assistance humanitaire}

L'assistance humanitaire est la manifestation la plus directe de la solidarité internationale. Elle constitue un \textbf{devoir moral} et une \textbf{obligation juridique} (Droit International Humanitaire - DIH).

\begin{definition}[Assistance humanitaire]
L'\textbf{assistance humanitaire} est l'ensemble des actions d'urgence et de reconstruction menées par des acteurs étatiques et non étatiques pour sauver des vies, soulager les souffrances et préserver la dignité humaine face aux crises (conflits, catastrophes naturelles, épidémies), dans le respect des principes humanitaires.
\end{definition}

\begin{propriete}[Les quatre principes humanitaires fondamentaux (Résolution ONU 46/182, 1991)]
\begin{itemize}
    \item \textbf{Humanité :} Sauver des vies, soulager souffrances, où que l'on soit.
    \item \textbf{Neutralité :} Ne pas prendre parti dans les hostilités ou les controverses politiques, raciales, religieuses, idéologiques.
    \item \textbf{Impartialité :} Agir uniquement selon les besoins, sans discrimination (priorité aux plus vulnérables).
    \item \textbf{Indépendance :} Autonomie des objectifs humanitaires par rapport aux objectifs politiques, économiques, militaires ou autres.
\end{itemize}
\end{propriete}

\begin{exemplebox}[Contexte humanitaire au Cameroun : Trois crises simultanées (2024)]
\begin{center}
\begin{tabularx}{\linewidth}{|l|X|X|X|}\hline
\textbf{Crise} & \textbf{Régions concernées} & \textbf{Populations affectées (estim. OCHA)} & \textbf{Acteurs clés} \\ \hline
\textbf{Extrême-Nord (Boko Haram)} & Extrême-Nord & \textasciitilde 1,2 million (PDI, réfugiés nigérians, retournés) & HCR, OCHA, PAM, UNICEF, ICRC, MSF, ONG locales \\ \hline
\textbf{Crise anglophone (NOSO)} & Nord-Ouest, Sud-Ouest & \textasciitilde 600 000 PDI, 70 000 réfugiés au Nigeria & OCHA, HCR, ONG internationales (NRC, DRC), Croix-Rouge camerounaise \\ \hline
\textbf{Réfugiés Centrafricains} & Est, Adamaoua, Nord & \textasciitilde 350 000 réfugiés + communautés d'accueil & HCR, PAM, UNICEF, ONG, Gouvernement (MINAS, MINSANTE) \\ \hline
\end{tabularx}
\end{center}
\legende{PDI = Personnes Déplacées Internement. Source : OCHA Cameroun, Aperçu des besoins humanitaires 2024.}
\end{exemplebox}

\begin{attention}[Piège : Confondre aide humanitaire et aide au développement]
\begin{itemize}
    \item \textbf{Aide humanitaire (Urgence) :} Court terme, sauvetage, gratuité totale, principes stricts (neutralité/impartialité), logistique lourde (avions, camions, kits). Ex. : Distribution vivres aux PDI de Mora.
    \item \textbf{Aide au développement (Structural) :} Long terme, renforcement capacités, infrastructures, gouvernance, co-financement possible, conditionnalités politiques. Ex. : Projet d'adduction d'eau potable sur 5 ans (AFD).
    \item \textbf{Nexus Humanitaire-Développement-Paix (Triple Nexus) :} Approche actuelle (OCDE, ONU) pour briser les silos : lier l'urgence à la résilience et à la prévention des conflits. Ex. : Programme « Résilience Extrême-Nord » (BM/UE) : filets sociaux + appui agricole + dialogue communautaire.
\end{itemize}
\end{attention}

\begin{exoresolu}[Analyse d'une intervention humanitaire : La réponse à l'afflux de réfugiés centrafricains dans l'Est (2014-2024)]
\textbf{Document :} « Le Cameroun maintient sa politique de porte ouverte. Le HCR et ses partenaires appuient le gouvernement pour l'enregistrement, la protection (violences basées sur le genre, protection de l'enfance), l'éducation (intégration écoles publiques), la santé (intégration centres de santé) et les moyens d'existence (agriculture, formation pro). »
\textbf{Questions :}
\begin{enumerate}
    \item Identifiez les formes de coopération mobilisées.
    \item Citez deux principes humanitaires illustrés par l'intégration dans les services publics (écoles, santé).
    \item Quel est le rôle de l'État camerounais dans cette architecture ?
\end{enumerate}
\tcblower
\textbf{Solution :}
\begin{enumerate}
    \item \textbf{Multilatérale} (HCR, PAM, UNICEF, OMS, OIM), \textbf{Bilatérale} (financements directs pays donateurs via HCR/ONG), \textbf{Non gouvernementale} (MSF, Plan International, Caritas, Croix-Rouge camerounaise), \textbf{Décentralisée} (Communes d'accueil impliquées dans l'intégration foncière/scolaire).
    \item \textbf{Impartialité} (accès aux services pour réfugiés et populations hôtes sans distinction) et \textbf{Humanité} (préservation de la dignité par l'autonomie, pas seulement l'assistance).
    \item \textbf{Responsable premier :} Garant de la protection (loi 2005 sur le statut de réfugié), coordination via MINAS/MINREX, sécurité des camps/sites, délivrance documents d'identité, inclusion dans les plans de développement locaux (PDL).
\end{enumerate}
\end{exoresolu}

\courssub{Les principes de la diplomatie camerounaise}

La diplomatie camerounaise, définie par le Préambule de la Constitution de 1996 et le \textbf{Livre Blanc de la Politique Extérieure (2007, actualisé)}, s'inscrit dans la continuité de l'héritage de l'indépendance (Ahmadou Ahidjo, Paul Biya). Elle vise à garantir la \cle{souveraineté}, la \cle{sécurité} et le \cle{développement}.

\begin{propriete}[Les piliers de la diplomatie camerounaise]
\begin{enumerate}
    \item \textbf{Le non-alignement et la diversification des partenariats :} Refus de tout blocage idéologique (Guerre froide). Politique « tous azimuts » : relations actives avec l'Ouest (France, UE, USA), l'Est (Chine, Russie), le Sud (Brésil, Inde, Turquie, Maroc, Afrique du Sud) et les organisations multilatérales. Objectif : maximiser les marges de manœuvre, éviter la dépendance unique.
    \item \textbf{Le bon voisinage et la résolution pacifique des différends :} Privilégier le dialogue, la médiation, le recours à la justice internationale (CIJ). Ex. : Affaire de la péninsule de Bakassi (CIJ 2002, Commission mixte Cameroun-Nigéria, retrait pacifique 2006-2008). Différend frontalier avec le Nigeria (Lac Tchad), avec la Guinée équatoriale (zone maritime) gérés par commissions mixtes.
    \item \textbf{L'intégration régionale et continentale (Panafricanisme) :} Moteur de la CEMAC (monnaie commune, marché commun), de la CEEAC (sécurité, marché), de l'UA (Agenda 2063, ZLECAf, CID). Cameroun = « Charnière » Afrique centrale / Afrique de l'Ouest.
    \item \textbf{Le respect du droit international et des engagements :} Respect des traités (Convention de Vienne, Charte ONU, DIH, Droits de l'Homme, Accords de Paris sur le climat). Participation aux OMP (Casques bleus).
    \item \textbf{La diplomatie économique au service de la SND30 :} L'action extérieure est un levier pour l'émergence (SND30). Priorités : attraction IDE, accès aux marchés (ZLECAf, AGOA, EPA UE), transfert technologies, mobilisation diaspora (transferts \textasciitilde 300 Mds FCFA/an), promotion « Made in Cameroon ».
    \item \textbf{La protection des ressortissants à l'étranger :} Assistance consulaire, rapatriements (crise Libye 2011, Ukraine 2022, Soudan 2023), lutte contre la traite des personnes.
\end{enumerate}
\end{propriete}

\begin{savaistu}[Le « Livre Blanc » de la politique extérieure du Cameroun (2007)]
Document de référence officiel qui formalise la doctrine. Il structure l'action autour de 5 axes : (1) Paix et Sécurité, (2) Intégration sous-régionale/régionale, (3) Coopération bilatérale diversifiée, (4) Coopération multilatérale renforcée, (5) Diplomatie économique et culturelle. Il est enseigné à l'IRIC (Institut des Relations Internationales du Cameroun).
\end{savaistu}

\begin{exemplebox}[Illustration : La gestion de la crise de Bakassi (2002-2008)]
\begin{itemize}
    \item \textbf{Contexte :} Arrêt CIJ (10 oct. 2002) attribuant la péninsule au Cameroun. Tensions vives, populations nigérianes majoritaires sur place.
    \item \textbf{Diplomatie camerounaise :} (1) \textbf{Legalité :} Respect strict de l'arrêt. (2) \textbf{Négociation :} Accord de Greentree (2006) sous égide ONU (Kofi Annan) organisant le transfert progressif. (3) \textbf{Humanitaire :} Protection des populations nigérianes restées (droit de rester, nationalité camerounaise optionnelle). (4) \textbf{Sécurisation :} Déploiement BIR/Gendarmerie sans exactions majeures.
    \item \textbf{Bilan :} Modèle international de résolution pacifique. Le Cameroun a gagné la souveraineté territoriale sans guerre.
\end{itemize}
\end{exemplebox}

\courssub{TP : Cartographier et argumenter l'action internationale du Cameroun}

\begin{methode}[Réaliser une carte d'analyse géopolitique (Croquis commenté) — Exigence Bac/Probatoire]
\begin{enumerate}
    \item \textbf{Choix du thème :} Ex. « Les partenariats stratégiques du Cameroun », « Flux d'aide humanitaire dans l'Extrême-Nord », « Projets d'infrastructures chinois », « Contribution aux OMP ».
    \item \textbf{Fond de carte :} Contour simplifié du Cameroun (10 régions), pays limitrophes, capitale Yaoundé, métropole économique Douala, ports (Kribi, Douala), axes routiers structurants (N1, N2, N3, N10).
    \item \textbf{Sélection de l'information (3 à 5 items max) :} Chiffres clés, localisations précises, logos partenaires.
    \item \textbf{Sémiologie graphique (Légende organisée) :}
    \begin{itemize}
        \item \textbf{Points (proportionnels) :} Montant investissements, effectifs troupes, volume échanges.
        \item   \textbf{Flèches (épaisseur/couleur) :} Flux commerciaux, flux réfugiés, axes de coopération (ex. : route Yaoundé-Nsimalen, pipeline Tchad-Cameroun).
        \item   \textbf{Surfaces (trames/couleurs) :} Zones d'influence (ex. : zone CEMAC, zone d'insécurité Boko Haram), projets agricoles (Ndop, Mbé).
        \item   \textbf{Signes ponctuels :} Ambassades, bases militaires, sièges OI, sites UNESCO.
    \end{itemize}
    \item \textbf{Titre, Légende, Échelle, Orientation, Source, Auteur/Date.}
    \item \textbf{Commentaire argumenté (15-20 lignes) :} Problématique \rightarrow Description spatiale \rightarrow Analyse des logiques (pourquoi là ? quels intérêts ?) \rightarrow Conclusion prospective.
\end{enumerate}
\end{methode}

\begin{exercice}[Sujet de TP noté : « Le Cameroun, carrefour de coopérations en Afrique centrale »]
\textbf{Consigne :} Sur un fond de carte du Cameroun et de son environnement immédiat (Nigéria, Tchad, RCA, Congo, Gabon, Guinée équatoriale, Golfe de Guinée), réaliser un croquis commenté mettant en évidence la multiplicité des formes et aspects de la coopération internationale dont bénéficie le Cameroun.
\textbf{Contraintes :}
\begin{enumerate}
    \item Représenter au minimum : 2 projets chinois (infrastructures), 1 projet français/UE (eau/santé/gouvernance), 1 projet multilatéral (BAD/BM transport/énergie), 1 zone d'action humanitaire (Extrême-Nord ou Est), 1 base/école militaire régionale (EMPABB Douala, FMM Extrême-Nord), le siège de la CEMAC (Bangui) et de la CEEAC (Libreville) par flèches de rattachement.
    \item Utiliser une sémiologie rigoureuse (figés, flèches, surfaces, signes).
    \item Rédiger un commentaire structuré (Introduction \rightarrow Analyse par forme/aspect \rightarrow Synthèse sur la diversification) d'une demi-page.
\end{enumerate}
\end{exercice}

\begin{corrige}[Éléments attendus pour le croquis : « Le Cameroun, carrefour de coopérations »]
\textbf{Fond de carte :} Contour Cameroun + 7 voisins + Golfe de Guinée. Villes : Yaoundé, Douala, Kribi, Maroua, Garoua, Bertoua, Ngaoundéré, Bafoussam, Bamenda, Buea.
\textbf{Légende suggérée :}
\begin{itemize}
    \item \textbf{Flèches épaisses \textcolor{popblue}{BLEUES} (Économique/IDE) :} Douala $\leftarrow$ Chine (Port Kribi, Barrages) ; Yaoundé $\leftarrow$ France/UE (Assainissement, Budget) ; Douala $\leftarrow$ BAD/BM (Routes, Énergie).
    \item \textbf{Flèches \textcolor{popgreen}{VERTES} (Sanitaire/Environnement) :} Bassin du Congo (trame verte) $\rightarrow$ COMIFAC/REDD+ ; Extrême-Nord $\leftarrow$ OMS/UNICEF/HCR (Urgence) ; Est $\leftarrow$ HCR/PAM (Réfugiés RCA).
    \item \textbf{Flèches \textcolor{poporange}{ORANGES} (Militaire/Sécuritaire) :} Extrême-Nord $\rightarrow$ FMM (Secteur 1) ; Douala $\rightarrow$ EMPABB (Formation OMP) ; Yaoundé $\rightarrow$ EMIA/École de Guerre (Partenariat France/USA).
    \item \textbf{Points \textcolor{poppurple}{VIOLETS} (Diplomatie/Politique) :} Yaoundé (Siège MINREX, Ambassades) ; Bangui (Siège CEMAC - flèche rattachement) ; Libreville (Siège CEEAC - flèche rattachement) ; N'Djaména (QG FMM - flèche rattachement).
    \item \textbf{Signes \textcolor{popgold}{OR} (Culturel/Scientifique) :} Yaoundé/Douala (Instituts Français, Confucius, Goethe) ; Universités (Projets IRD/CIRAD/CNRS).
\end{itemize}
\textbf{Commentaire type :} Le croquis révèle une \textbf{diplomatie tous azimuts}. L'aspect \textbf{économique} (bleu) structure le territoire (corridor Douala-Kribi-Ngaoundéré, barrages). L'aspect \textbf{humanitaire/sécuritaire} (vert/orange) concentre les flux aux frontières (Extrême-Nord, Est), illustrant la vulnérabilité régionale. La \textbf{diplomatie multilatérale} (violet) ancre le Cameroun comme hub institutionnel (CEMAC, CEEAC, FMM). La \textbf{diversification des partenaires} (Chine, France, UE, BM, BAD, ONU, Pays du Golfe) assure l'autonomie stratégique, pilier de la SND30.
\end{corrige}

\begin{aretenir}[Ce qu'il faut retenir pour l'examen (Probatoire / Baccalauréat Technique)]
\begin{itemize}
    \item \textbf{Relations Internationales $\neq$ Coopération :} Les RI incluent le conflit ; la coopération est le volet solidaire, contractuel, gagnant-gagnant.
    \item \textbf{5 Aspects sectoriels :} Politique/Diplomatique, Économique/Commercial, Culturel/Scientifique, Militaire/Sécuritaire, Sanitaire/Environnemental. \textbf{Un projet réel en croise toujours plusieurs.}
    \item \textbf{5 Formes institutionnelles :} Bilatérale, Multilatérale, Décentralisée, ONG/Humanitaire, Sud-Sud/Triangulaire. Savoir classifier un accord (ex. Pêche UE = Multilatérale + Bilatérale États membres).
    \item \textbf{Dossier Humanitaire :} 4 principes (Humanité, Neutralité, Impartialité, Indépendance). Contexte Cameroun : 3 crises (Extrême-Nord, NOSO, Est). Nexus Urgence-Développement-Paix.
    \item \textbf{Principes Diplomatie Camerounaise :} Non-alignement/Diversification, Bon voisinage/Règlement pacifique (Bakassi), Intégration panafricaine (CEMAC/CEEAC/ZLECAf), Respect Droit international, Diplomatie économique (SND30), Protection diaspora.
    \item \textbf{Maîtriser 3 exemples chiffrés/datés par aspect} (ex. : Kribi 600 Mds, Casques bleus 700-1000 hommes, Vaccins COVAX 391 200 doses, APD France 500-800 Mds/an, CDN -35\%).
    \item \textbf{Savoir faire :} Analyser un document (texte, tableau, carte), rédiger une fiche pays/partenariat, réaliser un croquis commenté avec légende hiérarchisée et commentaire argumenté.
\end{itemize}
\end{aretenir}

\coursec{Les formes de la coopération internationale}

Dans la section précédente, nous avons vu que la coopération internationale peut prendre des \emph{aspects} variés : économique, culturel, militaire, scientifique, sanitaire. Mais une question reste en suspens : \emph{avec qui} le Cameroun coopère-t-il, et \emph{selon quelles configurations} ? Coopérer seul à seul avec la France n'a pas la même signification que coopérer avec 193 États à l'ONU, ou avec ses cinq voisins de la CEMAC. C'est précisément l'objet de cette section : classer la coopération internationale selon ses \cle{formes}, c'est-à-dire selon le nombre de partenaires engagés et selon leur niveau de développement.

\courssub{La coopération bilatérale}

\textbf{Intuition.} Imagine un apprenti mécanicien de Douala qui signe un partenariat avec un seul maître artisan allemand : échange de savoir-faire, fourniture d'outillage, formation. C'est un lien direct, de personne à personne. À l'échelle des États, c'est exactement cela la coopération bilatérale.

\begin{definition}[Coopération bilatérale et coopération multilatérale]
La \cle{coopération bilatérale} est une forme de coopération qui associe \textbf{deux États} (ou un État et une organisation internationale) liés par des accords directs. La \cle{coopération multilatérale} associe \textbf{plusieurs États} (au moins trois) réunis au sein d'une organisation internationale ou d'un cadre commun.
\end{definition}

\textbf{Énoncé rigoureux.} La coopération bilatérale repose sur des \emph{accords}, des \emph{conventions} ou des \emph{traités} signés directement entre deux partenaires. Elle est pilotée par les ambassades et les agences nationales de coopération. Elle porte sur l'aide financière, l'assistance technique, les investissements, la formation ou la défense.

\textbf{Exemples camerounais.}

\begin{itemize}
\item \textbf{Cameroun--France} : partenariat historique (accords de coopération signés dès l'indépendance en 1960), aide au développement via l'AFD, coopération militaire et culturelle, enseignement du français.
\item \textbf{Cameroun--Chine} : financement et construction d'infrastructures (stade de Olembé, barrages, routes, hôpital de référence de Yaoundé), prêts et investissements chinois croissants depuis les années 2000.
\item \textbf{Cameroun--Allemagne} : coopération technique via la GIZ (formation professionnelle, gouvernance locale, énergie), très utile pour les filières techniques.
\end{itemize}

\begin{exemplebox}[Un accord bilatéral concret]
En 2022, le Cameroun et la Chine signent des conventions de financement pour la construction de routes et d'infrastructures sportives. Deux États, un accord direct : c'est de la coopération \cle{bilatérale}. L'avantage : la négociation est rapide et ciblée. La limite : le Cameroun peut devenir dépendant d'un seul partenaire et s'endetter fortement envers lui.
\end{exemplebox}

\courssub{La coopération multilatérale}

\textbf{Intuition.} Reprends l'apprenti mécanicien : au lieu d'un seul maître, il adhère à une \emph{association nationale des artisans} qui fixe des règles, offre des formations et défend tous ses membres. Le lien n'est plus direct mais passe par une organisation. C'est la coopération multilatérale.

\textbf{Énoncé rigoureux.} La coopération multilatérale s'exerce dans le cadre d'\emph{organisations internationales} qui fixent des règles communes, collectent des contributions et redistribuent aides, financements ou expertise à leurs membres.

\textbf{Exemples camerounais.}

\begin{itemize}
\item \textbf{Cameroun--ONU} : le Cameroun, membre depuis le 20 septembre 1960, bénéficie des programmes du PNUD (développement), de l'UNICEF (enfance), du HCR (réfugiés) et de l'OMS (santé).
\item \textbf{Cameroun--Union africaine (UA)} : organisation panafricaine de 55 États, œuvrant pour la paix, l'intégration et le développement du continent.
\item \textbf{Cameroun--CEMAC} : Communauté Économique et Monétaire de l'Afrique Centrale, cadre d'intégration économique régionale.
\item \textbf{Cameroun--Commonwealth} : le Cameroun, pays bilingue, est membre du Commonwealth depuis novembre 1995, ce qui lui ouvre des programmes éducatifs et commerciaux anglophones.
\end{itemize}

\begin{exemplebox}[La CEMAC, un exemple chiffré]
La \cle{CEMAC} regroupe \textbf{6 États} : le Cameroun, le Tchad, la Centrafrique, le Gabon, la Guinée équatoriale et le Congo. Ces pays partagent une monnaie commune, le \textbf{franc CFA} (zone BEAC), une banque centrale commune et un marché régional. Le Cameroun y est le premier poids économique.
\end{exemplebox}

\textbf{Avantages et limites.} Le multilatéralisme offre au Cameroun une protection collective, des financements importants et une voix dans les grandes décisions mondiales. Mais les décisions y sont lentes, les aides souvent conditionnées, et le poids d'un seul État reste faible face aux grandes puissances.

\courssub{La coopération Nord-Sud}

\begin{definition}[Coopération Nord-Sud et coopération Sud-Sud]
La \cle{coopération Nord-Sud} désigne l'aide des pays développés (le « Nord » : Europe, Amérique du Nord, Japon...) vers les pays en développement (le « Sud » : Afrique, une partie de l'Asie et de l'Amérique latine). La \cle{coopération Sud-Sud} désigne les échanges et l'entraide \textbf{entre pays en développement} eux-mêmes.
\end{definition}

\textbf{Intuition.} Le vocabulaire « Nord » et « Sud » est géographiquement approximatif mais économiquement parlant : les pays riches industrialisés se situent surtout dans l'hémisphère Nord, les pays pauvres surtout au Sud. La coopération Nord-Sud est donc un flux d'aide qui descend ; la coopération Sud-Sud est un échange horizontal, entre pairs.

\textbf{Exemples de coopération Nord-Sud avec le Cameroun.}

\begin{itemize}
\item \textbf{AFD} (Agence Française de Développement) : financement de projets d'eau potable, d'électricité rurale, d'agriculture.
\item \textbf{GIZ} (coopération allemande) : formation professionnelle, appui aux filières techniques, énergies renouvelables.
\item \textbf{Coopération américaine} : USAID (santé, lutte contre le paludisme et le VIH), Peace Corps (volontaires dans les lycées et villages).
\end{itemize}

\textbf{Avantages et limites.} La coopération Nord-Sud apporte des capitaux, des technologies et une expertise que le Cameroun ne possède pas encore. Mais elle entretient parfois une \emph{dépendance} : aides conditionnées, endettement, ingérences dans les choix politiques nationaux.

\courssub{La coopération Sud-Sud}

\textbf{Exemples.}

\begin{itemize}
\item \textbf{Cameroun--Brésil} : coopération agricole (culture de la canne à sucre, du manioc), échanges techniques et formation.
\item \textbf{Cameroun--Inde} : bourses d'études, formation en informatique et en médecine, lignes de crédit pour l'équipement industriel.
\item \textbf{Coopération intra-africaine} : échanges commerciaux entre le Cameroun et le Nigeria, le Tchad ou le Gabon ; corridors routiers transafricains.
\end{itemize}

\textbf{Avantages et limites.} La coopération Sud-Sud est fondée sur l'\emph{égalité} entre partenaires et sur des expériences de développement comparables : les solutions brésiliennes ou indiennes sont souvent mieux adaptées au contexte camerounais. Sa limite : les pays du Sud disposent de moyens financiers plus faibles que ceux du Nord.

\begin{attention}[Erreur fréquente]
Ne crois pas que la coopération \cle{Sud-Sud} a \emph{remplacé} la coopération Nord-Sud. Les deux formes \textbf{coexistent} : le Cameroun reçoit à la fois l'aide de la France ou des États-Unis (Nord-Sud) et coopère avec le Brésil, l'Inde ou le Nigeria (Sud-Sud). Dire que l'une a supplanté l'autre est une erreur classique à l'examen.
\end{attention}

\courssub{La coopération régionale et sous-régionale}

\textbf{Intuition.} Avant de commercer avec l'autre bout du monde, on commence par ses voisins : c'est moins cher, plus rapide, et les besoins sont proches. La coopération régionale applique ce bon sens à l'échelle des États.

\textbf{Énoncé rigoureux.} La \cle{coopération régionale} (ou \cle{sous-régionale} pour un ensemble plus restreint) regroupe des États d'un même continent ou d'une même zone géographique pour intégrer leurs économies : libre circulation des biens et des personnes, monnaie commune, politiques sectorielles harmonisées.

\textbf{Les cadres d'intégration du Cameroun.}

\begin{itemize}
\item \textbf{CEMAC} (6 États) : intégration économique et monétaire \emph{sous-régionale}, monnaie commune (franc CFA de la zone BEAC).
\item \textbf{CEEAC} (Communauté Économique des États de l'Afrique Centrale, 11 États) : cadre \emph{régional} plus large incluant l'Angola, le Burundi, le Rwanda, Sao Tomé-et-Principe et la RDC, axé sur la paix, la sécurité et l'intégration.
\item \textbf{ZLECAf} (Zone de Libre-Échange Continentale Africaine, entrée en vigueur le 1er janvier 2021) : vaste marché unique africain visant à supprimer les droits de douane entre pays africains et à stimuler le commerce intra-africain.
\end{itemize}

\begin{popfigure}
\begin{center}
\begin{tikzpicture}[scale=1.0, font=\small]
% Fond océan
\fill[popblueL!60] (-0.5,-0.5) rectangle (11.5,7.5);
% Pays CEEAC non CEMAC (approximatif)
\fill[popgoldL] (5.2,0.2) -- (9.5,0.2) -- (10.8,1.6) -- (10.2,3.4) -- (8.6,3.6) -- (7.8,2.4) -- (6.4,2.2) -- (5.2,1.2) -- cycle; % RDC/Angola bloc sud-est
\fill[popgoldL] (8.8,3.8) -- (10.4,3.8) -- (10.8,5.4) -- (9.2,5.8) -- (8.4,4.8) -- cycle; % Burundi/Rwanda
% CEMAC
\fill[popgreenL] (3.2,4.6) -- (5.4,4.6) -- (5.6,6.4) -- (4.2,7.2) -- (2.8,6.6) -- (2.4,5.4) -- cycle; % Cameroun
\fill[popgreenL] (5.6,4.8) -- (8.4,4.8) -- (8.4,6.6) -- (5.8,6.6) -- cycle; % Tchad
\fill[popgreenL] (5.8,3.0) -- (8.2,3.0) -- (8.2,4.6) -- (5.8,4.6) -- cycle; % Centrafrique
\fill[popgreenL] (1.2,2.6) -- (3.0,2.6) -- (3.2,4.4) -- (1.4,4.4) -- cycle; % Gabon
\fill[popgreenL] (3.2,2.4) -- (5.6,2.4) -- (5.6,4.4) -- (3.4,4.4) -- cycle; % Congo
\fill[popgreenL] (1.6,4.6) -- (2.6,4.6) -- (2.6,5.2) -- (1.6,5.2) -- cycle; % Guinée équatoriale (Rio Muni)
% Contours
\draw[popdark, thick] (3.2,4.6) -- (5.4,4.6) -- (5.6,6.4) -- (4.2,7.2) -- (2.8,6.6) -- (2.4,5.4) -- cycle;
% Noms
\node at (4.0,5.9) {\textbf{Cameroun}};
\node at (7.0,5.7) {Tchad};
\node at (7.0,3.8) {Centrafrique};
\node at (2.1,3.5) {Gabon};
\node at (4.4,3.4) {Congo};
\node[align=center] at (0.9,4.9) {Guinée\\équat.};
\node at (8.0,1.8) {RDC};
\node at (9.6,4.9) {Rwanda, Burundi};
% Capitale
\fill[popink] (3.6,5.6) circle (0.09);
\node[right, popink] at (3.7,5.6) {Yaoundé};
% Nord
\draw[-{Stealth}, very thick, popdark] (0.2,6.4) -- (0.2,7.3);
\node[above] at (0.2,7.3) {\textbf{N}};
% Légende
\fill[popgreenL] (0.2,0.6) rectangle (0.7,1.0);
\node[right] at (0.8,0.8) {Membres de la CEMAC (6 États)};
\fill[popgoldL] (6.0,0.6) rectangle (6.5,1.0);
\node[right] at (6.6,0.8) {Autres membres de la CEEAC};
\end{tikzpicture}
\end{center}
\legende{Carte schématique de l'Afrique centrale : les 6 États de la CEMAC (en vert) autour du Cameroun, au sein de la CEEAC plus large (en jaune). Croquis simplifié, sans précision cartographique ; échelle indicative.}
\end{popfigure}

\textbf{Avantages et limites.} L'intégration régionale élargit les débouchés pour les entreprises camerounaises, facilite les déplacements et renforce la stabilité. Mais le commerce intra-régional reste faible (économies semblables, routes dégradées, lenteurs douanières), et la ZLECAf peine encore à décoller concrètement.

\courssub{Identifier une forme de coopération : méthode et tableau de synthèse}

\begin{methode}[Comment identifier la forme de coopération dans un exemple]
\begin{enumerate}
\item \textbf{Compter les partenaires} : deux États liés directement $\Rightarrow$ coopération \cle{bilatérale} ; trois États ou plus dans une organisation $\Rightarrow$ coopération \cle{multilatérale}.
\item \textbf{Situer le niveau de développement} : aide d'un pays riche vers un pays pauvre $\Rightarrow$ \cle{Nord-Sud} ; échange entre pays en développement $\Rightarrow$ \cle{Sud-Sud}.
\item \textbf{Vérifier la proximité géographique} : partenaires d'une même région organisés en communauté $\Rightarrow$ coopération \cle{régionale ou sous-régionale}.
\item \textbf{Conclure en justifiant} : toujours citer le critère utilisé (nombre de partenaires, niveau de développement, zone géographique).
\end{enumerate}
\end{methode}

\begin{center}
\begin{tabularx}{\linewidth}{|l|X|X|}
\hline
\rowcolor{popblue} \textbf{Forme} & \textbf{Critère d'identification} & \textbf{Exemples camerounais} \\
\hline
Bilatérale & 2 partenaires directs (accords directs) & Cameroun--France, Cameroun--Chine, Cameroun--Allemagne \\
\hline
Multilatérale & Plusieurs États (3+) dans une organisation & ONU, UA, Commonwealth, CEMAC, CEEAC \\
\hline
Nord-Sud & Pays développé $\rightarrow$ Pays en développement & AFD (France), GIZ (Allemagne), USAID (USA) \\
\hline
Sud-Sud & Entre pays en développement (pairs) & Cameroun--Brésil, Cameroun--Inde, Cameroun--Nigeria \\
\hline
Régionale / Sous-régionale & Voisins d'une même zone géographique & CEMAC (sous-région), CEEAC (région), ZLECAf (continent) \\
\hline
\end{tabularx}
\end{center}
\legende{Tableau de synthèse : formes de coopération, critères d'identification et exemples camerounais.}

\begin{exoresolu}[Classer des exemples de coopération]
\textbf{Énoncé.} Pour chacune des situations suivantes, indique la forme de coopération en justifiant : (a) la GIZ finance un centre de formation professionnelle à Bafoussam ; (b) le Cameroun participe à une réunion des 55 États de l'Union africaine ; (c) l'Inde offre des bourses d'études à des étudiants camerounais ; (d) le Cameroun et le Gabon harmonisent leurs politiques douanières dans la CEMAC.
\tcblower
\textbf{Solution détaillée.}
\begin{itemize}
\item (a) Coopération \textbf{Nord-Sud} (et bilatérale) : l'Allemagne, pays développé, aide le Cameroun, pays en développement, via son agence GIZ.
\item (b) Coopération \textbf{multilatérale} : 55 États réunis au sein d'une organisation, l'UA.
\item (c) Coopération \textbf{Sud-Sud} : l'Inde et le Cameroun sont tous deux des pays en développement ; l'échange se fait entre pairs du Sud.
\item (d) Coopération \textbf{régionale (sous-régionale)} : deux États voisins agissent dans le cadre d'une communauté régionale, la CEMAC.
\end{itemize}
Critères utilisés : nombre de partenaires, niveau de développement, proximité géographique.
\end{exoresolu}

\begin{savaistu}[Le Cameroun, « Afrique en miniature » de la coopération]
Le Cameroun est l'un des rares pays africains à cumuler presque toutes les formes de coopération : membre de la Francophonie \emph{et} du Commonwealth, de la CEMAC \emph{et} de la CEEAC, partenaire traditionnel de la France \emph{et} premier partenaire commercial de la Chine en Afrique centrale. Cette position charnière est un atout diplomatique majeur.
\end{savaistu}

\begin{exercice}[TP guidé : construire un tableau comparatif]
À partir des exemples vus dans cette leçon (Cameroun--France, ONU, AFD, Cameroun--Brésil, CEMAC), construis un tableau à quatre colonnes : \emph{Forme de coopération / Partenaires / Avantages pour le Cameroun / Limites}. Rédige ensuite un paragraphe de conclusion répondant à la question : « De quelle forme de coopération le Cameroun tire-t-il le plus profit ? » Justifie ta réponse avec deux arguments.
\end{exercice}

\begin{corrige}[TP guidé]
Éléments de correction : le tableau doit distinguer correctement bilatérale (2 partenaires, ex. Cameroun--France : rapidité, mais dépendance), multilatérale (ONU : financements et protection, mais lenteur et conditions), Nord-Sud (AFD : capitaux et technologie, mais endettement), Sud-Sud (Brésil : égalité et adaptation, mais moyens limités) et régionale (CEMAC : marché élargi, monnaie stable, mais faible commerce interne). La conclusion est libre mais doit être \emph{justifiée} : par exemple, la coopération régionale car elle renforce l'autonomie du Cameroun (marché commun, monnaie stable) tout en limitant la dépendance extérieure.
\end{corrige}

\begin{aretenir}[L'essentiel de la section]
La coopération internationale prend cinq formes : \cle{bilatérale} (deux États), \cle{multilatérale} (plusieurs États dans une organisation), \cle{Nord-Sud} (aide des pays riches), \cle{Sud-Sud} (échanges entre pays en développement) et \cle{régionale} (intégration entre voisins : CEMAC, CEEAC, ZLECAf). Pour identifier une forme, on compte les partenaires, on situe leur niveau de développement et leur zone géographique. Toutes ces formes coexistent dans la diplomatie camerounaise.
\end{aretenir}

\coursec{Dossier 1 : L'assistance humanitaire}

Après avoir examiné les différentes formes que peut prendre la coopération internationale — bilatérale, multilatérale, décentralisée — nous nous penchons désormais sur une modalité particulière, souvent spectaculaire et toujours urgente : l'assistance humanitaire. Si la coopération au développement vise la transformation structurelle sur le long terme, l'action humanitaire répond à l'impératif moral immédiat de sauver des vies et de préserver la dignité humaine face à la catastrophe. Le Cameroun, par sa position géographique au carrefour de l'Afrique centrale et de l'Ouest, en est un théâtre majeur où s'appliquent concrètement les principes de sa diplomatie : hospitalité, non-ingérence, solidarité africaine et respect du droit international humanitaire (DIH).

\courssub{Définition et fondements de l'assistance humanitaire}

\begin{definition}[L'assistance humanitaire]
L'\cle{assistance humanitaire} est l'aide d'urgence (vivres, soins, abris, protection, éducation) apportée aux populations en détresse à la suite de conflits armés, de catastrophes naturelles ou d'épidémies, \textbf{sans discrimination} fondée sur la nationalité, l'origine ethnique, la religion ou l'opinion politique.
\end{definition}

\begin{definition}[Réfugié et Déplacé Interne]
Un \cle{réfugié} est une personne qui, craignant avec raison d'être persécutée du fait de sa race, de sa religion, de sa nationalité, de son appartenance à un certain groupe social ou de ses opinions politiques, se trouve hors du pays dont elle a la nationalité (Convention de Genève, 1951). Un \cle{déplacé interne} (DI) a fui son domicile pour les mêmes raisons mais \textbf{n'a pas franchi de frontière internationale} ; il reste sous la juridiction de son propre État.
\end{definition}

\begin{attention}[Urgence vs Développement]
Ne confondez jamais \textbf{assistance humanitaire} (urgence, survie, court terme, logique de « sauvetage ») et \textbf{aide au développement} (long terme, autonomie, infrastructures, renforcement des capacités). Une même organisation peut faire les deux, mais les logiques, les financements et les cycles de projet diffèrent radicalement.
\end{attention}

\courssub{Les quatre principes directeurs de l'action humanitaire (HNI)}

L'action humanitaire ne s'improvise pas ; elle obéit à un code éthique strict, adopté par la résolution 46/182 de l'Assemblée générale de l'ONU (1991) et porté par le Comité international de la Croix-Rouge (CICR) et les ONG signataires du Code de conduite.

\begin{propriete}[Les 4 piliers de l'action humanitaire (HNI)]
\begin{enumerate}
    \item \textbf{Humanité} : La souffrance humaine doit être soulagée partout où elle se trouve. Le but est de protéger la vie, la santé et d'assurer le respect de l'être humain.
    \item \textbf{Neutralité} : Les acteurs humanitaires ne doivent pas prendre parti dans les hostilités ni s'engager dans des controverses d'ordre politique, racial, religieux ou idéologique.
    \item \textbf{Impartialité} : L'aide est apportée \textbf{uniquement en fonction des besoins}, en donnant la priorité aux cas les plus urgents. Aucune distinction adverse (nationalité, race, sexe, croyance, classe, opinion politique) n'est admise.
    \item \textbf{Indépendance} : L'action humanitaire doit être autonome par rapport aux objectifs politiques, économiques, militaires ou autres que les acteurs (États, bailleurs) pourraient avoir dans les zones où elle est mise en œuvre.
\end{enumerate}
\end{propriete}

\begin{exemplebox}[Application concrète : Le dilemme de l'accès]
Dans le Nord-Ouest/Sud-Ouest (crise anglophone), une ONG médicale soigne des blessés des deux camps (séparatistes et forces de défense) sans les dénoncer. C'est l'application de l'\textbf{impartialité} (soigner selon le besoin médical) et de la \textbf{neutralité} (ne pas juger la cause de la blessure). Si l'ONG acceptait de ne soigner que les civils d'un côté de la ligne de front sous pression militaire, elle violerait ces principes et perdrait sa légitimité humanitaire.
\end{exemplebox}

\courssub{Les acteurs humanitaires au Cameroun : une architecture multipartite}

L'assistance au Cameroun repose sur un écosystème complexe où s'articulent acteurs onusiens, mouvement de la Croix-Rouge, ONG internationales/nationales et l'État.

\begin{center}
\begin{tabularx}{\linewidth}{|l|X|X|}
\hline
\rowcolor{popblueL}
\textbf{Catégorie} & \textbf{Principaux acteurs} & \textbf{Rôles / Mandats spécifiques} \\
\hline
\textbf{Agences ONU} & \textbf{HCR} (Haut-Commissariat aux Réfugiés) & Protection internationale, enregistrement, solutions durables (rapatriement, intégration, réinstallation). \\
& \textbf{PAM} (Programme Alimentaire Mondial) & Assistance alimentaire (distributions vivres, cash/vouchers), logistique humanitaire (UNHAS). \\
& \textbf{UNICEF} & Protection de l'enfance, éducation d'urgence, nutrition, WASH (Eau, Hygiène, Assainissement). \\
& \textbf{OMS} (Organisation Mondiale de la Santé) & Coordination santé, riposte épidémies (choléra, COVID-19, rougeole), appui aux structures de santé. \\
& \textbf{OCHA} (Bureau de la coordination des affaires humanitaires) & Coordination générale, plaidoyer, financement (Fonds humanitaire), cartographie des besoins (HNO/HRP). \\
\hline
\textbf{Mouvement Croix-Rouge} & \textbf{CICR} (Comité International) & Mandat de protection (détention, rétablissement liens familiaux, DIH), assistance aux victimes de conflits armés. \\
& \textbf{Croix-Rouge Camerounaise (CRC)} & Auxiliaire des pouvoirs publics, premiers secours, gestion des corps, sensibilisation communautaire, première ligne locale. \\
\hline
\textbf{ONG Internationales} & \textbf{MSF} (Médecins Sans Frontières), \textbf{NRC} (Norwegian Refugee Council), \textbf{Plan International}, \textbf{IRC}, \textbf{ACF}... & Mise en œuvre opérationnelle directe (santé, nutrition, protection, abris, éducation) souvent dans les zones les plus difficiles d'accès. \\
\hline
\textbf{ONG Nationales / Locales} & \textbf{ALDEPA}, \textbf{CAMENA}, \textbf{ADRA}, réseaux de femmes, organisations de base... & Ancrage communautaire, acceptation locale, interface culturelle, réponse rapide de proximité (localisation de l'aide). \\
\hline
\textbf{État du Cameroun} & \textbf{MINAS} (Affaires Sociales), \textbf{MINSANTE}, \textbf{Protection Civile}, \textbf{Gouvernorats}, \textbf{Communes} & Responsabilité primaire de protection des populations sur son territoire, coordination nationale, sécurité, facilitation d'accès (visas, laisser-passer), politiques d'asile. \\
\hline
\end{tabularx}
\end{center}

\begin{savaistu}[Le Cameroun, pays « double face » humanitaire]
Le Cameroun présente une particularité rare : c'est simultanément un \textbf{pays d'accueil} majeur pour les réfugiés (Nigéria à l'Extrême-Nord, RCA à l'Est) et un \textbf{pays d'origine} de déplacés internes (crise du Bassin du Lac Tchad, crise anglophone au NO/SO). L'État assume sa responsabilité via la loi n°2005/006 du 27 juillet 2005 relative au statut de réfugié, tout en faisant appel à la solidarité internationale pour combler les gaps de capacité.
\end{savaistu}

\courssub{Cartographie de l'intervention humanitaire au Cameroun}

L'action humanitaire se concentre sur quatre axes géographiques majeurs, dictés par les crises.

\begin{popfigure}
\begin{tikzpicture}[scale=0.85]
% Contour très simplifié du Cameroun (polygone approximatif)
\draw[popdark, thick] (0,0) -- (2,0.5) -- (4,0.2) -- (5.5,1) -- (6,3) -- (5.5,5) -- (4.5,6.5) -- (3,7) -- (1.5,6.5) -- (0.5,5.5) -- (-0.5,4) -- (-0.5,2) -- (0,0.5) -- cycle;
% Régions / Zones
% Extrême-Nord (Nord)
\fill[poporange!30] (0.5,0.5) -- (2.5,1) -- (3,2) -- (1.5,2) -- (0.5,1.5) -- cycle;
\node[poporange, font=\bfseries\small] at (1.8,1.3) {Extrême-Nord};
% Nord
\fill[poporange!15] (2,1) -- (4,1.5) -- (4.5,3.5) -- (3,2.5) -- cycle;
\node[font=\small] at (3.5,2.2) {Nord};
% Adamaoua
\fill[popmuted!20] (3.5,2.5) -- (5,3) -- (5,5) -- (3.5,4.5) -- cycle;
\node[font=\small] at (4.3,3.8) {Adamaoua};
% Est (RCA)
\fill[popred!30] (4.5,3) -- (6,3.5) -- (5.5,5.5) -- (4.5,5) -- cycle;
\node[popred, font=\bfseries\small] at (5.2,4.2) {Est};
% Centre / Sud / Littoral (Yaoundé, Douala)
\fill[popgreen!15] (2.5,3.5) -- (4.5,4) -- (4,5.5) -- (2,5.5) -- (1,4.5) -- cycle;
\node[popdark, font=\small] at (3,4.5) {Centre/Sud/Littoral};
% Nord-Ouest / Sud-Ouest (Crise anglophone)
\fill[poppurple!30] (0.5,3.5) -- (2.5,4) -- (2,5.5) -- (0.5,5) -- cycle;
\node[poppurple, font=\bfseries\small] at (1.3,4.5) {NO / SO};
% Villes repères
\node[circle, fill=popdark, inner sep=1.5pt, label=above:\tiny Yaoundé] (Y) at (2.8,4.2) {};
\node[circle, fill=popdark, inner sep=1.5pt, label=right:\tiny Douala] (D) at (3.2,3.2) {};
\node[circle, fill=popdark, inner sep=1.5pt, label=below:\tiny Maroua] (M) at (1.5,1.2) {};
\node[circle, fill=popdark, inner sep=1.5pt, label=left:\tiny Bertoua] (B) at (5.0,3.8) {};
\node[circle, fill=popdark, inner sep=1.5pt, label=above:\tiny Bamenda] (Ba) at (1.0,4.8) {};
\node[circle, fill=popdark, inner sep=1.5pt, label=below:\tiny Buea] (Bu) at (1.8,3.8) {};
% Flèches flux réfugiés
\draw[->, thick, popblue] (0.5,0.2) -- (1.5,1.0) node[midway, above, sloped, font=\tiny, popblue] {Nigérians (Boko Haram)};
\draw[->, thick, popblue] (6.5,4.0) -- (5.2,4.2) node[midway, above, sloped, font=\tiny, popblue] {Centrafricains};
% Nord et Échelle
\node[font=\small, popdark] at (-1.5, 3.5) {N $\uparrow$};
\draw[|-|, popdark] (-1.5, -0.5) -- (-0.5, -0.5) node[midway, below, font=\tiny] {$\sim$ 200 km};
\end{tikzpicture}
\legende{Carte schématique du Cameroun localisant les quatre principales zones d'intervention humanitaire (Extrême-Nord, Est, Nord-Ouest, Sud-Ouest) et les flux de population concernés. Légende : Orange = Zone réfugiés nigérians + DI (Boko Haram) ; Rouge = Zone réfugiés centrafricains ; Violet = Zone déplacés internes (Crise anglophone) ; Bleu = Flux réfugiés transfrontaliers.}
\end{popfigure}

\courssub{Focus terrain : Le camp de Minawao (Extrême-Nord)}

\begin{exemplebox}[Minawao : De l'urgence à la gestion durable]
Ouvert en juillet 2013, le camp de \textbf{Minawao} (département du Mayo-Tsanaga, près de Maroua) est l'archétype de la réponse humanitaire structurée.
\begin{itemize}
    \item \textbf{Population :} Environ \num{70000} réfugiés nigérians (chiffres HCR 2023-2024), fuyant l'insurrection de Boko Haram au Nord-Est du Nigeria.
    \item \textbf{Gouvernance :} Administration camerounaise (Préfet, Délégué MINAS) + HCR (coordination protection) + ONG partenaires (Plan International, IMC, CARE, CRC...).
    \item \textbf{Secteurs couverts :}
    \begin{itemize}
        \item \textbf{Abri :} Transition des tentes d'urgence vers des abris semi-durables (briques de terre stabilisée), construction participative.
        \item \textbf{Education :} Écoles primaires et secondaires intégrées au système camerounais (programme bilingue), centres d'apprentissage accéléré pour les hors-scolarisation longue durée.
        \item \textbf{Santé/Nutrition :} Centre de santé intégré (réfugiés + population hôte), prise en charge de la malnutrition aiguë (CRENAS/CRENI).
        \item \textbf{Moyens d'existence :} Jardins maraîchers, formations professionnelles (couture, menuiserie, informatique), groupements d'épargne-crédit (AVEC).
        \item \textbf{Protection :} Enregistrement biométrique, prévention VBG (violences basées sur le genre), protection de l'enfance (espaces amis des enfants).
    \end{itemize}
    \item \textbf{Relation hôte-réfugiés :} Partage des ressources (eau, bois, pâturages), tensions gérées par des comités de concertation. Le marché hebdomadaire de Minawao est un pôle d'échange économique vital pour les deux populations.
\end{itemize}
\end{exemplebox}

\courssub{Les formes de l'assistance : la « boîte à outils » humanitaire}

L'assistance ne se limite pas à jeter des sacs de riz depuis un camion. Elle suit des standards techniques (Sphère, INEE) et s'adapte au contexte.

\begin{center}
\begin{tabularx}{\linewidth}{|l|X|X|}
\hline
\rowcolor{popblueL}
\textbf{Secteur} & \textbf{Interventions types} & \textbf{Indicateurs de qualité (Standards Sphère/INEE)} \\
\hline
\textbf{Sécurité Alimentaire} & Distributions vivres (céréales, légumineuses, huile, sel), Transferts monétaires (Cash/Vouchers), Appui agricole (semences, outils). & 2100 Kcal/pers/jour ; Diversité alimentaire (score HDDS) ; Accès marché fonctionnel pour le Cash. \\
\hline
\textbf{Abri / Articles Ménagers Essentiels (AME)} & Tentes, kits abri (bâches, corde, outils), kits AME (nattes, couvertures, bidons, ustensiles, moustiquaires). & Surface couverte $\ge$ 3,5 m²/pers ; Intimité, protection climatique ; Distribution dignifiée. \\
\hline
\textbf{Santé / Nutrition} & Cliniques mobiles, appui centres de santé, vaccination (rougeole, polio, choléra), santé reproductive, santé mentale (MHPSS), prise en charge malnutrition (PECMA). & Taux de mortalité < 1/10 000/j ; Couverture vaccinale > 95\% ; Taux de guérison malnutrition > 75\%. \\
\hline
\textbf{EAH (Eau, Assainissement, Hygiène)} & Forages/réhabilitation points d'eau, latrines (familiales/communautaires), promotion hygiène (lavage mains), gestion déchets, distribution savon. & 15 L/eau/pers/jour ; 1 latrine / 20 pers ; Distance point d'eau < 500m ; Pas de défécation à l'air libre. \\
\hline
\textbf{Protection} & Enregistrement/Documentation, Prévention VBG, Protection de l'enfance (espaces sûrs, gestion cas), Rétablissement liens familiaux (RLF), Assistance juridique. & Accès aux documents d'identité ; Mécanismes plaintes (CFM) fonctionnels ; Meilleurs intérêts de l'enfant. \\
\hline
\textbf{Education d'urgence (EiE)} & Écoles temporaires (TLS), kits scolaires, formation enseignants (PSS, pédagogie urgence), programmes accélérés, récréation. & Accès école < 3km ; Ratio élèves/enseignant < 50:1 ; Programme certifié (INEE). \\
\hline
\end{tabularx}
\end{center}

\courssub{Étude de documents guidée : Photographies et Données chiffrées}

\begin{methode}[Analyser une photographie humanitaire (Type Bac/Probatoire)]
Pour exploiter un document iconographique (photo de camp, distribution, centre de santé) :
\begin{enumerate}
    \item \textbf{Identifier le document :} Nature (photo reportage, satellite, drone), source (HCR, ONG, presse), date, lieu précis (légende).
    \item \textbf{Décrire objectivement ce qui est visible :} Plans (avant/arrière-plan), acteurs (hommes/femmes/enfants, humanitaires en gilets, uniformes), objets (bâches, bidons, sacs PAM, registres), environnement (poussière, boue, végétation, bâtiments).
    \item \textbf{Analyser la mise en scène / le cadrage :} Quel message l'auteur veut-il faire passer ? (Détresse ? Résilience ? Efficacité de l'aide ? Tension ? Ordre/Désordre ?).
    \item \textbf{Relier aux connaissances du cours :} Quels principes (HNI) sont illustrés ? Quels secteurs (EAH, Abri, Protection) ? Quels acteurs ? Quel type de population (réfugiés/DI) ? Quels défis apparaissent (surpeuplement, insalubrité, file d'attente) ?
    \item \textbf{Conclure / Argumenter :} Répondre à la question posée en s'appuyant \textbf{uniquement} sur l'image et le cours.
\end{enumerate}
\end{methode}

\begin{exoresolu}[Analyse d'une photographie de distribution au camp de Minawao]
\textbf{Énoncé :} Vous observez une photographie prise en 2022 au camp de Minawao. On y voit une file de femmes, pour la plupart voilées, tenant des cartes de rationnement et des bidons jaunes vides. Elles attendent sous un soleil intense (ombres courtes). Des agents en gilets orange (logo PAM/ONG) vérifient les cartes sur une table sous une tente. Au fond, des rangées d'abris en bâches blanches et en banco. Aucune ombre d'arbre.

\textbf{Question :} En vous appuyant sur le document et vos connaissances, montrez que cette image illustre à la fois l'efficacité de l'assistance humanitaire et la persistance des vulnérabilités structurelles.

\textbf{Solution détaillée :}
\begin{enumerate}
    \item \textbf{Preuves de l'efficacité de l'assistance (Mise en œuvre des principes) :}
    \begin{itemize}
        \item \textit{Impartialité / Humanité :} La distribution cible les femmes (souvent chefs de ménage vulnérables), sans distinction apparente.
        \item \textit{Organisation / Standards :} Présence de cartes de rationnement (enregistrement biométrique HCR), agents identifiés (gilets), lieu abrité (tente), processus ordonné (file). Cela montre une \textbf{logistique maîtrisée} (PAM/Partenaires).
        \item \textit{Secteur EAH / Abri :} Les bidons jaunes (standard 10-20L) et les abris (bâches UNHCR / banco local) prouvent la couverture des besoins de base (eau, abri).
    \end{itemize}
    \item \textbf{Preuves des vulnérabilités structurelles persistantes (Défis) :}
    \begin{itemize}
        \item \textit{Environnement hostile :} Soleil intense, absence d'arbres (défaut d'ombrage, risque coup de chaleur, pénurie bois de chauffe $\rightarrow$ protection/environnement).
        \item \textit{Genre / Charge de soin :} Quasi-exclusivité des femmes dans la file $\rightarrow$ charge domestique genrée, exposition aux risques (VBG sur le trajet, fatigue).
        \item \textit{Abri précaire :} Mélange bâches (urgence) / banco (durable) $\rightarrow$ transition inachevée, dépendance aux distributions, insécurité foncière.
        \item \textit{Dépendance :} Attente passive $\rightarrow$ absence d'autonomie économique (moyens d'existence insuffisants), nécessité de passer du « care » au « cash » ou activités génératrices de revenus (AGR).
    \end{itemize}
    \item \textbf{Synthèse :} L'image capture l'instant « distribution » (réussite logistique humanitaire) mais le décor (soleil, abris précaires, genre) rappelle que l'assistance humanitaire \textbf{ne résout pas les causes profondes} (conflit, pauvreté, changement climatique) et maintient souvent les bénéficiaires dans une dépendance prolongée.
\end{enumerate}
\end{exoresolu}

\courssub{Évolution chiffrée : Réfugiés et Déplacés Internes au Cameroun}

L'analyse des données du HCR et de l'OIM (Matrice de suivi des déplacements - DTM) permet de mesurer l'ampleur et la dynamique des crises.

\begin{popfigure}
\begin{tikzpicture}
\begin{axis}[
    ybar,
    bar width=12pt,
    height=9cm,
    width=\linewidth,
    enlarge x limits=0.25,
    ylabel={Effectifs (en millions)},
    ylabel style={font=\small},
    xlabel={Années / Situations},
    xlabel style={font=\small},
    symbolic x coords={Ref 2015, Ref 2020, Ref 2024, DI 2018, DI 2020, DI 2024},
    xtick=data,
    xticklabel style={rotate=45, anchor=east, font=\tiny},
    ymin=0,
    ymax=1.6,
    legend style={at={(0.5,-0.25)}, anchor=north, legend columns=2, font=\small, draw=popline, fill=popcream},
    ymajorgrids=true,
    grid style=dashed,
]
\addplot[popblue, fill=popblue!60] coordinates {(Ref 2015,0.30) (Ref 2020,0.44) (Ref 2024,0.48) (DI 2018,0) (DI 2020,0) (DI 2024,0)};
\addplot[poporange, fill=poporange!60] coordinates {(Ref 2015,0) (Ref 2020,0) (Ref 2024,0) (DI 2018,0.44) (DI 2020,1.0) (DI 2024,1.1)};
\legend{Réfugiés (HCR), Déplacés Internes (OIM/Gouv)}
\end{axis}
\end{tikzpicture}
\legende{Évolution du nombre de réfugiés (principalement Nigérians et Centrafricains) et de déplacés internes (crises Extrême-Nord et NO/SO) au Cameroun. Sources : HCR Data Portal, OIM DTM, Rapports HNO Cameroun. Note : Les chiffres DI incluent les retournés dans certaines années ; 2024 est une estimation récente.}
\end{popfigure}

\begin{aretenir}[Lecture du graphique]
\begin{itemize}
    \item La population \textbf{réfugiée} est relativement stable (\textasciitilde 450 000 - 480 000), structurelle (prolongation des crises RCA et Boko Haram).
    \item La population \textbf{déplacée interne (DI)} a \textbf{plus que doublé} entre 2018 et 2024, passant de \textasciitilde 440 000 à plus de \textbf{1,1 million}. C'est la crise la plus dynamique et la plus complexe (double foyer : Extrême-Nord + NO/SO).
    \item Le Cameroun gère ainsi une \textbf{crise de déplacement mixte} massive, pesant lourd sur les budgets nationaux et l'architecture humanitaire.
\end{itemize}
\end{aretenir}

\courssub{Les défis majeurs de l'assistance humanitaire au Cameroun}

Malgré l'engagement des acteurs, l'efficacité de l'aide se heurte à des obstacles structurels et conjoncturels.

\begin{propriete}[Les 4 défis critiques]
\begin{enumerate}
    \item \textbf{Le financement humanitaire insuffisant et imprévisible :} Le Plan de réponse humanitaire (HRP) Cameroun est chroniquement sous-financé (\textasciitilde 40-50\% couverts en moyenne). Conséquence : rationnement alimentaire (ex: 70\% de la ration PAM), fermeture de cliniques, réduction des kits abris. La dépendance aux donateurs institutionnels (ECHO, USAID/BHA, FCDO, Allemagne, Japon) crée une volatilité.
    \item \textbf{L'accès humanitaire (Access) :} Contraintes physiques (routes impraticables saison des pluies, ponts coupés), contraintes administratives (visas, laissez-passer, check-points multiples), contraintes sécuritaires (IED/EEI sur axes routiers NO/SO, enlèvements, attaques de convois). Le principe de \textbf{nécessité d'accès sans entrave} (DIH) est souvent violé.
    \item \textbf{La sécurité des humanitaires :} Le Cameroun figure parmi les pays les plus dangereux pour les travailleurs humanitaires (enlèvements, meurtres, vols de véhicules). Cela force au \textit{remote management} (gestion à distance depuis Yaoundé/Douala), réduisant la qualité du monitoring et la reddition de comptes.
    \item \textbf{La coordination et la « localisation » :} Multiplicité des clusters/secteurs, lourdeur bureaucratique (enregistrement ONG, MINEPAT/MINAS), faible leadership des acteurs locaux (ONG nationales, CRC) dans la prise de décision et l'allocation des ressources (objectif Grand Bargain : 25\% financement direct aux locaux non atteint).
\end{enumerate}
\end{propriete}

\begin{attention}[Piège d'examen : « L'État ne fait rien »]
Ne jamais écrire que « l'État est absent ». L'État camerounais assure : la sécurité (FDS), l'administration civile (préfets, sous-préfets), la politique d'asile (loi 2005), l'intégration des réfugiés dans les services publics (santé, éducation), la fourniture de terrains pour les camps/sites. L'aide internationale \textbf{supplémente} l'État, elle ne le substitue pas (sauf cas d'effondrement total, ce qui n'est pas le cas au Cameroun).
\end{attention}

\courssub{Les principes de la diplomatie camerounaise à l'épreuve de l'humanitaire}

L'action humanitaire au Cameroun est l'application concrète des grands principes qui fondent sa politique étrangère, tels que définis par la Constitution et la pratique diplomatique depuis l'indépendance.

\begin{propriete}[Principes directeurs de la diplomatie camerounaise]
\begin{enumerate}
    \item \textbf{Non-ingérence et respect de la souveraineté :} Le Cameroun refuse toute intervention extérieure dans ses affaires intérieures (crise NO/SO) mais assume sa responsabilité de protection (R2P) sur son territoire. L'aide humanitaire internationale est acceptée sur la base du consentement de l'État hôte.
    \item \textbf{Règlement pacifique des différends :} Privilégie le dialogue (Commission mixte Cameroun-Nigéria, Comité de suivi de l'Accord de cessez-le-feu RCA) et le recours aux juridictions internationales (CIJ, CPI) plutôt que la force.
    \item \textbf{Pan-africanisme et Intégration régionale :} Moteur de la CEMAC, de la CEEAC, de la CLC (Commission du Bassin du Lac Tchad). L'accueil des réfugiés centrafricains et nigérians illustre la solidarité africaine (Convention de l'OUA de 1969 sur les réfugiés).
    \item \textbf{Multilatéralisme actif :} Fidélité à l'ONU, à la Francophonie, au Commonwealth. Candidatures aux organes onusiens (Conseil de sécurité, CDH). Respect des traités (DIH, Droit des réfugiés, CIDE, CEDAW).
    \item \textbf{Diplomatie économique et de développement :} L'aide humanitaire est perçue comme un levier de stabilité régionale, condition sine qua non du développement (SND30). Le Nexus Humanitaire-Développement-Paix (HDP Nexus) est l'outil opérationnel de cette vision.
\end{enumerate}
\end{propriete}

\begin{exemplebox}[Illustration : La gestion de la péninsule de Bakassi]
Le règlement pacifique du différend frontalier avec le Nigeria (Arrêt CIJ 2002, Accord de Greentree 2006) a permis d'éviter une guerre. La diplomatie camerounaise a ensuite géré le retour des populations et l'assistance humanitaire aux déplacés (côté camerounais) avec l'appui de l'ONU, illustrant le lien indissociable entre \textbf{règlement juridique}, \textbf{sécurité} et \textbf{action humanitaire}.
\end{exemplebox}

\courssub{Exercices d'application}

\begin{exercice}[Exercice 1 : Qualification juridique et principe]
\textbf{Contexte :} Une ONG internationale souhaite distribuer des vivres dans un village du Nord-Ouest contrôlé par un groupe armé non-étatique (GANE). Le GANE exige que l'ONG embauche 5 de ses combattants comme « agents de sécurité » pour la distribution et que 20\% des vivres soient remis au groupe pour « l'effort de guerre ».
\textbf{Questions :}
\begin{enumerate}
    \item Quels principes humanitaires (HNI) sont violés par ces exigences ? Justifiez.
    \item Quelle doit être la réponse de l'ONG pour rester fidèle à son mandat ?
\end{enumerate}
\end{exercice}

\begin{corrige}[Exercice 1]
\begin{enumerate}
    \item \textbf{Neutralité violée :} Embaucher des combattants et donner des ressources à une partie au conflit équivaut à prendre parti / soutenir l'effort de guerre. \textbf{Indépendance violée :} L'ONG subit une ingérence dans sa gestion (recrutement imposé, prélèvement sur stocks). \textbf{Impartialité menacée :} La distribution risque d'être détournée au profit du GANE et non selon les besoins réels des civils.
    \item L'ONG doit \textbf{refuser catégoriquement} ces conditions. Elle doit négocier l'accès \textit{sans condition} (accès humanitaire fondé sur le DIH). Si le refus entraîne l'impossibilité d'opérer en sécurité, elle doit \textbf{suspendre ou relocaliser} l'activité, documenter l'entrave (rapport OCHA/Cluster) et plaider pour un accès principié. Elle ne doit jamais céder au chantage.
\end{enumerate}
\end{corrige}

\begin{exercice}[Exercice 2 : Analyse de données et argumentation]
\textbf{Document :} Extrait du graphique « Évolution du nombre de réfugiés et de déplacés internes au Cameroun » (ci-dessus).
\textbf{Question :} « Entre 2018 et 2024, la crise des déplacés internes a pris une ampleur inédite au Cameroun. » À l'aide du graphique et de vos connaissances, vérifiez cette affirmation et expliquez les deux causes principales de cette évolution.
\end{exercice}

\begin{corrige}[Exercice 2]
\textbf{Vérification chiffrée :} Le graphique montre une courbe « Déplacés Internes » passant de \textasciitilde 0,44 million (2018) à \textasciitilde 1,1 million (2024), soit une \textbf{multiplication par 2,5} en 6 ans. L'affirmation est \textbf{exacte}.
\textbf{Causes principales :}
\begin{enumerate}
    \item \textbf{Crise anglophone (Nord-Ouest / Sud-Ouest) :} Déclenchée fin 2016, militarisée à partir de 2017/2018. Combats, destructions de villages, « lockdowns », violences ciblées $\rightarrow$ fuite massive des populations vers les chefs-lieux de département, les villes (Bamenda, Buea, Douala, Yaoundé) et la brousse.
    \item \textbf{Insécurité dans le Bassin du Lac Tchad (Extrême-Nord) :} Activité persistante de Boko Haram / ISWAP (attaques, enlèvements, pillages) + opérations militaires + inondations récurrentes $\rightarrow$ déplacements pendulaires et prolongés des populations riveraines du Lac et de la frontière nigériane.
\end{enumerate}
\textbf{Conséquence :} Le Cameroun passe d'une situation « réfugiés majoritaire » (2015) à une situation « déplacés internes majoritaires » (2024), complexifiant la réponse (protection DI vs Réfugiés, juridictions différentes, solutions durables plus difficiles pour les DI).
\end{corrige}

\begin{savaistu}[Le « Nexus » Humanitaire-Développement-Paix (HDP Nexus)]
Face à la chronicité des crises (Minawao existe depuis 2013, crise NO/SO depuis 2016), l'approche purement humanitaire montre ses limites. Le \cle{Nexus HDP} vise à faire travailler ensemble, dès le début de la crise : les humanitaires (urgence), les acteurs du développement (infrastructures, services publics, emploi) et les acteurs de la paix (médiation, cohésion sociale, DDR - Désarmement Démobilisation Réintégration). Au Cameroun, c'est la stratégie du « Plan de Réponse Humanitaire » aligné sur la « Stratégie Nationale de Développement (SND30) » et les plans de relèvement (Plan de Reconstruction NO/SO, Plan d'Urgence Extrême-Nord).
\end{savaistu}

\coursec{Les principes de la diplomatie camerounaise}

Après avoir analysé les mécanismes de l'assistance humanitaire, forme particulière de la coopération internationale mobilisée dans l'urgence, nous nous tournons maintenant vers le cadre permanent qui guide l'action extérieure de notre pays. La diplomatie camerounaise ne se résume pas à des réactions ponctuelles ; elle repose sur une doctrine cohérente, inscrite dans la Constitution et éprouvée par l'histoire. Comprendre ces principes, c'est saisir comment le Cameroun, « Afrique en miniature », défend ses intérêts et promeut la paix sur la scène mondiale.

\courssub{Qu'est-ce que la diplomatie ? Définition et fondement constitutionnel}

\begin{definition}[La diplomatie]
La \cle{diplomatie} est l'art et la pratique de conduire les relations d'un État avec les autres États et les organisations internationales par des moyens pacifiques : négociations, conclusion de traités, représentation permanente ou temporaire. Elle vise à défendre les intérêts nationaux, à préserver la paix et à favoriser la coopération.
\end{definition}

Au Cameroun, cette définition trouve sa source suprême dans le \textbf{Préambule de la Constitution du 18 janvier 1996}. Celui-ci proclame l'attachement du peuple camerounais aux \cle{principes du droit international} tels qu'édictés dans la Charte des Nations Unies : égalité souveraine des États, règlement pacifique des différends, non-ingérence, coopération. La diplomatie camerounaise n'est donc pas une simple habitude administrative ; c'est une \textbf{obligation constitutionnelle}.

\begin{aretenir}[Fondement juridique]
L'article 45 de la Constitution confie au Président de la République la charge de définir la politique de la Nation et de la mettre en œuvre. Il est le chef de la diplomatie : il accrédite les ambassadeurs, ratifie les traités, représente le pays aux sommets. Le Ministère des Relations Extérieures (MINREX) est l'instrument technique d'exécution de cette volonté politique.
\end{aretenir}

\courssub{Les cinq piliers de la doctrine diplomatique camerounaise}

La politique étrangère du Cameroun repose sur cinq principes directeurs, constants depuis l'indépendance, qui forment un tout cohérent. Le schéma ci-dessous les synthétise avec une illustration concrète pour chacun.

\begin{popfigure}
\begin{tikzpicture}[
    node distance=1.2cm and 2.5cm,
    principle/.style={rectangle, rounded corners, draw=popblue, thick, fill=popblueL, minimum width=3.2cm, minimum height=1.8cm, text width=3cm, align=center, font=\small},
    example/.style={rectangle, rounded corners, draw=popgreen, fill=popgreenL, minimum width=3.5cm, minimum height=1.5cm, text width=3.2cm, align=center, font=\small\itshape},
    arrow/.style={-Stealth, thick, popdark}
]
% Central node
\node[principle, fill=popgoldL, draw=popgold, minimum width=4cm, minimum height=2cm, font=\normalsize\bfseries, align=center] (center){Les 5 Principes\\ de la Diplomatie\\ Camerounaise};

% Principle 1
\node[principle, above left=of center, xshift=-1cm, align=center] (p1){\textbf{1. Souveraineté\\ \& Non-ingérence}};
\node[example, left=of p1, xshift=-0.5cm, align=center] (e1){Refus de toute ingérence\\ dans la gestion interne\\ (ex : élections, justice)};

% Principle 2
\node[principle, above right=of center, xshift=1cm, align=center] (p2){\textbf{2. Coexistence pacifique\\ \& Règlement pacifique}};
\node[example, right=of p2, xshift=0.5cm, align=center] (e2){Affaire Bakassi (CIJ 2002,\\ Accord de Greentree 2006)};

% Principle 3
\node[principle, below right=of center, xshift=1cm, align=center] (p3){\textbf{3. Non-alignement\\ \& Ouverture}};
\node[example, right=of p3, xshift=0.5cm, align=center] (e3){Membre du Commonwealth \&\\ de la Francophonie ;\\ partenariats Chine, UE,\\ Russie, Monde arabe};

% Principle 4
\node[principle, below left=of center, xshift=-1cm, align=center] (p4){\textbf{4. Solidarité africaine\\ \& Intégration}};
\node[example, left=of p4, xshift=-0.5cm, align=center] (e4){Moteur de l'UA, CEMAC,\\ CEEAC ; Casques bleus\\ au Mali, en RCA};

% Principle 5
\node[principle, below=of center, yshift=-1.5cm, align=center] (p5){\textbf{5. Coopération pour\\ le Développement}};
\node[example, below=of p5, yshift=-0.5cm, align=center] (e5){Projets structurants :\\ Barrage de Nachtigal,\\ Port de Kribi, routes,\\ santé, éducation};

% Arrows
\draw[arrow] (center.north west) -- (p1.south east);
\draw[arrow] (center.north east) -- (p2.south west);
\draw[arrow] (center.south east) -- (p3.north west);
\draw[arrow] (center.south west) -- (p4.north east);
\draw[arrow] (center.south) -- (p5.north);

\draw[arrow, popblue] (p1.west) -- (e1.east);
\draw[arrow, popblue] (p2.east) -- (e2.west);
\draw[arrow, popblue] (p3.east) -- (e3.west);
\draw[arrow, popblue] (p4.west) -- (e4.east);
\draw[arrow, popblue] (p5.south) -- (e5.north);
\end{tikzpicture}
\legende{Les cinq principes directeurs de la diplomatie camerounaise et leurs illustrations concrètes.}
\end{popfigure}

\courssub{Principe 1 : Respect de la souveraineté et non-ingérence}

\begin{definition}[Non-ingérence]
La \cle{non-ingérence} est le principe selon lequel aucun État ne doit intervenir, directement ou indirectement, pour quelque motif que ce soit, dans les affaires intérieures d'un autre État souverain (politique, économique, social, culturel).
\end{definition}

Ce principe est le corollaire de l'égalité souveraine des États (article 2, paragraphe 1 de la Charte de l'ONU). Pour le Cameroun, cela signifie :
\begin{itemize}
    \item le refus des ingérences étrangères dans son processus électoral, sa justice ou sa gestion administrative ;
    \item l'engagement à ne pas soutenir de mouvements subversifs chez ses voisins ;
    \item la défense de l'intégrité territoriale comme domaine réservé exclusif de l'État.
\end{itemize}

\begin{attention}[Piège à éviter]
La non-ingérence n'interdit pas la \emph{critique} diplomatique ni la \emph{conditionnalité} de l'aide au développement (gouvernance, droits de l'homme). Elle interdit l'\textbf{action coercitive} (intervention militaire, subversion, pression économique déloyale) visant à imposer un choix interne à un État.
\end{attention}

\courssub{Principe 2 : Coexistence pacifique et règlement pacifique des différends}

C'est la marque de fabrique du Cameroun. Face à un litige, le réflexe diplomatique est : \textbf{négociation $\rightarrow$ médiation $\rightarrow$ arbitrage ou justice internationale}. Jamais la force.

\begin{exemplebox}[Le modèle Bakassi : de la Cour à la paix]
Le différend frontalier avec le Nigeria sur la péninsule de Bakassi (richesses halieutiques, pétrole, position stratégique) a duré des décennies. Le Cameroun a saisi la \cle{Cour Internationale de Justice (CIJ)} en 1994.
\begin{itemize}
    \item \textbf{Octobre 2002 :} la CIJ attribue la souveraineté sur Bakassi au Cameroun (fondée notamment sur les traités anglo-allemands de 1913).
    \item \textbf{Juin 2006 (Accord de Greentree) :} sous l'égide de l'ONU (Secrétaire général Kofi Annan), le Nigeria accepte le retrait progressif de ses forces et de son administration. Transfert effectif de l'administration en août 2006.
    \item \textbf{Août 2008 :} retrait total des forces nigérianes. Aucune guerre, aucun mort au combat pour la récupération du territoire.
\end{itemize}
C'est un \textbf{cas d'école} de règlement pacifique (article 33 de la Charte de l'ONU).
\end{exemplebox}

\begin{exoresolu}[Analyse juridique du règlement de Bakassi]
\textbf{Énoncé :} En vous appuyant sur le principe du règlement pacifique des différends, expliquez pourquoi l'affaire Bakassi est considérée comme un succès de la diplomatie camerounaise.
\tcblower
\textbf{Solution détaillée :}
\begin{enumerate}
    \item \textbf{Identification du principe :} le Cameroun a appliqué l'article 33 de la Charte des Nations Unies (négociation, enquête, médiation, conciliation, arbitrage, règlement judiciaire).
    \item \textbf{Choix de la voie judiciaire :} après l'échec des négociations bilatérales, le Cameroun a saisi la CIJ (organe judiciaire principal de l'ONU), acceptant par avance le caractère obligatoire de son arrêt (article 94 de la Charte).
    \item \textbf{Respect du droit :} l'arrêt de 2002 s'est fondé sur les titres juridiques (traité anglo-allemand de 1913), et non sur la force ou sur l'effectivité nigériane.
    \item \textbf{Mise en œuvre négociée :} face à la réticence nigériane, le Cameroun a accepté la médiation du Secrétaire général de l'ONU, d'où l'Accord de Greentree (2006) fixant un calendrier de retrait.
    \item \textbf{Bilan :} récupération du territoire sans effusion de sang, préservation des relations bilatérales (Commission mixte Cameroun-Nigeria), respect du droit international. \textbf{Un modèle pour l'Afrique.}
\end{enumerate}
\end{exoresolu}

\begin{popfigure}
\begin{tikzpicture}[scale=0.9, every node/.style={font=\small}]
% Coastline approximation
\draw[popblue, very thick, fill=popblue!10] (0,0) .. controls (1,0.5) and (3,1) .. (5,0.5) .. controls (6,0) and (7,-0.5) .. (8,0) .. controls (9,0.5) and (10,1) .. (11,0) -- (11,-4) -- (0,-4) -- cycle;

% Bakassi Peninsula
\draw[poporange, very thick, fill=poporange!20] (5,0.5) .. controls (6,1.5) and (7.5,2) .. (8.5,1.5) .. controls (9,1) and (9.5,0) .. (10,0) -- (10,-1) -- (5,-1) -- cycle;
\node[poporange, fill=white, fill opacity=0.8, text opacity=1, rounded corners, inner sep=2pt] at (7,0.5) {\textbf{Péninsule de Bakassi}};

% Border line (CIJ 2002)
\draw[popdark, dashed, thick] (5,0.5) -- (5,-1) node[midway, left] {Frontière terrestre};
\draw[popdark, dashed, thick] (10,0) -- (10,-1) node[midway, right] {Frontière maritime};

% Cities
\node[circle, fill=popdark, inner sep=1.5pt, label=above:{Rio del Rey}] at (3.5,-0.5) {};
\node[circle, fill=popdark, inner sep=1.5pt, label=above:{Bakassi (ville)}] at (7.5,0.2) {};
\node[circle, fill=popdark, inner sep=1.5pt, label=below right:{Calabar (Nigeria)}] at (10.5,-1.5) {};
\node[circle, fill=popdark, inner sep=1.5pt, label=left:{Limbe}] at (2,-2) {};
\node[circle, fill=popdark, inner sep=1.5pt, label=left:{Douala}] at (1,-3) {};

% Country labels
\node[font=\bfseries\large, popblue] at (1.5,-1.5) {CAMEROUN};
\node[font=\bfseries\large, poporange] at (10,-2.5) {NIGERIA};

% Gulf of Guinea
\node[font=\itshape, popblue] at (5.5,-3) {Golfe de Guinée};

% North Arrow
\draw[-Stealth, thick] (10.5,2.5) -- (10.5,3.2) node[above] {N};

% Scale
\draw[|-|, thick] (0.5,-3.5) -- (3.5,-3.5) node[midway, below] {$\approx$ 100 km};
\end{tikzpicture}
\legende{Croquis schématique de la péninsule de Bakassi et de la frontière Cameroun-Nigeria (arrêt CIJ de 2002). La frontière maritime prolonge la frontière terrestre selon le principe d'équidistance. Croquis indicatif, sans prétention à la précision cartographique.}
\end{popfigure}

\courssub{Principe 3 : Non-alignement et ouverture à tous les partenaires}

\begin{definition}[Non-alignement]
Le \cle{non-alignement} est la politique d'un État qui refuse de s'intégrer dans un bloc militaire ou idéologique (hier : OTAN et Pacte de Varsovie ; aujourd'hui : blocs Ouest/Est) pour préserver sa liberté de décision et coopérer avec tous sur la base de l'intérêt mutuel.
\end{definition}

\begin{savaistu}[Double héritage, double appartenance]
Le Cameroun est l'un des rares pays au monde à être membre à part entière du \textbf{Commonwealth} (depuis 1995, héritage britannique) et de la \textbf{Francophonie} (héritage français). Cette originalité historique (réunification de 1961) est un atout diplomatique majeur : le Cameroun sert de \textbf{pont} entre l'Afrique anglophone et francophone et dispose de réseaux diplomatiques doublés.
\end{savaistu}

Concrètement, le Cameroun entretient des relations soutenues avec :
\begin{itemize}
    \item \textbf{l'Occident (UE, France, USA, Allemagne, Canada) :} aide au développement, formation, commerce, sécurité (lutte antiterroriste) ;
    \item \textbf{la Chine :} partenaire infrastructurel majeur (barrages, routes, stades, hôpital de référence), sans conditionnalité politique ;
    \item \textbf{la Russie :} coopération militaire (formation, équipement), énergie, éducation (boursiers) ;
    \item \textbf{le monde arabe (Arabie Saoudite, Qatar, Koweït, Émirats) :} financement de projets (routes, barrages, édifices religieux et universitaires), emploi de travailleurs camerounais.
\end{itemize}

\begin{attention}[Nuance cruciale]
Non-alignement $\neq$ neutralité $\neq$ isolement. La neutralité (comme celle de la Suisse) implique des obligations juridiques strictes en temps de guerre. Le non-alignement est une \textbf{liberté d'appréciation} : le Cameroun vote à l'ONU selon sa conscience et ses intérêts, condamne l'agression lorsqu'il l'estime nécessaire, mais maintient le dialogue avec toutes les parties.
\end{attention}

\courssub{Principe 4 : Solidarité africaine et intégration régionale}

Le Préambule de la Constitution affirme l'attachement du peuple camerounais aux idéaux panafricains et à l'unité africaine. Cela se traduit par un engagement institutionnel fort :

\begin{center}
\begin{tabularx}{\linewidth}{|l|X|X|}\hline
\rowcolor{popblueL}
\textbf{Organisation} & \textbf{Rôle du Cameroun} & \textbf{Retombées concrètes} \\ \hline
\textbf{Union Africaine (UA)} & Forte présence camerounaise dans les organes ; participation active aux sommets des Chefs d'État. & Promotion de la ZLECAf (Zone de Libre-Échange Continentale Africaine) et de l'Agenda 2063. \\ \hline
\textbf{CEMAC} (Communauté Économique et Monétaire de l'Afrique Centrale) & Siège de la BEAC (Banque des États de l'Afrique Centrale) à \textbf{Yaoundé} ; leadership économique du pays dans la sous-région. & Monnaie commune (franc CFA), marché commun, surveillance multilatérale des économies. \\ \hline
\textbf{CEEAC} (Communauté Économique des États de l'Afrique Centrale) & Contribution majeure aux mécanismes de paix et de sécurité (FOMAC) ; intégration physique. & Paix et sécurité en RCA et au Tchad ; corridors routiers transafricains. \\ \hline
\end{tabularx}
\end{center}

\begin{exemplebox}[Casques bleus : la solidarité en action]
Le Cameroun est un \textbf{pourvoyeur de contingents} aux opérations de maintien de la paix de l'ONU (MINUSCA en République centrafricaine, MINUSMA au Mali jusqu'à son retrait, MONUSCO en RDC). Des centaines de militaires et de policiers camerounais servent sous drapeau onusien. C'est la diplomatie par l'action : protéger des civils, stabiliser des États frères, gagner en crédibilité internationale.
\end{exemplebox}

\courssub{Principe 5 : La coopération internationale au service du développement national}

La diplomatie n'est pas une fin en soi ; elle est \textbf{instrumentale}. Le but ultime : le \cle{développement} (émergence à l'horizon 2035, SND30). Cela se traduit par la \textbf{diplomatie économique} :

\begin{itemize}
    \item \textbf{recherche de financements :} négociation d'accords de prêts et de dons (Banque mondiale, FMI, BAD, BEI, Exim Bank of China, fonds saoudien et koweïtien, BADEA) ;
    \item \textbf{attraction des IDE (Investissements Directs Étrangers) :} missions de promotion, ambassades comme antennes commerciales ;
    \item \textbf{transfert de technologie et formation :} bourses d'études (Chine, Russie, France, Allemagne, Maroc, Turquie, Inde), centres de formation professionnelle ;
    \item \textbf{accès aux marchés :} négociation d'accords commerciaux (accords de partenariat avec l'UE, AGOA avec les USA, ZLECAf).
\end{itemize}

\begin{aretenir}[Diplomatie économique : exemples phares]
\begin{itemize}
    \item \textbf{Barrage hydroélectrique de Nachtigal (420 MW) :} partenariat public-privé (État camerounais, EDF, IFC) — diplomatie énergétique.
    \item \textbf{Port en eau profonde de Kribi (phases 1 et 2) :} financement chinois (Exim Bank), exploitation en partenariat international — diplomatie portuaire.
    \item \textbf{Autoroute Yaoundé-Nsimalen, axe Yaoundé-Douala, routes du Nord :} cofinancements multiples (BAD, Banque islamique de développement, Chine, Corée).
\end{itemize}
\end{aretenir}

\courssub{Les instruments de la diplomatie camerounaise}

Comment ces principes s'opèrent-ils concrètement ? Par une architecture institutionnelle et humaine.

\begin{methode}[L'architecture diplomatique camerounaise]
\begin{enumerate}
    \item \textbf{Le décideur suprême :} le Président de la République (Chef de l'État, chef de la diplomatie). Il définit les grandes orientations, nomme les ambassadeurs, ratifie les traités.
    \item \textbf{Le ministère technique :} le \textbf{MINREX} (Ministère des Relations Extérieures). Il conçoit, coordonne et met en œuvre la politique extérieure. Il comprend des directions géographiques (Afrique, Europe, Amériques, Asie-Océanie, Moyen-Orient) et thématiques (affaires juridiques, coopération, affaires économiques, affaires consulaires, protocole).
    \item \textbf{Le réseau de représentation :} ambassades, hauts-commissariats (dans les pays du Commonwealth), missions permanentes (ONU à New York et Genève, UA à Addis-Abeba, UNESCO à Paris, FAO à Rome), consulats généraux. \textbf{Rôle :} représentation, négociation, protection des ressortissants, veille informationnelle, promotion économique.
    \item \textbf{La diplomatie parlementaire :} Assemblée Nationale et Sénat (groupes d'amitié parlementaire, Union interparlementaire, Assemblée parlementaire de la Francophonie).
    \item \textbf{La diplomatie décentralisée :} jumelages entre villes et entre régions (coopération décentralisée).
    \item \textbf{L'instrument militaire :} participation aux opérations de maintien de la paix (casques bleus), accords de coopération de défense, formation d'officiers étrangers dans les écoles militaires camerounaises.
\end{enumerate}
\end{methode}

\courssub{Étude de cas guidée : la résolution pacifique du conflit de Bakassi, modèle de diplomatie camerounaise}

Cette étude de cas synthétise \textbf{tous les principes} vus ci-dessus. Utilisez la grille d'analyse ci-dessous pour structurer votre réponse au Bac.

\begin{experience}[Grille d'analyse du dossier Bakassi — Démarche d'investigation]
\textbf{Documents de base :} arrêt de la CIJ du 10 octobre 2002 ; Accord de Greentree du 12 juin 2006 ; discours officiels lors du transfert de souveraineté (août 2006).

\textbf{Étape 1 : Identifier le différend et les acteurs.}
\begin{itemize}
    \item Objet : souveraineté sur la péninsule de Bakassi (presqu'île pétrolifère et halieutique).
    \item Parties : République du Cameroun contre République fédérale du Nigeria.
    \item Contexte : tensions militaires récurrentes, présence nigériane effective dans la péninsule depuis plusieurs décennies.
\end{itemize}

\textbf{Étape 2 : Qualifier la démarche camerounaise (principe 2).}
\begin{itemize}
    \item Choix de la \textbf{voie judiciaire} (CIJ) et non militaire. Pourquoi ? Respect du droit international, recherche de légitimité internationale, refus d'une escalade coûteuse.
    \item Argumentaire juridique : traité anglo-allemand du 11 mars 1913 (tracé de la frontière), effectivités administratives camerounaises.
\end{itemize}

\textbf{Étape 3 : Analyser la gestion de la victoire juridique (principes 1 et 2).}
\begin{itemize}
    \item Arrêt de 2002 : victoire sur le fond (souveraineté reconnue au Cameroun).
    \item Problème : le Nigeria tarde à se retirer (populations, administration, militaires).
    \item Réaction : \textbf{patience stratégique}, refus de l'ultimatum guerrier, acceptation de la \textbf{médiation onusienne} (Secrétaire général Kofi Annan).
\end{itemize}

\textbf{Étape 4 : Étudier l'Accord de Greentree (principes 2 et 5).}
\begin{itemize}
    \item Mécanisme : retrait progressif en deux phases (administrative puis militaire) sur deux ans (2006-2008).
    \item Garanties : protection des populations nigérianes restées sur place (droits de l'homme), commission mixte de suivi.
    \item Résultat : transfert pacifique de l'administration le 14 août 2006 (cérémonie du drapeau), retrait militaire total en août 2008.
\end{itemize}

\textbf{Étape 5 : Mesurer la portée diplomatique (principes 3, 4 et 5).}
\begin{itemize}
    \item \textbf{Non-alignement et ouverture :} médiation acceptée par les deux parties, soutenue par l'ensemble de la communauté internationale (UE, USA, UA, France, Chine, Russie).
    \item \textbf{Solidarité africaine :} soutien de l'UA ; modèle cité pour d'autres conflits frontaliers africains, dans le respect du principe d'intangibilité des frontières héritées de la colonisation.
    \item \textbf{Développement :} sécurisation de la zone économique exclusive (ZEE) pour l'exploitation pétrolière et gazière, développement de Kribi (port, zone industrielle) en arrière-pays.
\end{itemize}
\end{experience}

\begin{exoresolu}[Sujet type Bac : « La résolution du conflit de Bakassi illustre l'efficacité de la diplomatie camerounaise. » Discutez.]
\textbf{Énoncé :} Construisez une dissertation structurée sur ce sujet.
\tcblower
\textbf{Plan structuré (méthode : principe $\rightarrow$ explication $\rightarrow$ exemple)}

\textbf{Introduction :} accroche (conflit frontalier majeur en Afrique centrale), définition de la diplomatie camerounaise (moyens pacifiques), thèse (un modèle abouti), annonce du plan.

\textbf{I. Le choix du droit comme arme : le principe du règlement pacifique (principe 2).}
\begin{itemize}
    \item Saisine de la CIJ en 1994 (article 33 de la Charte de l'ONU).
    \item Argumentation juridique solide (traité de 1913) face à l'argument nigérian de l'effectivité.
    \item Arrêt de 2002 : victoire juridique indiscutable.
\end{itemize}

\textbf{II. La force de la patience et du dialogue : non-ingérence et coexistence (principes 1 et 2).}
\begin{itemize}
    \item Refus de l'option militaire malgré les provocations ; fermeté sur la souveraineté.
    \item Acceptation de la médiation de l'ONU (Kofi Annan) $\rightarrow$ Accord de Greentree (2006).
    \item Gestion humaine du retrait : protection des populations, commission mixte de suivi.
\end{itemize}

\textbf{III. Un rayonnement diplomatique global (principes 3, 4 et 5).}
\begin{itemize}
    \item \textbf{Non-alignement :} soutien unanime de la communauté internationale, tous horizons confondus.
    \item \textbf{Solidarité africaine :} modèle de l'UA pour la prévention des conflits, respect de l'intangibilité des frontières héritées.
    \item \textbf{Développement :} récupération de la ZEE, base pour l'émergence (Kribi, gaz, pêche).
\end{itemize}

\textbf{Conclusion :} bilan — territoire récupéré, paix préservée, image renforcée, droit respecté. Limite — le défi restant de l'intégration socio-économique des populations de Bakassi (état civil, services publics, développement local). Ouverture — application de la même méthode aux autres différends frontaliers.
\end{exoresolu}

\begin{methode}[Réussir l'exercice de dissertation ou d'analyse de texte au Bac ECM]
Pour justifier une réponse sur la diplomatie camerounaise, \textbf{articulez toujours votre raisonnement en trois temps} :
\begin{enumerate}
    \item \textbf{Énoncer le principe} (vocabulaire précis : « non-ingérence », « coexistence pacifique », « non-alignement », « solidarité africaine », « coopération pour le développement »).
    \item \textbf{Expliquer le mécanisme} (en quoi consiste ce principe ? quelle est sa base juridique ?).
    \item \textbf{Illustrer par un exemple camerounais concret, daté et chiffré si possible} (Bakassi 2002/2006, casques bleus en RCA et au Mali, adhésion au Commonwealth en 1995, barrage de Nachtigal, ZLECAf, etc.).
\end{enumerate}
\textit{Exemple de phrase réussie : « Le principe de non-alignement (1) permet au Cameroun de diversifier ses partenariats sans subir de conditionnalité politique (2) ; ainsi, il coopère simultanément avec la France sur le barrage de Nachtigal et avec la Chine sur le port de Kribi (3). »}
\end{methode}

\begin{exercice}[Entraînement — Sujets probables Probatoire/Bac]
\begin{enumerate}
    \item Définissez la diplomatie et citez son fondement constitutionnel au Cameroun.
    \item Énumérez les cinq principes directeurs de la diplomatie camerounaise. En choisissant \textbf{deux} d'entre eux, montrez, par des exemples précis, comment ils s'appliquent dans la pratique.
    \item « Le règlement de l'affaire Bakassi est un modèle de diplomatie préventive et curative. » Commentez cette assertion en mobilisant les notions de souveraineté, de règlement pacifique et de solidarité africaine.
    \item En quoi la double appartenance du Cameroun (Commonwealth et Francophonie) est-elle un atout pour sa politique de non-alignement ?
    \item Présentez les principaux instruments de la diplomatie camerounaise et expliquez le rôle du MINREX.
\end{enumerate}
\end{exercice}

\begin{corrige}[Exercice n°2 — Éléments de réponse]
\textbf{Les cinq principes :} 1. Souveraineté et non-ingérence. 2. Coexistence pacifique et règlement pacifique des différends. 3. Non-alignement et ouverture. 4. Solidarité africaine et intégration régionale. 5. Coopération pour le développement.

\textbf{Exemple développé pour le principe 2 (règlement pacifique) :} affaire Bakassi. Saisine de la CIJ en 1994 $\rightarrow$ arrêt de 2002 (souveraineté reconnue au Cameroun) $\rightarrow$ médiation de l'ONU $\rightarrow$ Accord de Greentree en 2006 $\rightarrow$ retrait effectif total en 2008. Aucune guerre.

\textbf{Exemple développé pour le principe 4 (solidarité africaine) :} engagement de casques bleus (MINUSCA en RCA, ex-MINUSMA au Mali) ; siège de la BEAC à Yaoundé ; promotion de la ZLECAf ; contribution à la stabilité régionale (RCA, Tchad).
\end{corrige}

\coursec{TP : cartographier et argumenter l'action internationale du Cameroun}

Après avoir analysé les principes qui fondent la diplomatie camerounaise — non-ingérence, règlement pacifique des différends, coopération mutuellement avantageuse, respect du droit international —, il est indispensable de passer de la théorie à la pratique. Le Baccalauréat technique évalue votre capacité à \textbf{mobiliser des connaissances précises} (chiffres, dates, acteurs, lieux) et à \textbf{structurer une démonstration} rigoureuse, qu'elle soit cartographique, statistique ou rédactionnelle. Cette séance de travaux pratiques vous guide pas à pas dans la réalisation de trois productions types de l'épreuve d'ECM : un croquis de synthèse, un diagramme statistique commenté et un paragraphe argumenté de niveau Bac.

\courssub{TP 1 : Construction d'un croquis « Le Cameroun et ses partenaires dans le monde »}

\begin{methode}[Construire un croquis géographique de synthèse (niveau Bac)]
Un croquis n'est pas un dessin artistique : c'est une \textbf{carte schématique codée} qui répond à une problématique précise. Au Bac, il est noté sur la rigueur de sa légende et la pertinence des informations localisées. Suivez impérativement ces 5 étapes :
\begin{enumerate}
    \item \textbf{Analyser le sujet} : identifier la thématique (ici : coopération, partenaires, aide humanitaire) et l'échelle (monde / Afrique / Cameroun).
    \item \textbf{Préparer la légende \emph{avant} de tracer} : la classer par \textbf{thèmes} (ex. : Partenaires bilatéraux, Organisations internationales, Zones humanitaires) et choisir des \textbf{signes cartographiques normalisés} (flèches proportionnelles, points, hachures, couleurs distinctes).
    \item \textbf{Tracer le fond de carte} : contours simplifiés (Cameroun, pays limitrophes, océans, méridiens/parallèles de repère), orientation (flèche Nord), échelle graphique.
    \item \textbf{Reporter les informations} : localiser avec précision (capitales, ports, régions), tracer les flèches (direction, épaisseur proportionnelle à l'intensité), poser les symboles.
    \item \textbf{Finaliser} : titre complet (Thème + Échelle + Date), légende encadrée et ordonnée, propreté (règle, crayons de couleur, pas de ratures).
\end{enumerate}
\end{methode}

\begin{attention}[Piège fréquent : la surcharge graphique]
Ne représentez \textbf{que} les informations utiles au sujet. Inutile de dessiner tous les fleuves, toutes les villes ou tous les pays du monde. Un croquis lisible comporte 8 à 12 éléments de légende maximum. Évitez les flèches qui se croisent de manière illisible ; regroupez les partenaires par zone géographique (Union européenne, Asie, Afrique) si nécessaire.
\end{attention}

\textbf{Consignes du TP 1.}\par
À partir d'un fond de carte Afrique-Monde fourni (ou tracé au tableau), réalisez le croquis « \emph{Le Cameroun et ses partenaires dans le monde} ». Votre légende doit comporter obligatoirement trois rubriques :
\begin{itemize}
    \item \textbf{Partenaires bilatéraux majeurs} (flèches d'épaisseur proportionnelle aux échanges/aide) : France, Chine, États-Unis, Nigeria, Allemagne.
    \item \textbf{Organisations internationales et régionales} (symboles ponctuels sur les capitales ou sièges) : ONU (New York, Yaoundé), UA (Addis-Abeba), CEMAC (Bangui), BAD (Abidjan), FMI/Banque mondiale (Washington).
    \item \textbf{Zones d'intervention humanitaire} (hachures ou couleur sur le territoire camerounais) : Extrême-Nord (crise Boko Haram, réfugiés nigérians), Est (réfugiés centrafricains), Nord-Ouest/Sud-Ouest (crise anglophone, déplacés internes).
\end{itemize}

\begin{popfigure}
\begin{tikzpicture}[scale=0.85, every node/.style={font=\small}]
    % Fond de carte simplifié : Afrique + Europe/Proche-Orient esquissés
    \draw[popline, line width=0.8pt, fill=popcream!30]
        (-5,-6) -- (12,-6) -- (12,8) -- (-2,10) -- (-6,4) -- cycle;
    \draw[popline, line width=0.5pt, fill=popblueL!40]
        (-5,8) -- (12,8) -- (12,14) -- (-5,14) -- cycle;

    % Cameroun (position approximative)
    \coordinate (CMR) at (2.5, 1.5);
    \fill[poporange!60] (CMR) circle (0.45cm);
    \node[below=0.1cm of CMR, font=\small\bfseries, text=popdark] {Cameroun};
    \node[below=0.6cm of CMR, font=\tiny, text=popmuted] {Yaoundé};

    % Villes partenaires (coordonnées relatives approximatives)
    \coordinate (PAR) at (-3, 11);
    \coordinate (BEI) at (10, 9);
    \coordinate (WAS) at (-8, 8);
    \coordinate (LAG) at (-1, 2);
    \coordinate (BER) at (-1, 10.5);
    \coordinate (NY) at (-9, 7.5);
    \coordinate (ADD) at (6, 3);
    \coordinate (BAN) at (1, -1);
    \coordinate (ABI) at (-3, 1);

    % Flèches de coopération bilatérale (épaisseur variable)
    \draw[popblue, line width=2.5pt, -{Stealth[length=3mm]}, opacity=0.7] (CMR) to[bend left=15] (PAR);
    \draw[poppink, line width=2.2pt, -{Stealth[length=3mm]}, opacity=0.7] (CMR) to[bend right=10] (BEI);
    \draw[popblue!60!black, line width=1.8pt, -{Stealth[length=3mm]}, opacity=0.7] (CMR) to[bend left=20] (WAS);
    \draw[popgreen!60!black, line width=1.5pt, -{Stealth[length=2.5mm]}, opacity=0.7] (CMR) -- (LAG);
    \draw[poppurple, line width=1.5pt, -{Stealth[length=2.5mm]}, opacity=0.7] (CMR) to[bend left=10] (BER);

    % Symboles Organisations Internationales (carrés)
    \foreach \pos/\label/\col in {NY/ONU/popblue, ADD/UA/popgreen, BAN/CEMAC/poporange, ABI/BAD/poppurple} {
        \node[draw=\col, fill=\col!20, rectangle, inner sep=1.5pt, font=\tiny\bfseries] at (\pos) {\label};
    }

    % Capitales partenaires (ronds petits)
    \foreach \pos/\name in {PAR/Paris, BEI/Pékin, WAS/Washington, LAG/Lagos, BER/Berlin} {
        \fill[popdark] (\pos) circle (1.5pt);
        \node[above=1pt of \pos, font=\tiny, anchor=south] {\name};
    }

    % Zones humanitaires (hachures sur le Cameroun - zones relatives)
    \pattern[pattern=north east lines, pattern color=poppink!60]
        ($(CMR)+(0.2,0.4)$) rectangle ++(0.6,0.5);
    \pattern[pattern=north west lines, pattern color=poporange!60]
        ($(CMR)+(0.7,0.0)$) rectangle ++(0.5,0.6);
    \pattern[pattern=crosshatch dots, pattern color=poppurple!60]
        ($(CMR)+(-0.5,0.3)$) rectangle ++(0.5,0.5);

    % Légende intégrée au TikZ (simplifiée pour l'exemple)
    \begin{scope}[shift={(-7.5,-4)}, font=\tiny]
        \node[anchor=west, font=\bfseries] at (0,2.5) {LÉGENDE SIMPLIFIÉE};
        \draw[popblue, line width=2.5pt, -{Stealth[length=2mm]}] (0,2) -- (1.5,2); \node[anchor=west] at (1.7,2) {France (coop. dense)};
        \draw[poppink, line width=2.2pt, -{Stealth[length=2mm]}] (0,1.6) -- (1.5,1.6); \node[anchor=west] at (1.7,1.6) {Chine (investissements)};
        \draw[popblue!60!black, line width=1.8pt, -{Stealth[length=2mm]}] (0,1.2) -- (1.5,1.2); \node[anchor=west] at (1.7,1.2) {USA (aide santé/sécurité)};
        \draw[popgreen!60!black, line width=1.5pt, -{Stealth[length=2mm]}] (0,0.8) -- (1.5,0.8); \node[anchor=west] at (1.7,0.8) {Nigeria (échanges régionaux)};
        \node[draw=popblue, fill=popblue!20, rectangle, inner sep=1pt] at (0.7,0.3) {ONU}; \node[anchor=west] at (1.7,0.3) {Sièges organisations};
        \fill[pattern=north east lines, pattern color=poppink!60] (0,-0.2) rectangle (0.4,0.2); \node[anchor=west] at (0.6,-0.1) {Extrême-Nord (Boko Haram)};
        \fill[pattern=north west lines, pattern color=poporange!60] (0,-0.6) rectangle (0.4,-0.2); \node[anchor=west] at (0.6,-0.4) {Est (réfugiés RCA)};
        \fill[pattern=crosshatch dots, pattern color=poppurple!60] (0,-1.0) rectangle (0.4,-0.6); \node[anchor=west] at (0.6,-0.8) {NOSO (déplacés internes)};
    \end{scope}

    % Orientation et Échelle
    \node[anchor=south east, font=\bfseries] at (11, -5.5) {N $\uparrow$};
    \draw[|-|, thick] (8, -5.8) -- (10, -5.8); \node[below, font=\tiny] at (9, -5.8) {$\approx$ 3000 km};
\end{tikzpicture}
\legende{Modèle de croquis attendu : « Le Cameroun et ses partenaires dans le monde ». Flèches bilatérales (épaisseur proportionnelle au volume de coopération), symboles des organisations (carrés), hachures pour les zones humanitaires. Titre, légende ordonnée, orientation et échelle sont obligatoires. Fond de carte volontairement schématique.}
\end{popfigure}

\begin{exoresolu}[Croquis commenté : éléments attendus dans la copie idéale]
\textbf{Énoncé :} Présentez les éléments qui doivent figurer dans la légende finale du croquis pour obtenir la note maximale.
\tcblower
\textbf{Solution détaillée :}
Le croquis remis au correcteur doit présenter une légende structurée en trois parties distinctes, encadrée, placée en bas à droite (ou à gauche) de la feuille.

\textbf{1. Partenaires bilatéraux (flèches orientées, épaisseur proportionnelle)}
\begin{itemize}
    \item \textbf{Flèche large bleue} vers la France (Paris) : premier partenaire historique, aide au développement (AFD), coopération militaire, culturelle et linguistique.
    \item \textbf{Flèche large rose} vers la Chine (Pékin) : premier partenaire commercial (import/export), infrastructures (barrages, routes, stades), prêts concessionnels.
    \item \textbf{Flèche moyenne bleu foncé} vers les États-Unis (Washington) : aide humanitaire (USAID), santé (PEPFAR), coopération militaire (formation), investissements privés.
    \item \textbf{Flèche moyenne verte} vers le Nigeria (Lagos/Abuja) : échanges régionaux (CEDEAO), pétrole, produits agricoles, sécurité transfrontalière (lutte contre Boko Haram).
    \item \textbf{Flèche moyenne violette} vers l'Allemagne (Berlin) : coopération technique (GIZ), énergies renouvelables, formation professionnelle.
\end{itemize}

\textbf{2. Organisations internationales et régionales (symboles ponctuels : carrés)}
\begin{itemize}
    \item \textbf{Carré bleu ONU} : siège à New York + agences à Yaoundé (PNUD, UNICEF, HCR, OMS, FAO). Rôle : aide humanitaire, ODD, maintien de la paix.
    \item \textbf{Carré vert UA} : siège à Addis-Abeba. Rôle : paix et sécurité, intégration africaine (ZLECAf), Agenda 2063.
    \item \textbf{Carré orange CEMAC} : siège à Bangui. Rôle : intégration monétaire (BEAC), marché commun, libre circulation.
    \item \textbf{Carré violet BAD} : siège à Abidjan. Rôle : financement des infrastructures (transport, énergie, eau), appui au secteur privé.
    \item \textbf{Carré gris FMI/BM} : siège à Washington. Rôle : facilités de crédit, appui budgétaire, réformes structurelles.
\end{itemize}

\textbf{3. Zones d'intervention humanitaire (hachures / surfaces colorées sur le territoire national)}
\begin{itemize}
    \item \textbf{Hachures roses -- Extrême-Nord} : crise Boko Haram (depuis 2014), plusieurs centaines de milliers de déplacés internes et environ 100\,000 réfugiés nigérians (camp de Minawao). Acteurs : HCR, PAM, ONG, MINAT.
    \item \textbf{Hachures oranges -- Région de l'Est} : réfugiés centrafricains (crises de 2013 et suivantes), environ 300\,000 personnes (sites de Gado-Badzéré, Lolo, Mbile). Acteurs : HCR, UNICEF, MSF, Croix-Rouge camerounaise.
    \item \textbf{Hachures violettes -- Nord-Ouest / Sud-Ouest} : crise anglophone (depuis 2016-2017), plusieurs centaines de milliers de déplacés internes. Acteurs : OCHA, ONG internationales, Comité International de la Croix-Rouge (CICR).
\end{itemize}

\textbf{Éléments cartographiques obligatoires (hors légende) :} titre centré \og Le Cameroun et ses partenaires dans le monde : coopération et aide humanitaire \fg, flèche du Nord, échelle graphique indicative, contours des pays limitrophes (Nigeria, Tchad, RCA, Congo, Gabon, Guinée équatoriale), repères (Équateur, méridien de Greenwich).
\end{exoresolu}

\begin{exercice}[TP 1 -- À vous de jouer]
Sur une feuille A4 (format paysage), tracez le fond de carte Afrique-Monde (simplifié) et réalisez le croquis complet avec sa légende ordonnée en 3 rubriques. Respectez la méthode en 5 étapes. Temps imparti : 25 minutes.
\end{exercice}

\begin{corrige}[Exercice TP 1]
\emph{Grille d'évaluation indicative (sur 6 pts) :}
\begin{itemize}
    \item Fond de carte propre, orientation, échelle, titre complet (1 pt).
    \item Légende structurée en 3 rubriques titrées, signes distincts et cohérents (1,5 pt).
    \item Localisation correcte de 5 partenaires bilatéraux + 4 organisations + 3 zones humanitaires (2 pts).
    \item Proportionnalité des flèches (qualitative) et choix des hachures adaptés (1 pt).
    \item Propreté générale, lisibilité, absence de surcharge (0,5 pt).
\end{itemize}
\end{corrige}

\courssub{TP 2 : Réalisation d'un diagramme en barres comparant les flux de coopération}

\begin{methode}[Construire et commenter un diagramme en barres groupées]
Un diagramme en barres compare des \textbf{valeurs quantitatives} pour des \textbf{catégories discrètes} (pays, années, secteurs). Pour le Bac :
\begin{enumerate}
    \item \textbf{Données} : choisir des séries homogènes (même unité, même année). Citer la source (ex. : MINFI, Banque mondiale, OCDE, Douanes).
    \item \textbf{Type} : barres \emph{groupées} (côte à côte) pour comparer deux variables (aide reçue vs échanges commerciaux) par partenaire ; barres \emph{empilées} pour montrer la composition d'un total.
    \item \textbf{Axes} : axe des ordonnées = valeur + unité (milliards de FCFA) ; axe des abscisses = partenaires. Échelle régulière, partant de 0.
    \item \textbf{Forme} : largeur des barres constante, espace entre les groupes, légende des séries (couleurs), titre informatif.
    \item \textbf{Commentaire type (3 phrases)} : 1. Lecture globale (qui domine ?). 2. Comparaison interne (écart aide/échanges par pays). 3. Interprétation (nature de la relation : commerciale, aidée, mixte).
\end{enumerate}
\end{methode}

\textbf{Données (ordres de grandeur indicatifs, à titre pédagogique) -- Unité : milliards de FCFA.}\par
\begin{center}
\begin{tabularx}{\linewidth}{|l|X|X|X|}\hline
\rowcolor{popblueL}
\textbf{Partenaire} & \textbf{Aide Publique au Développement (APD) reçue} & \textbf{Échanges commerciaux totaux (export + import)} & \textbf{Nature dominante} \\ \hline
France & 45 & 850 & Mixte (commerce ++ / aide ciblée) \\ \hline
Chine & 15 & 1\,200 & Commerciale (import massif / export bois, coton) \\ \hline
Union européenne (hors France) & 60 & 600 & Aide + commerce (APE) \\ \hline
États-Unis & 35 & 250 & Aide (santé/sécurité) + commerce (pétrole) \\ \hline
Nigeria & 2 & 400 & Commerciale régionale (pétrole, informel) \\ \hline
Allemagne & 20 & 120 & Aide technique (GIZ/KfW) + commerce \\ \hline
\end{tabularx}
\end{center}

\begin{attention}[Valeurs indicatives]
Ces chiffres sont des \textbf{ordres de grandeur pédagogiques} construits pour l'entraînement : ils ne remplacent pas les statistiques officielles (INS, MINFI, Douanes). Au Bac, si des données vous sont fournies dans le sujet, utilisez \emph{exclusivement} celles du document.
\end{attention}

\begin{popfigure}
\begin{tikzpicture}
\begin{axis}[
    width=\linewidth,
    height=0.55\linewidth,
    ybar,
    bar width=0.3cm,
    ymin=0, ymax=1300,
    ymajorgrids=true,
    grid style={popline, dashed},
    ylabel={Valeur (milliards de FCFA)},
    xlabel={Principaux partenaires (ordres de grandeur)},
    symbolic x coords={France, Chine, UE, USA, Nigeria, Allemagne},
    xtick=data,
    xticklabel style={rotate=30, anchor=east, font=\footnotesize},
    yticklabel style={font=\footnotesize},
    legend style={at={(0.5,-0.32)}, anchor=north, legend columns=2, font=\footnotesize, draw=popline, fill=popcream!50},
    enlarge x limits=0.12,
]
\addplot[popblue, fill=popblue!60] coordinates {
    (France,850) (Chine,1200) (UE,600) (USA,250) (Nigeria,400) (Allemagne,120)
};
\addplot[poporange, fill=poporange!60] coordinates {
    (France,45) (Chine,15) (UE,60) (USA,35) (Nigeria,2) (Allemagne,20)
};
\legend{Échanges commerciaux (export + import), APD reçue (dons + prêts concessionnels)}
\end{axis}
\end{tikzpicture}
\legende{Diagramme en barres groupées : comparaison APD reçue / échanges commerciaux avec les 6 principaux partenaires du Cameroun (ordres de grandeur pédagogiques). L'écart d'échelle (échanges très supérieurs à l'APD) illustre le poids croissant des partenariats commerciaux.}
\end{popfigure}

\begin{exoresolu}[Commentaire type du diagramme (niveau Bac -- 8 à 10 lignes)]
\textbf{Énoncé :} Rédigez un commentaire structuré du diagramme ci-dessus.
\tcblower
\textbf{Solution détaillée :}
Le diagramme met en évidence une \textbf{asymétrie marquée} entre le volume des échanges commerciaux et le montant de l'Aide Publique au Développement (APD) pour l'ensemble des partenaires du Cameroun. En valeur absolue, la \textbf{Chine s'impose comme le premier partenaire commercial} (1\,200 milliards de FCFA), devançant la France (850 milliards) et le Nigeria (400 milliards), alors que son APD reste modeste (15 milliards), traduisant une relation \textbf{essentiellement marchande} (importation de biens manufacturés, exportation de matières premières). À l'inverse, l'\textbf{Union européenne (hors France) et la France} sont les premiers pourvoyeurs d'APD (60 et 45 milliards), confirmant une coopération \textbf{structurelle et sectorielle} (gouvernance, monde rural, éducation) couplée à des échanges significatifs. Les \textbf{États-Unis} présentent un profil d'\textbf{aide ciblée} (35 milliards d'APD contre 250 milliards d'échanges), l'aide étant concentrée sur la santé (VIH/sida, paludisme) et la sécurité. Le \textbf{Nigeria} illustre l'intégration régionale : échanges soutenus (pétrole, commerce frontalier) pour une APD quasi nulle. Globalement, le Cameroun opère une \textbf{diversification de ses partenariats} : la part des pays émergents (Chine, Inde, Turquie) dans les échanges croît, tandis que les partenaires traditionnels (UE, France) conservent le leadership sur le financement concessionnel du développement.
\end{exoresolu}

\begin{exercice}[TP 2 -- À vous de jouer]
À partir du tableau de données, tracez sur papier millimétré le diagramme en barres groupées (échelle : 1 cm pour 100 milliards de FCFA). Rédigez ensuite un commentaire de 8 lignes maximum en appliquant la méthode en 3 phrases. Temps imparti : 20 minutes.
\end{exercice}

\begin{corrige}[Exercice TP 2]
\emph{Points de vigilance pour la correction :}
\begin{itemize}
    \item Échelle des ordonnées partant de 0 (sinon distorsion visuelle).
    \item Barres de même largeur, espacement régulier, légende claire (Échanges / APD).
    \item Commentaire : phrase 1 = lecture globale (échanges très supérieurs à l'APD partout) ; phrase 2 = contraste Chine (commerce) vs UE/France (aide + commerce) ; phrase 3 = interprétation (diversification, nature des relations).
    \item Citation d'au moins 3 chiffres précis (ex. : 1\,200 contre 15 pour la Chine ; 60 d'APD pour l'UE).
\end{itemize}
\end{corrige}

\courssub{TP 3 : Rédaction guidée d'un paragraphe argumenté -- Le règlement du conflit de Bakassi}

\begin{methode}[Rédiger un paragraphe argumenté de niveau Baccalauréat (ECM)]
L'épreuve écrite d'ECM exige souvent une \textbf{réponse construite} (10 à 15 lignes). La structure « thèse -- arguments -- exemples -- ouverture » est la norme.
\begin{enumerate}
    \item \textbf{Thèse (1 phrase) :} réponse directe, affirmative et nuancée à la question, en reprenant les termes du sujet.
    \item \textbf{Arguments juridiques et politiques (2 à 3 phrases) :} mobiliser les \textbf{principes de la diplomatie camerounaise} vus en cours (droit international, règlement pacifique, non-usage de la force, dialogue).
    \item \textbf{Preuves factuelles précises (2 à 3 phrases) :} citer \textbf{dates, actes juridiques, acteurs, résultats}. C'est ce qui distingue la copie « Bac » de la copie moyenne.
    \item \textbf{Ouverture (1 phrase) :} portée générale, leçon pour l'avenir, lien avec un autre dossier (ex. : frontière du lac Tchad, ZLECAf).
\end{enumerate}
\textbf{Connecteurs logiques obligatoires :} \emph{en effet, par conséquent, qui plus est, illustrons par, cela démontre que, dès lors, par ailleurs.}
\end{methode}

\begin{exemplebox}[Paragraphe modèle : « Le règlement du conflit de Bakassi illustre les principes de la diplomatie camerounaise »]
\textbf{Thèse :} Le règlement du différend de la péninsule de Bakassi constitue l'\textbf{illustration exemplaire} de l'application rigoureuse des principes de la diplomatie camerounaise, alliant respect du droit international, refus de la force et recherche du consensus.

\textbf{Argument 1 -- Légalisme et recours à la justice internationale :} face au contentieux frontalier avec le Nigeria, le Cameroun a saisi la \textbf{Cour Internationale de Justice (CIJ)} dès 1994, refusant l'escalade militaire. L'\textbf{arrêt du 10 octobre 2002} a reconnu la souveraineté camerounaise sur Bakassi, notamment sur la base des accords coloniaux germano-britanniques.

\textbf{Argument 2 -- Mise en œuvre pacifique et diplomatie préventive :} plutôt que d'exiger un retrait immédiat par la force, Yaoundé a accepté la \textbf{médiation du Secrétaire général de l'ONU (Kofi Annan)}, aboutissant aux \textbf{Accords de Greentree (12 juin 2006)}. Ce texte prévoyait un transfert progressif (retrait de l'administration nigériane, puis transfert définitif en 2008), avec des garanties pour les populations nigérianes de la péninsule.

\textbf{Argument 3 -- Respect des populations et stabilité régionale :} la \textbf{Commission Mixte Cameroun-Nigeria}, soutenue par l'ONU, a supervisé le retrait effectif des troupes nigérianes en \textbf{août 2008} et le processus de démarcation de la frontière terrestre et maritime. Le transfert de souveraineté s'est déroulé sans affrontement majeur, illustrant les principes de \textbf{règlement pacifique des différends} et de \textbf{bon voisinage}.

\textbf{Ouverture :} ce « modèle Bakassi » inspire aujourd'hui la gestion d'autres questions frontalières (lac Tchad) et renforce la crédibilité du Cameroun comme acteur de la \textbf{paix et de la sécurité en Afrique centrale}.
\end{exemplebox}

\begin{attention}[Erreurs rédactionnelles sanctionnées au Bac]
\begin{itemize}
    \item \textbf{Confondre la CIJ} (organe judiciaire de l'ONU, siège à La Haye, règle les différends entre États) \textbf{et la CPI} (Cour pénale internationale, juge des individus). Bakassi = CIJ.
    \item \textbf{Omettre les dates clés :} 1994 (saisine), 2002 (arrêt de la CIJ), 2006 (Accords de Greentree), 2008 (transfert définitif).
    \item \textbf{Écrire « le Cameroun a gagné la guerre » :} c'est faux. C'est une victoire \emph{juridique et diplomatique}, obtenue sans guerre.
    \item \textbf{Paragraphe trop court (moins de 6 lignes) ou trop long (plus de 15 lignes) sans structure visible.}
    \item \textbf{Absence de connecteurs logiques} (suite de phrases juxtaposées).
\end{itemize}
\end{attention}

\begin{savaistu}[Le Bac ECM et les exemples camerounais]
Les correcteurs du Baccalauréat technique valorisent systématiquement les copies qui \textbf{ancrent l'argumentation dans la réalité camerounaise récente et vérifiable}. Un élève qui cite « l'Accord de Greentree de 2006 » ou « le transfert définitif d'août 2008 » obtient une meilleure note qu'un élève qui écrit seulement « le Cameroun a récupéré Bakassi par la diplomatie ». De même, pour l'aide humanitaire, citer « le camp de Minawao (Extrême-Nord) » ou « la réponse humanitaire à la crise du Nord-Ouest et du Sud-Ouest » fait la différence. \textbf{La précision factuelle est la marque de fabrique du candidat reçu.}
\end{savaistu}

\begin{exercice}[TP 3 -- Sujet d'entraînement Bac]
\textbf{Sujet :} « Le règlement du conflit de Bakassi illustre les principes de la diplomatie camerounaise. » En vous appuyant sur vos connaissances, rédigez un paragraphe argumenté d'une dizaine de lignes pour justifier cette affirmation.
\end{exercice}

\begin{corrige}[Exercice TP 3 -- Grille d'attente détaillée]
\emph{Éléments obligatoires pour la note maximale (sur 6 pts) :}
\begin{center}
\begin{tabularx}{\linewidth}{|l|X|}\hline
\rowcolor{popblueL}
\textbf{Critère} & \textbf{Indicateurs de réussite} \\ \hline
Thèse (1 pt) & Phrase d'ouverture claire : « Le règlement de Bakassi est l'illustration exemplaire des principes de la diplomatie camerounaise... » \\ \hline
Argument 1 : Droit (1 pt) & Citation de la CIJ + dates 1994/2002 + base juridique (accords coloniaux). Vocabulaire : « souveraineté », « arrêt contraignant ». \\ \hline
Argument 2 : Paix/Négociation (1 pt) & Accords de Greentree (12 juin 2006), médiation ONU (Kofi Annan), transfert progressif, protection des civils. Vocabulaire : « règlement pacifique », « diplomatie préventive ». \\ \hline
Argument 3 : Voisinage/Stabilité (1 pt) & Commission Mixte Cameroun-Nigeria, démarcation de la frontière, retrait d'août 2008, bon voisinage. \\ \hline
Ouverture (1 pt) & Lien avec le lac Tchad, la ZLECAf, la crédibilité internationale du Cameroun. \\ \hline
Forme/Langue (1 pt) & 8 à 12 lignes, connecteurs logiques, orthographe soignée, termes techniques (CIJ, APD, souveraineté). \\ \hline
\end{tabularx}
\end{center}
\end{corrige}

\courssub{Correction collective et grille d'auto-évaluation des productions}

La séance se termine par une \textbf{correction collective} (15 minutes) au cours de laquelle 3 volontaires présentent leur croquis, leur diagramme et leur paragraphe au tableau ou au vidéoprojecteur. La classe évalue à l'aide de la grille ci-dessous. L'enseignant valide, complète et note les points forts et les axes de progression.

\begin{aretenir}[Grille d'auto-évaluation TP Diplomatie -- À coller dans le cahier]
\begin{center}
\begin{tabularx}{\linewidth}{|l|X|X|X|}\hline
\rowcolor{popblueL}
\textbf{Compétence évaluée} & \textbf{Non acquis} & \textbf{En cours d'acquisition} & \textbf{Acquis / Expert} \\ \hline
\textbf{Croquis (TP1)} & Légende absente ou incomplète ; pas d'échelle ni d'orientation ; flèches aléatoires ; surcharge illisible. & Légende en 3 rubriques ; 4-5 partenaires localisés ; flèches qualitatives ; titre, orientation et échelle présents. & Légende complète et structurée (3 rubriques, signes normalisés) ; 5 partenaires + 4 organisations + 3 zones humanitaires ; proportionnalité des flèches ; propreté exemplaire. \\ \hline
\textbf{Diagramme (TP2)} & Axe non gradué ou ne partant pas de 0 ; barres de largeurs variables ; pas de légende des séries ; commentaire absent ou purement descriptif. & Axes corrects ; barres groupées régulières ; légende des séries ; commentaire en 3 phrases (global + contraste + interprétation) avec 2 chiffres. & Graphique parfait (échelle, unités, source) ; commentaire structuré (thèse -- chiffres -- interprétation) ; 4 chiffres précis cités ; analyse de la nature des partenariats. \\ \hline
\textbf{Paragraphe (TP3)} & Moins de 6 lignes ; pas de thèse ; confusion CIJ/CPI ; pas de dates ; pas de connecteurs ; ouverture absente. & 8-10 lignes ; thèse claire ; 2 arguments + 1 exemple daté (2002 ou 2006) ; connecteurs présents ; ouverture simple. & 10-12 lignes ; thèse nuancée ; 3 arguments juridiques et politiques ; 4 dates/actes précis (CIJ 2002, Greentree 2006, retrait 2008, Commission Mixte) ; connecteurs variés ; ouverture géopolitique. \\ \hline
\end{tabularx}
\end{center}
\end{aretenir}

\begin{experience}[Démarche d'investigation : « Enquête de terrain » (prolongement hors classe)]
Pour ancrer ces apprentissages dans le réel, chaque binôme d'élèves réalise une \textbf{mini-enquête} (1 semaine) auprès d'un acteur local de la coopération :
\begin{enumerate}
    \item \textbf{Cible :} responsable d'une ONG locale (ex. : Croix-Rouge camerounaise), agent d'une délégation régionale (MINEPAT, MINRESI), entrepreneur import-export, chef de projet de développement communal.
    \item \textbf{Questionnaire type (5 questions) :} 1) Quel partenaire ou financeur ? 2) Quel type de flux (don, prêt, contrat commercial) ? 3) Quel secteur (santé, infrastructures, formation) ? 4) Quelles difficultés rencontrées (procédures, sécurité, change) ? 5) Quelle perception de l'action internationale du Cameroun (efficacité, visibilité) ?
    \item \textbf{Restitution :} affichage en classe « Carte des acteurs de la coopération dans ma région » + synthèse orale de 3 minutes.
\end{enumerate}
Cette démarche développe l'esprit critique et fournit des exemples \emph{vivants} et locaux, réutilisables dans les copies du Probatoire et du Baccalauréat technique.
\end{experience}

\newpage

\coursec{Activité d'intégration}

\coursec{Les principes de la diplomatie camerounaise}

\courssub{Objectifs de l'activité}
Cette activité d'intégration vise à mobiliser l'ensemble des savoirs et savoir-faire de la leçon dans une situation professionnelle simulée. Elle permet à l'élève de :
\begin{itemize}
    \item Identifier et qualifier les différentes \cle{formes} et \cle{aspects} de la coopération internationale à partir de documents authentiques.
    \item Appliquer les \cle{principes de la diplomatie camerounaise} pour justifier une décision d'État concrète.
    \item Produire une représentation graphique (croquis) synthétisant l'action extérieure du Cameroun.
    \item Rédiger un discours argumenté, structuré et juridiquement fondé, respectant les codes de la communication diplomatique.
\end{itemize}

\courssub{Situation}
\emph{Contexte :} Nous sommes en janvier 2025. Le Cameroun fait face à une conjoncture sous-régionale complexe. La République centrafricaine (RCA), pays frontalier lié au Cameroun par l'histoire, la géographie et les accords de la \cle{CEMAC}, traverse une nouvelle flambée de violence dans l'Ouest du pays (préfectures de l'Ouham-Pendé et de la Nana-Mambéré). Le Secrétaire général des Nations Unies, par la voix de la MINUSCA (Mission multidimensionnelle intégrée des Nations Unies pour la stabilisation en Centrafrique), sollicite officiellement le Cameroun pour fournir un contingent de 750 casques bleus (bataillon d'infanterie + appui logistique) et ouvrir un corridor humanitaire terrestre via Garoua-Boulaï pour acheminer 12 000 tonnes de vivres du PAM (Programme alimentaire mondial) vers les populations déplacées de l'Ouest centrafricain.

Le Ministre des Relations Extérieures (MINREX), S.E. Mbella Mbella, convoque une cellule de crise élargie. Vous êtes les \textbf{conseillers techniques} chargés de préparer le dossier de décision. Sur la table, cinq documents officiels constituent le \emph{dossier documentaire} :

\begin{enumerate}
    \item \textbf{Document 1 — Carte HCR « Situation des réfugiés au Cameroun, nov. 2024 » :} Le Cameroun accueille 473 827 réfugiés et demandeurs d'asile (HCR). Dont 342 117 en provenance de la RCA (72,2 \%), répartis principalement dans l'Est (sites de Gado-Badzéré, Lolo, Mbilé, Timangolo) et l'Adamaoua. 124 500 réfugiés nigérians dans l'Extrême-Nord (site de Minawao : 68 400 personnes). Le reste : Tchadiens, Soudanais, Congolais.
    \item \textbf{Document 2 — Fiche technique HCR Minawao (oct. 2024) :} Site ouvert en 2013. Superficie : 623 ha. Population : 68 400 réfugiés nigérians (52 \% femmes, 58 \% < 18 ans). Taux de scolarisation primaire : 84 \%. Accès à l'eau : 18 L/pers/jour (standard SPHERE : 15 L). Partenaires : HCR, UNICEF, WFP, ONG locales (ALDEPA, IEDA-Relief). Budget annuel : 14,2 M USD (gap de financement 38 \%).
    \item \textbf{Document 3 — Extrait de l'arrêt CIJ, 10 oct. 2002 (Affaire Bakassi) :} « La Cour, par 9 voix contre 3, décide que la souveraineté sur la presqu'île de Bakassi appartient au Cameroun... Le Nigeria doit retirer son administration et ses forces militaires... » Suivi de l'Accord de Greentree (2006) et du retrait effectif nigérian (2008). Le Cameroun a privilégié la \cle{voie juridique} et le \cle{dialogue} à la force.
    \item \textbf{Document 4 — Statistiques Douanes/INS « Échanges Cameroun–Chine 2023 » :} Volume total : 1 842 milliards FCFA. Exportations Cameroun vers Chine : 412 Mrd FCFA (bois, coton, aluminium, pétrole brut). Importations Chine vers Cameroun : 1 430 Mrd FCFA (machines, textiles, électronique, engins BTP). Investissements chinois (FDI stock 2023) : ~ 3 200 Mrd FCFA (barrage de Nachtigal, autoroute Kribi-Lolabé, stade Olembé, centres hospitaliers). Accord de partenariat stratégique global (2018).
    \item \textbf{Document 5 — Photo officielle : 16\textsuperscript{e} Sommet ordinaire des Chefs d'État de la CEMAC (Yaoundé, 17 mars 2024).} Communiqué final : « Les Chefs d'État réaffirment leur engagement pour la mutualisation des forces de sécurité (FORCE CEMAC), la libre circulation des personnes et des biens, et la mise en œuvre effective de l'Union douanière. » Présence du Président Paul Biya (Président en exercice de la CEMAC 2023-2024).
\end{enumerate}

\courssub{Consignes}
Travaillant en groupe de 4 à 5 élèves (cellule de conseil), vous devez produire le \textbf{dossier de plaidoyer} suivant, à remettre au Ministre avant 16h :

\begin{enumerate}
    \item \textbf{Analyse documentaire (tableau) :} Pour chacun des 5 documents, complétez le tableau ci-dessous en identifiant :
    \begin{itemize}
        \item L'\cle{aspect} de la coopération concerné (politique, économique, culturel, scientifique/technique, militaire/sécuritaire, humanitaire).
        \item La \cle{forme} de la coopération (bilatérale, multilatérale, régionale/sous-régionale, décentralisée, ONG/associative).
        \item Le \cle{principe de la diplomatie camerounaise} illustré (non-ingérence, règlement pacifique, solidarité africaine, non-alignement, coopération gagnant-gagnant, respect du droit international).
    \end{itemize}
    
    \item \textbf{Croquis commenté :} Sur une carte de base du Cameroun et de son environnement immédiat (fournie), réalisez un croquis légendé intitulé : \emph{« Le Cameroun, acteur de la paix et de la coopération »}. Faites apparaître :
    \begin{itemize}
        \item Les flux de réfugiés (flèches proportionnelles vers l'Est et l'Extrême-Nord).
        \item Les axes de coopération économique majeure (ex. corridor Douala-N'Djaména, projet ferroviaire, investissements chinois).
        \item Les zones d'engagement militaire/régional (Bakassi, zone CEMAC, zone LCBC/Commission du bassin du Lac Tchad).
        \item Le corridor humanitaire proposé vers la RCA (Garoua-Boulaï → Bouar/Bossangoa).
    \end{itemize}
    
    \item \textbf{Rédaction du discours (15 lignes maximum, police Times New Roman 12, interligne 1,15) :} En votre qualité de Porte-parole adjoint du MINREX, rédigez la déclaration officielle justifiant l'engagement du Cameroun dans la mission humanitaire et de maintien de la paix en Centrafrique. Le discours doit :
    \begin{itemize}
        \item S'appuyer explicitement sur \textbf{au moins 4 principes} de la diplomatie camerounaise.
        \item Intégrer \textbf{2 données chiffrées} issues du dossier documentaire.
        \item Adopter le ton formel de la diplomatie (allocution, engagement solennel, appel à la communauté internationale).
    \end{itemize}
\end{enumerate}

\courssub{Ressources à mobiliser}
\begin{aretenir}[Rappel synthétique : Principes directeurs de la diplomatie camerounaise (Constitution 1996, art. 45 ; Discours d'investiture 2018)]
\begin{itemize}
    \item \cle{Respect du droit international} et des traités ratifiés (Charte ONU, Acte constitutif UA, Traité CEMAC).
    \item \cle{Règlement pacifique des différends} (négociation, médiation, juridictions internationales — ex. CIJ Bakassi).
    \item \cle{Non-ingérence} dans les affaires intérieures des États souverains.
    \item \cle{Solidarité africaine} et panafricanisme (intégration régionale, soutien aux pays frères).
    \item \cle{Non-alignement} et diversification des partenariats (Sud-Sud, Nord-Sud, gagnant-gagnant).
    \item \cle{Coopération mutuellement avantageuse} (économique, technique, culturelle).
    \item \cle{Promotion des droits de l'homme} et du droit humanitaire (accueil réfugiés, protection civils).
\end{itemize}
\end{aretenir}

\begin{aretenir}[Aspects et Formes de la coopération — Grille de lecture]
\begin{center}
\begin{tabularx}{\linewidth}{|l|X|X|}\hline
\textbf{Aspect} & \textbf{Description} & \textbf{Exemples camerounais} \\ \hline
Politique & Relations État-État, diplomatie, traités & Accord de Greentree, Sommets UA/CEMAC \\ \hline
Économique & Commerce, investissements, financement & Partenariat Cameroun-Chine, ZLECAf \\ \hline
Culturel/Éducatif & Échanges universitaires, patrimoine, langues & Bourses Chine/France/Allemagne, IFE \\ \hline
Scientifique/Technique & Recherche, transfert tech, santé & IRAD, centres hospitaliers chinois \\ \hline
Militaire/Sécuritaire & Défense, formation, maintien paix & FORCE CEMAC, MINUSCA, BIR, LCBC \\ \hline
Humanitaire & Urgence, réfugiés, catastrophe & Accueil RCA/Nigeria, corridor PAM \\ \hline
\end{tabularx}
\end{center}

\begin{center}
\begin{tabularx}{\linewidth}{|l|X|X|}\hline
\textbf{Forme} & \textbf{Définition} & \textbf{Exemples camerounais} \\ \hline
Bilatérale & 2 États & Cameroun–France, Cameroun–Chine, Cameroun–Nigeria \\ \hline
Multilatérale & Organisations universelles & ONU, OMS, UNESCO, OMC, FMI, Banque Mondiale \\ \hline
Régionale/Sous-régionale & Organisations géographiques & CEMAC, CEEAC, UA, LCBC, CICOS \\ \hline
Décentralisée & Collectivités territoriales & Jumelages villes (Yaoundé–Paris, Douala–Marseille) \\ \hline
ONG/Associative & Acteurs non étatiques & Croix-Rouge, MSF, Plan International, ALDEPA \\ \hline
\end{tabularx}
\end{center}
\end{aretenir}

\courssub{Démarche et corrigé commenté}
\begin{exoresolu}[Dossier de plaidoyer — Corrigé type attendu]
\textbf{1. Tableau d'analyse documentaire}

\begin{center}
\begin{tabularx}{\linewidth}{|c|X|X|X|}\hline
\textbf{Document} & \textbf{Aspect(s) de coopération} & \textbf{Forme(s) de coopération} & \textbf{Principe(s) diplomatiques illustrés} \\ \hline
Doc 1 (Carte HCR réfugiés) & Humanitaire ; Social & Multilatérale (HCR/ONU) ; Bilatérale (accueil RCA, Nigeria) & Solidarité africaine ; Respect droit humanitaire ; Non-ingérence (asile = acte souverain) \\ \hline
Doc 2 (Fiche Minawao) & Humanitaire ; Éducatif ; Sanitaire & Multilatérale (HCR, UNICEF, PAM) ; ONG/Associative (ALDEPA, IEDA) & Coopération gagnant-gagnant (intégration locale) ; Droits de l'homme \\ \hline
Doc 3 (Arrêt CIJ Bakassi) & Politique ; Juridique ; Territorial & Multilatérale (CIJ/ONU) ; Bilatérale (Cameroun–Nigeria via Greentree) & Règlement pacifique des différends ; Respect droit international ; Non-violence \\ \hline
Doc 4 (Échanges Chine 2023) & Économique ; Scientifique/Technique ; Infrastructures & Bilatérale (Partenariat stratégique global) ; Multilatérale (FOCAC) & Non-alignement (diversification) ; Coopération mutuellement avantageuse \\ \hline
Doc 5 (Sommet CEMAC 2024) & Politique ; Militaire/Sécuritaire ; Économique (Union douanière) & Régionale/Sous-régionale (CEMAC) & Solidarité africaine ; Intégration régionale ; Sécurité collective \\ \hline
\end{tabularx}
\end{center}

\textbf{2. Croquis commenté — « Le Cameroun, acteur de la paix et de la coopération »}

\begin{methode}[Réalisation du croquis (étapes pour l'élève)]
\begin{enumerate}
    \item \textbf{Fond de carte :} Tracer le contour du Cameroun + pays limitrophes (Nigeria, Tchad, RCA, Congo, Guinée-Équatoriale, Gabon). Mer au sud-ouest. Capitales : Yaoundé, N'Djaména, Bangui, Abuja, Libreville, Brazzaville, Malabo.
    \item \textbf{Flux réfugiés (flèches rouges pleines, épaisseur ∝ effectifs) :}
    \begin{itemize}
        \item RCA → Est Cameroun (Gado, Lolo, Mbilé, Timangolo) : flèche large (342 117).
        \item Nigeria → Extrême-Nord (Minawao) : flèche moyenne (124 500).
    \end{itemize}
    \item \textbf{Coopération économique (flèches bleues double trait) :}
    \begin{itemize}
        \item Douala → N'Djaména (corridor routier/ferroviaire CEMAC).
        \item Kribi (port en eau profonde) → Chine (navires) : investissements (barrage Nachtigal, autoroute, stade).
        \item Douala ↔ Monde (export bois, coton, aluminium).
    \end{itemize}
    \item \textbf{Engagements sécuritaires/régionaux (symboles casque bleu / bouclier) :}
    \begin{itemize}
        \item Bakassi (sud-ouest) : « Règlement pacifique CIJ 2002 ».
        \item Zone LCBC (Extrême-Nord, Lac Tchad) : « Lutte anti-Boko Haram, gestion eau ».
        \item Zone CEMAC (arc Douala–Libreville–Bangui–N'Djaména) : « FORCE CEMAC, Union douanière ».
    \end{itemize}
    \item \textbf{Nouveau corridor humanitaire RCA (flèche verte pointillée + picto camion/blé) :}
    \item Garoua-Boulaï (frontière) → Bouar → Bossangoa (Ouest RCA). Légende : « 12 000 t vivres PAM / 750 casques bleus ».
    \item \textbf{Légende organisée} (3 colonnes : Peuplement/Réfugiés ; Économie/Infrastructures ; Sécurité/Paix). Titre, échelle, orientation (Nord).
\end{enumerate}
\end{methode}

\textbf{3. Discours du Porte-parole adjoint du MINREX (Proposition de corrigé — 15 lignes)}

\begin{quote}
\textbf{Déclaration de la République du Cameroun sur l'engagement en Centrafrique}

« Monsieur le Secrétaire général, Excellences, Mesdames, Messieurs,

Au nom de Son Excellence Paul Biya, Président de la République, Chef de l'État, le Cameroun réaffirme son attachement indéfectible aux principes fondateurs de notre diplomatie : le \textbf{règlement pacifique des différends}, la \textbf{solidarité africaine} et le \textbf{respect du droit international humanitaire}.

Fort de son expérience d'État hôte de \textbf{473 827 réfugiés} (dont 342 117 Centrafricains), le Cameroun ne saurait rester indifférent à la détresse de ses frères centrafricains. C'est pourquoi, conformément à l'esprit de l'Acte constitutif de l'Union africaine et du Traité de la CEMAC, nous acceptons de déployer un bataillon de \textbf{750 casques bleus} au sein de la MINUSCA et d'ouvrir le corridor humanitaire de Garoua-Boulaï pour l'acheminement de 12 000 tonnes de vivres.

Cet engagement illustre notre doctrine de \textbf{non-ingérence} — nous intervenons sur demande légitime — et de \textbf{coopération gagnant-gagnant} : la stabilité de la RCA conditionne la sécurité de notre propre frontière est.

Le Cameroun appelle la communauté internationale à combler le gap de financement humanitaire, gage d'une paix durable en Centrafrique comme dans le bassin du Lac Tchad.

Je vous remercie. »
\end{quote}

\textbf{Analyse du corrigé (grille d'évaluation) :}
\begin{itemize}
    \item \textbf{Principes explicités (4+) :} Règlement pacifique, Solidarité africaine, Respect droit international humanitaire, Non-ingérence, Coopération gagnant-gagnant. (✓)
    \item \textbf{Données chiffrées intégrées (2+) :} 473 827 réfugiés (342 117 Centrafricains) ; 750 casques bleus ; 12 000 tonnes vivres. (✓)
    \item \textbf{Ton diplomatique :} Formules protocolaires (« Au nom de... », « Excellences », « Je vous remercie »), engagement solennel, référence aux textes fondateurs. (✓)
    \item \textbf{Concision :} 15 lignes exactement (format Times 12, interligne 1,15). (✓)
\end{itemize}
\end{exoresolu}

\courssub{Bilan de l'activité}
Cette activité d'intégration a permis de vérifier l'appropriation opérationnelle des concepts de la leçon. Les élèves ont démontré leur capacité à :
\begin{itemize}
    \item \textbf{Discriminer} les \cle{aspects} (humanitaire, économique, politique, sécuritaire) et les \cle{formes} (bilatérale, multilatérale, régionale, ONG) de coopération à travers des documents hétérogènes (carte, statistiques, texte juridique, photo).
    \item \textbf{Articuler} les \cle{principes de la diplomatie camerounaise} non comme une liste abstraite, mais comme des \emph{arguments juridiques et moraux} mobilisables dans une décision d'État réelle.
    \item \textbf{Synthétiser} spatialement l'action extérieure du Cameroun via un croquis géopolitique rigoureux (flux, corridors, zones d'engagement, symboles adaptés).
    \item \textbf{Rédiger} dans le registre formel de la diplomatie, en respectant les contraintes de format (15 lignes), de ton (solennel, engageant) et de preuve (données chiffrées, références textuelles).
\end{itemize}

\begin{attention}[Pièges fréquents relevés lors de la correction collective]
\begin{itemize}
    \item Confondre \emph{aspect} (secteur thématique) et \emph{forme} (cadre institutionnel). Ex. : « L'aide du HCR est une forme multilatérale » → \textbf{Non}, c'est une forme multilatérale \emph{portant sur l'aspect humanitaire}.
    \item Oublier le principe de \cle{non-ingérence} dans le discours : l'intervention en RCA se fait \emph{sur demande} (ONU/État hôte), pas unilatéralement.
    \item Croquis illisible : pas de légende hiérarchisée, flèches non proportionnelles, absence d'échelle/orientation.
    \item Discours trop « journalistique » ou « militant » : absence de formules protocolaires, pas de référence aux textes (Charte ONU, Traité CEMAC, Constitution).
    \item Données chiffrées inventées ou approximatives : obligation de puiser dans le dossier documentaire fourni.
\end{itemize}
\end{attention}

\begin{aretenir}[Compétences certificatives visées (Probatoire / Baccalauréat Technique)]
\begin{itemize}
    \item \textbf{Analyser} un dossier documentaire hétérogène pour en extraire des informations pertinentes sur les relations internationales.
    \item \textbf{Caractériser} la coopération internationale (aspects, formes, acteurs) dans le contexte camerounais.
    \item \textbf{Justifier} une position diplomatique en mobilisant les principes constitutionnels et les engagements internationaux du Cameroun.
    \item \textbf{Communiquer} par écrit (discours officiel) et par l'image (croquis géopolitique) dans un registre professionnel.
\end{itemize}
\end{aretenir}

\newpage

\coursec{Méthodes et exercices résolus}

\courssub{Méthodes et exercices résolus}

\begin{methode}[Analyser une situation de relations internationales]
\textbf{Objectif :} Appliquer le concept de relations internationales (RI) à un fait d'actualité ou un cas concret.
\begin{enumerate}
    \item \textbf{Identifier les acteurs :} États, organisations internationales (OI), ONG, entreprises multinationales, groupes de pression.
    \item \textbf{Qualifier la nature de l'interaction :} Coopération (pacifique), conflit (tension, guerre), concurrence, dépendance, interdépendance.
    \item \textbf{Repérer les enjeux :} Sécurité, économie, environnement, droits de l'homme, influence géopolitique, identité.
    \item \textbf{Situer dans l'espace et le temps :} Bilatéral / Multilatéral ; Ponctuel / Structurel ; Régional / Mondial.
    \item \textbf{Conclure :} Résumer la dynamique (vers l'apaisement, l'escalade, l'institutionnalisation).
\end{enumerate}
\cle{Acteurs} -- \cle{Interactions} -- \cle{Enjeux} -- \cle{Niveau d'analyse}
\end{methode}

\begin{exoresolu}[Analyse d'une situation : le différend frontalier de Bakassi]
\textbf{Énoncé :} En vous référant à l'affaire de la péninsule de Bakassi entre le Cameroun et le Nigeria (règlement par la CIJ en 2002, retrait nigérian en 2006 sous l'égide de l'ONU), analysez cette situation en mobilisant les concepts de relations internationales.
\tcblower
\textbf{Solution détaillée :}
\begin{enumerate}
    \item \textbf{Acteurs principaux :} Deux États souverains : la \textbf{République du Cameroun} et la \textbf{République fédérale du Nigeria}. Acteurs tiers : la \textbf{Cour internationale de Justice (CIJ)} (organe judiciaire principal de l'ONU), le \textbf{Secrétaire général de l'ONU} (Kofi Annan), la \textbf{Commission mixte Cameroun-Nigeria} et la \textbf{Commission de suivi de Greentree}.
    \item \textbf{Nature de l'interaction :} Initialement, un \textbf{conflit frontalier} ponctué d'affrontements armés (années 1980-1990) pour le contrôle d'un territoire riche en pétrole et en ressources halieutiques. Suite au recours à la CIJ (1994), l'interaction bascule vers une \textbf{procédure juridictionnelle pacifique} (article 33 de la Charte de l'ONU : règlement pacifique des différends). L'exécution de l'arrêt (2002) par le retrait progressif (Accord de Greentree, 2006) illustre une \textbf{coopération contrainte mais acceptée} sous médiation onusienne.
    \item \textbf{Enjeux :}
    \begin{itemize}
        \item \emph{Territorial et souveraineté :} Définition précise de la frontière terrestre et maritime.
        \item \emph{Économique :} Contrôle des gisements pétroliers offshore et de la zone de pêche.
        \item \emph{Humain :} Sort des populations (Bakassiens), droit d'option, protection des personnes.
        \item \emph{Juridique :} Respect du principe \emph{pacta sunt servanda} (les traités doivent être respectés) et de l'autorité de la chose jugée (article 94 de la Charte de l'ONU).
    \end{itemize}
    \item \textbf{Niveaux d'analyse :} \textbf{Bilatéral} (négociations directes, commission mixte), \textbf{Régional} (stabilité du Golfe de Guinée), \textbf{International/Universel} (CIJ, Conseil de sécurité, Secrétaire général de l'ONU).
    \item \textbf{Conclusion :} L'affaire Bakassi est un \textbf{modèle de règlement pacifique} d'un différend territorial majeur en Afrique. Elle illustre la primauté du droit international sur la force, le rôle central de l'ONU et la complexité de la mise en œuvre (aspects humanitaires, transferts de souveraineté). Le Cameroun a privilégié la \textbf{diplomatie juridique} et le dialogue.
\end{enumerate}
\textbf{Phrase de conclusion :} Ce cas d'école démontre que les relations internationales ne se résument pas à la puissance militaire, mais s'appuient sur un ordre juridique et des institutions capables de transformer un conflit en coopération encadrée.
\end{exoresolu}

\begin{methode}[Classifier les aspects de la coopération internationale]
\textbf{Objectif :} Distinguer les secteurs (aspects) d'intervention de la coopération pour qualifier un accord ou un projet.
\begin{enumerate}
    \item \textbf{Lister les domaines d'action} présents dans le document ou la situation (santé, éducation, infrastructure, défense, culture, environnement, gouvernance...).
    \item \textbf{Rattacher chaque action à son aspect type :}
    \begin{itemize}
        \item \textbf{Technique/Scientifique :} Transfert de technologie, formation d'experts, recherche commune.
        \item \textbf{Financière/Économique :} Prêts, dons, investissements, annulation de dette, appui budgétaire.
        \item \textbf{Humanitaire/Sociale :} Urgence, réfugiés, santé de base, sécurité alimentaire.
        \item \textbf{Culturelle/Éducative :} Bourses, échanges universitaires, patrimoine, médias, sport.
        \item \textbf{Militaire/Sécuritaire :} Formation, équipement, opérations de maintien de la paix, lutte antiterroriste.
        \item \textbf{Institutionnelle/Gouvernance :} Appui à l'État de droit, réforme de l'administration, élections, décentralisation.
    \end{itemize}
    \item \textbf{Justifier le classement} par des mots-clés du texte (ex. : « construction d'un barrage » $\rightarrow$ Technique/Économique ; « envoi de médecins » $\rightarrow$ Technique/Sociale).
    \item \textbf{Synthétiser :} Un projet réel est souvent \textbf{multisectoriel} (ex. : un hôpital = infrastructure technique + formation médicale + financement + gestion administrative).
\end{enumerate}
\cle{Multisectoriel} -- \cle{Appui budgétaire} -- \cle{Transfert de compétences}
\end{methode}

\begin{exoresolu}[Classification d'un programme de coopération Cameroun -- Union européenne]
\textbf{Énoncé :} Le 11\textsuperscript{e} Fonds européen de développement (FED) pour le Cameroun (période 2014-2020) prévoit une enveloppe d'environ 286 millions d'euros axée sur trois secteurs : (1) Gouvernance et État de droit, (2) Développement rural durable et sécurité alimentaire, (3) Connectivité régionale et infrastructures de transport. Un appui budgétaire conditionné aux réformes complète le dispositif. Classez les aspects de coopération mobilisés et justifiez.
\tcblower
\textbf{Solution détaillée :}
\begin{center}
\begin{tabularx}{\linewidth}{|l|X|l|}
\hline
\textbf{Secteur / Action} & \textbf{Justification / Mots-clés} & \textbf{Aspect(s) de coopération} \\
\hline
Gouvernance et État de droit & Réformes judiciaires, lutte contre la corruption, droits de l'homme, appui aux élections, décentralisation. & \textbf{Institutionnelle / Gouvernance} \\
\hline
Développement rural durable & Sécurité alimentaire, agriculture, élevage, gestion des ressources naturelles, résilience climatique. & \textbf{Économique / Technique / Environnementale} \\
\hline
Connectivité régionale & Routes, corridors (ex. : Bamenda-Mamfe-Ekok), facilitation des transports, intégration CEMAC. & \textbf{Technique / Infrastructurale / Économique (Intégration)} \\
\hline
Appui budgétaire conditionné & Transfert direct au Trésor public contre respect d'indicateurs (macroéconomie, secteur social, finances publiques). & \textbf{Financière / Institutionnelle (Conditionnalité)} \\
\hline
\end{tabularx}
\end{center}
\textbf{Analyse transversale :} Ce programme est \textbf{multisectoriel}. L'aspect \textbf{financier} (don, non remboursable) est le vecteur. L'aspect \textbf{institutionnel} est la condition (réformes). Les aspects \textbf{technique} et \textbf{économique} sont les finalités visibles (routes, fermes, tribunaux).
\textbf{Phrase de conclusion :} La coopération Cameroun-UE illustre l'approche \textbf{globale et sectorielle} moderne : l'aide financière sert de levier pour des réformes institutionnelles (gouvernance) qui permettent l'efficacité des investissements techniques (rural, transport).
\end{exoresolu}

\begin{methode}[Différencier les formes de la coopération internationale]
\textbf{Objectif :} Identifier la forme juridique et géopolitique d'un accord de coopération.
\begin{enumerate}
    \item \textbf{Compter le nombre de parties signataires :}
    \begin{itemize}
        \item 2 États (ou 1 État + 1 OI) $\rightarrow$ \textbf{Bilatérale}.
        \item Plusieurs États (conférence, sommet) ou États + OI $\rightarrow$ \textbf{Multilatérale}.
    \end{itemize}
    \item \textbf{Analyser la géographie des partenaires :}
    \begin{itemize}
        \item Pays développés (Nord) $\leftrightarrow$ Pays en développement (Sud) $\rightarrow$ \textbf{Nord-Sud}.
        \item Pays en développement entre eux $\rightarrow$ \textbf{Sud-Sud} (ex. : Cameroun -- Brésil, Cameroun -- Chine, CEMAC).
        \item Pays développés entre eux $\rightarrow$ \textbf{Nord-Nord} (ex. : UE -- USA).
        \item Donateur Nord + pays pivot Sud + bénéficiaire Sud $\rightarrow$ \textbf{Triangulaire}.
    \end{itemize}
    \item \textbf{Identifier le cadre juridique :} Traité, Accord-cadre, Protocole d'accord (MoU), Résolution de l'ONU, Décision d'organisation régionale.
    \item \textbf{Préciser la nature de l'obligation :} Contraignante (traité, règlement) ou non contraignante (déclaration politique, MoU, engagements volontaires).
\end{enumerate}
\cle{Bilatéral} -- \cle{Multilatéral} -- \cle{Nord-Sud} -- \cle{Sud-Sud} -- \cle{Triangulaire} -- \cle{Contraignant}
\end{methode}

\begin{exoresolu}[Typologie d'accords : cas pratiques]
\textbf{Énoncé :} Pour chacun des accords suivants, précisez la forme (bilatérale/multilatérale), l'orientation géographique (Nord-Sud/Sud-Sud/Triangulaire) et le degré d'obligation juridique.
\begin{enumerate}
    \item Accord de siège entre le Cameroun et le PNUD (Programme des Nations Unies pour le développement).
    \item Traité de l'OHADA (Organisation pour l'harmonisation en Afrique du droit des affaires).
    \item Mémorandum d'entente (MoU) Cameroun -- Chine pour la construction du port en eau profonde de Kribi (phase 1).
    \item Initiative « Compact with Africa » du G20 avec le Cameroun (depuis 2018).
\end{enumerate}
\tcblower
\textbf{Solution détaillée :}
\begin{enumerate}
    \item \textbf{Accord de siège Cameroun -- PNUD :}
    \begin{itemize}
        \item \textbf{Forme :} \textbf{Bilatérale} (État hôte $\leftrightarrow$ Organisation internationale).
        \item \textbf{Orientation :} \textbf{Nord-Sud} (PNUD financé principalement par les pays développés, opérant au Sud).
        \item \textbf{Obligation :} \textbf{Contraignante} (accord international : immunités, privilèges, obligations fiscales).
    \end{itemize}
    \item \textbf{Traité OHADA (1993, révisé 2008) :}
    \begin{itemize}
        \item \textbf{Forme :} \textbf{Multilatérale} (17 États parties + institutions communes : CCJA, Commission, UNIDA).
        \item \textbf{Orientation :} \textbf{Sud-Sud} (exclusivement États africains).
        \item \textbf{Obligation :} \textbf{Contraignante} (traité international, droit uniforme directement applicable dans les États parties, primauté sur le droit national).
    \end{itemize}
    \item \textbf{MoU Cameroun -- Chine (Port de Kribi) :}
    \begin{itemize}
        \item \textbf{Forme :} \textbf{Bilatérale} (Gouvernement camerounais $\leftrightarrow$ Gouvernement chinois / Exim Bank of China / entreprises chinoises).
        \item \textbf{Orientation :} \textbf{Sud-Sud} (Chine = pays émergent, Cameroun = pays en développement).
        \item \textbf{Obligation :} \textbf{Non contraignante (politique/morale)} pour le MoU lui-même. \textbf{Attention :} les contrats commerciaux (EPC : Engineering, Procurement, Construction) et accords de prêt qui en découlent sont, eux, \textbf{contraignants}.
    \end{itemize}
    \item \textbf{Compact with Africa (G20 -- Cameroun) :}
    \begin{itemize}
        \item \textbf{Forme :} \textbf{Multilatérale} (cadre G20 + institutions de Bretton Woods + pays africains partenaires + investisseurs privés).
        \item \textbf{Orientation :} \textbf{Triangulaire} (le G20 facilite / les institutions internationales encadrent / le Cameroun réforme / les privés investissent).
        \item \textbf{Obligation :} \textbf{Non contraignante} (engagement politique, pas de sanction juridique directe, mais conditionnalité pour l'accès aux financements).
    \end{itemize}
\end{enumerate}
\textbf{Phrase de conclusion :} La distinction forme/orientation/obligation est cruciale : un MoU Sud-Sud (Kribi) engage financièrement via des contrats dérivés, tandis qu'un traité multilatéral Sud-Sud (OHADA) crée un ordre juridique directement applicable.
\end{exoresolu}

\begin{methode}[Évaluer une action d'assistance humanitaire (Dossier 1)]
\textbf{Objectif :} Juger de la conformité d'une intervention humanitaire aux 4 principes fondamentaux (ONU, OCHA, CICR).
\begin{enumerate}
    \item \textbf{Rappeler les 4 piliers (HNII) :}
    \begin{itemize}
        \item \textbf{Humanité :} Secourir les vies, soulager les souffrances, respecter la dignité (priorité aux plus vulnérables).
        \item \textbf{Neutralité :} Ne pas prendre parti dans les hostilités, ni dans les controverses politiques, religieuses ou idéologiques.
        \item \textbf{Impartialité :} Agir \textbf{uniquement selon les besoins}, sans discrimination (nationalité, race, sexe, religion, opinion). Priorité = urgence vitale.
        \item \textbf{Indépendance :} Autonomie des objectifs humanitaires par rapport aux objectifs politiques, économiques, militaires ou autres des acteurs.
    \end{itemize}
    \item \textbf{Tester l'action au crible de chaque principe :} Y a-t-il discrimination ? Prise de parti ? Instrumentalisation politique ? Conditionnalité d'accès ?
    \item \textbf{Identifier les dilemmes :} Accès humanitaire vs souveraineté (consentement de l'État) ; neutralité vs témoignage (plaidoyer) ; sécurité du personnel vs accès aux victimes.
    \item \textbf{Conclure :} L'action est-elle « puriste » (CICR/MSF) ou « intégrée » (ONU/OCHA en mission intégrée) ? Quels risques pour les principes ?
\end{enumerate}
\cle{Humanité} -- \cle{Neutralité} -- \cle{Impartialité} -- \cle{Indépendance} -- \cle{Accès humanitaire} -- \cle{Instrumentalisation}
\end{methode}

\begin{exoresolu}[Dilemme humanitaire : accès aux zones NOSO (Nord-Ouest/Sud-Ouest du Cameroun)]
\textbf{Énoncé :} Depuis 2017, des ONG internationales (MSF, CICR) et des agences de l'ONU (OCHA) opèrent dans les régions anglophones en crise. Le gouvernement exige une coordination stricte et des autorisations de mouvement, et accuse certaines ONG de partialité. Les groupes armés non étatiques (GANE) imposent des « villes mortes », bloquent les routes et taxent l'aide. En 2021, MSF suspend ses activités dans le NOSO après l'arrestation de membres de son personnel et l'attaque d'un convoi. Analysez le respect des principes humanitaires par les parties prenantes.
\tcblower
\textbf{Solution détaillée :}
\begin{center}
\begin{tabularx}{\linewidth}{|l|X|X|}
\hline
\textbf{Principe} & \textbf{Respect / Violation par l'État (gouvernement/forces de défense)} & \textbf{Respect / Violation par les GANE / contexte} \\
\hline
\textbf{Humanité} & \textbf{Violation partielle :} restrictions d'accès (couvre-feu, interdiction des motos), attaques de structures de santé. & \textbf{Violation grave :} attaques directes sur le personnel humanitaire (enlèvements, meurtres), pillage de dépôts, « villes mortes » empêchant l'accès aux soins. \\
\hline
\textbf{Neutralité} & \textbf{Risque d'instrumentalisation :} conditionner l'accès à l'étiquetage des bénéficiaires (« population civile » vs « séparatistes »). & \textbf{Violation :} perception de l'aide comme soutien à l'ennemi ; taxation de l'aide = détournement à des fins de guerre. \\
\hline
\textbf{Impartialité} & \textbf{Tension :} pression pour orienter l'aide vers les zones contrôlées par l'armée ; refus d'accès aux zones tenues par les GANE. & \textbf{Violation :} les GANE interdisent l'aide aux populations perçues comme « collaboratrices » ou aux zones tenues par l'armée. \\
\hline
\textbf{Indépendance} & \textbf{Violation :} subordination de l'action humanitaire à l'agenda sécuritaire/militaire ; arrestation de personnel de MSF. & \textbf{Violation :} tentative de contrôle de l'aide par les GANE pour la légitimité politique ou la ressource financière. \\
\hline
\end{tabularx}
\end{center}
\textbf{Synthèse :} L'espace humanitaire se \textbf{rétrécit drastiquement} (\emph{shrinking humanitarian space}). La suspension de MSF (2021) est la conséquence directe de l'impossibilité de garantir l'\textbf{indépendance} et la \textbf{sécurité} (condition de l'humanité).
\textbf{Nuance importante :} Le CICR maintient un dialogue confidentiel bilatéral avec \textbf{toutes} les parties (État + GANE) pour préserver l'accès, illustrant la \textbf{neutralité opérationnelle} stricte. L'ONU (OCHA) cherche une coordination civilo-militaire (CMCoord) plus difficile.
\textbf{Phrase de conclusion :} Dans les conflits internes asymétriques, le respect des 4 principes exige un \textbf{dialogue humanitaire constant} avec les porteurs d'armes ; toute confusion entre objectifs militaires/politiques et aide humanitaire tue l'action humanitaire elle-même.
\end{exoresolu}

\begin{methode}[Appliquer les principes de la diplomatie camerounaise]
\textbf{Objectif :} Justifier la position ou l'action du Cameroun dans une situation internationale donnée en s'appuyant sur ses constantes doctrinales.
\begin{enumerate}
    \item \textbf{Lister les 5 piliers officiels (Charte de l'ONU, Acte constitutif de l'UA, Constitution, discours officiels) :}
    \begin{enumerate}
        \item \textbf{Non-alignement / Neutralité active :} refus des blocs, liberté de vote, médiation.
        \item \textbf{Non-ingérence / Égalité souveraine :} respect de l'intégrité territoriale, intangibilité des frontières (héritage colonial), règlement pacifique.
        \item \textbf{Coopération / Solidarité africaine :} intégration sous-régionale (CEMAC, CEEAC), panafricanisme, paix et sécurité (Force multinationale mixte MNJTF).
        \item \textbf{Dialogue / Négociation / Règlement pacifique :} primauté du droit (CIJ, OHADA), diplomatie préventive.
        \item \textbf{Respect du droit international / Engagement multilatéral :} ratification des traités, contribution de Casques bleus, candidature aux organes de l'ONU.
    \end{enumerate}
    \item \textbf{Analyser la situation :} Quel principe est mobilisé ? Y a-t-il tension entre principes ? (ex. : non-ingérence vs droits de l'homme / paix).
    \item \textbf{Argumenter :} « Le Cameroun agit conformément au principe de [X] car [fait précis]... ». Citer un texte (Préambule de la Constitution, discours UA/ONU, traité).
    \item \textbf{Nuancer :} La diplomatie est pragmatique : alliances de circonstance, intérêts nationaux (sécurité, économie), réalisme.
\end{enumerate}
\cle{Non-alignement} -- \cle{Non-ingérence} -- \cle{Intangibilité des frontières} -- \cle{Règlement pacifique} -- \cle{Pragmatisme}
\end{methode}

\begin{exoresolu}[Position du Cameroun face à la crise en RCA (2013-2014) et la MNJTF]
\textbf{Énoncé :} En 2013, la RCA s'effondre (Séléka, Anti-Balaka). Le Cameroun, frontalier, accueille des centaines de milliers de réfugiés, ferme temporairement sa frontière, déploie le BIR à la frontière, et participe à la médiation pour la transition (Djotodia $\rightarrow$ Samba-Panza). Parallèlement, le Cameroun est un contributeur majeur de la Force multinationale mixte (MNJTF) contre Boko Haram dans le bassin du Lac Tchad. Montrez comment ces deux dossiers illustrent les principes de la diplomatie camerounaise.
\tcblower
\textbf{Solution détaillée :}

\textbf{Dossier 1 : Crise RCA (volet humanitaire, sécuritaire, médiation)}
\begin{itemize}
    \item \textbf{Solidarité africaine / Humanité :} accueil massif de réfugiés (politique de « porte ouverte » malgré la pression sur les ressources), aide logistique (vivres, santé).
    \item \textbf{Sécurité nationale / Non-ingérence (protection) :} fermeture de la frontière (mesure de souveraineté, protection contre l'infiltration d'armes et de combattants), déploiement du BIR (défense du territoire). \textbf{La non-ingérence} ne signifie pas la passivité face à la menace extérieure.
    \item \textbf{Règlement pacifique / Médiation :} rôle de la CEEAC (Conférence des chefs d'État) pour imposer une transition constitutionnelle (Samba-Panza) et organiser les élections de 2015-2016. \textbf{Diplomatie du dialogue} en action.
    \item \textbf{Multilatéralisme :} contribution de troupes et de policiers à la MINUSCA (Casques bleus), coordination ONU/UA/CEEAC.
\end{itemize}

\textbf{Dossier 2 : MNJTF / Boko Haram (volet sécurité collective, intégration)}
\begin{itemize}
    \item \textbf{Coopération / Solidarité sous-régionale :} engagement camerounais (mise à disposition de troupes : BIR, renseignement, logistique), mutualisation des moyens (Nigeria, Tchad, Niger, Bénin).
    \item \textbf{Intangibilité des frontières / Souveraineté :} lutte contre un groupe terroriste transfrontalier niant les frontières. Opérations conjointes transfrontalières (nécessitant des accords bilatéraux/multilatéraux pour le « droit de poursuite »).
    \item \textbf{Non-alignement / Partenariats diversifiés :} appui logistique et de renseignement de la France, des USA (formation), équipements divers. Liberté de choix des partenaires.
    \item \textbf{Droit international / Droits de l'homme :} engagement de respect du droit international humanitaire par les forces nationales.
\end{itemize}

\textbf{Synthèse comparative :}
\begin{center}
\begin{tabularx}{\linewidth}{|l|X|X|}
\hline
\textbf{Principe} & \textbf{Crise RCA (volet médiation/aide)} & \textbf{MNJTF (volet sécurité collective)} \\
\hline
Non-ingérence & Respect de la souveraineté de la RCA (médiation acceptée) & Action conjointe \textbf{avec consentement} (mandat MNJTF 2014-2015) = non-violation. \\
\hline
Règlement pacifique & \textbf{Cœur de l'action} (médiation politique) & \textbf{Préliminaire :} la force crée les conditions de la paix (négociation ultérieure). \\
\hline
Solidarité africaine & Accueil de réfugiés + médiation + Casques bleus & \textbf{Archétype :} mutualisation des forces, financement partagé. \\
\hline
Pragmatisme / Intérêt national & Stabilité de la frontière = sécurité du Cameroun + image de leader & Protection de l'arrière-pays (Extrême-Nord) + leadership régional. \\
\hline
\end{tabularx}
\end{center}
\textbf{Phrase de conclusion :} La diplomatie camerounaise allie \textbf{fidélité aux principes} (droit, paix, solidarité) et \textbf{réalisme sécuritaire} (défense des intérêts vitaux, leadership régional assumé). Le Cameroun est un « État-pivot » en Afrique centrale : sa stabilité conditionne celle de la sous-région.
\end{exoresolu}

\begin{methode}[Structurer une réponse argumentée type Bac (rédaction / justification)]
\textbf{Objectif :} Rédiger une réponse structurée, juridiquement fondée et nuancée pour une question de dissertation ou d'analyse de document.
\begin{enumerate}
    \item \textbf{Analyser le sujet :} mots-clés, instruction (Montrer, Expliquer, Dans quelle mesure, Apprécier), délimitation (période, acteurs).
    \item \textbf{Problématiser :} transformer le sujet en question ouverte ou en tension (ex. : « Le principe de non-ingérence est-il compatible avec l'obligation de protection des populations ? »).
    \item \textbf{Construire le plan (2 ou 3 parties équilibrées) :}
    \begin{itemize}
        \item \textbf{Thèse (I) :} l'affirmation du principe / la réalité de l'action.
        \item \textbf{Antithèse (II) :} les limites / les exceptions / les tensions / le pragmatisme.
        \item \textbf{Synthèse (III, si dissertation) :} l'évolution / l'articulation / le modèle camerounais spécifique.
    \end{itemize}
    \item \textbf{Rédiger par paragraphes argumentés (A.R.E.) :}
    \begin{itemize}
        \item \textbf{A}ffirmation (phrase-thèse du paragraphe).
        \item \textbf{R}éférence / Preuve (texte constitutionnel, traité, fait historique précis, chiffre, citation officielle).
        \item \textbf{E}xplication / Analyse (lien avec le problème, portée, nuance).
    \end{itemize}
    \item \textbf{Conclure :} réponse directe à la problématique + ouverture (perspective future, défi actuel).
\end{enumerate}
\cle{Problématique} -- \cle{Plan dialectique} -- \cle{Argumentation ARE} -- \cle{Preuve constitutionnelle}
\end{methode}

\begin{exoresolu}[Sujet type Bac : « Le Cameroun, acteur de la paix et de la sécurité en Afrique centrale : réalités et limites »]
\textbf{Énoncé :} En vous appuyant sur vos connaissances et sur les principes de la diplomatie camerounaise, montrez que le Cameroun est un acteur majeur de la paix et de la sécurité en Afrique centrale, tout en soulignant les limites de son action.
\tcblower
\textbf{Solution détaillée :}

\textbf{Problématique :} Comment le Cameroun, par sa diplomatie et son engagement militaire, s'impose-t-il comme un pilier de la stabilité en Afrique centrale, malgré les contraintes structurelles et les contradictions inhérentes à son rôle régional ?

\textbf{Introduction :}
\emph{Accroche :} « L'Afrique centrale ne peut se construire sans le Cameroun » (idée récurrente des discours officiels camerounais).
\emph{Définition :} Diplomatie camerounaise = non-alignement, règlement pacifique, solidarité africaine, intangibilité des frontières (Préambule de la Constitution).
\emph{Annonce du plan :} I. Un leadership diplomatique et militaire assumé (réalités). II. Des limites structurelles et contextuelles qui fragilisent l'action (limites). III. Vers une refondation du leadership par l'intégration et la gouvernance (synthèse/perspective).

\textbf{Développement :}

\textbf{I. Le Cameroun, pilier de la paix et de la sécurité : un leadership multidimensionnel (réalités)}
\begin{itemize}
    \item \textbf{A. La médiation politique et le règlement pacifique des différends (principe : dialogue/négociation).}
    \item \emph{Preuve :} rôle dans la transition en RCA (2013-2016), participation aux médiations sous-régionales, facilitation du dialogue entre voisins.
    \item \emph{Analyse :} la « diplomatie du dialogue » fait du Cameroun un médiateur de confiance grâce à sa stabilité relative et à son non-alignement.
    \item \textbf{B. L'engagement militaire au service de la sécurité collective (principe : solidarité africaine / intégration).}
    \item \emph{Preuve :} contribution importante de troupes et de policiers à la MINUSCA (RCA) et à la MNJTF (Lac Tchad). Ratification des traités CEMAC/CEEAC (protocoles de non-agression et de défense mutuelle).
    \item \emph{Analyse :} la mutualisation des forces (MNJTF) concrétise la sécurité collective (article 4 de l'Acte constitutif de l'UA). Le Cameroun assume une part majeure du fardeau (\emph{burden sharing}).
    \item \textbf{C. L'ancrage institutionnel et normatif (principe : droit international / intangibilité des frontières).}
    \item \emph{Preuve :} ratification quasi systématique des instruments UA/ONU ; résolution pacifique de Bakassi (CIJ, 2002) = modèle jurisprudentiel pour la région.
\end{itemize}

\textbf{II. Les limites structurelles et contextuelles de l'action camerounaise (limites)}
\begin{itemize}
    \item \textbf{A. La crise interne (NOSO) comme paradoxe du « gendarme » régional (tension souveraineté / droits de l'homme).}
    \item \emph{Preuve :} crise anglophone depuis 2016 (milliers de morts, centaines de milliers de déplacés, écoles fermées). Allégations de violations des droits de l'homme (Amnesty International, Human Rights Watch). Suspension partielle de l'accès du Cameroun à l'AGOA (2019) par les USA.
    \item \emph{Analyse :} perte de légitimité morale pour donner des leçons de gouvernance aux pairs. Le principe de non-ingérence invoqué pour rejeter les enquêtes internationales $\rightarrow$ isolement diplomatique partiel.
    \item \textbf{B. La dépendance sécuritaire et financière vis-à-vis de partenaires extra-régionaux (limite du non-alignement / autonomie).}
    \item \emph{Preuve :} appui français (formation du BIR), américain (formation, renseignement), chinois (infrastructures). L'endettement infrastructurel (Kribi, autoroutes) limite la marge de manœuvre budgétaire de la défense.
    \item \emph{Analyse :} le « non-alignement actif » est mis à l'épreuve par la nécessité d'équipements de haute technologie non produits localement. Risque d'alignement de facto sur les agendas antiterroristes occidentaux.
    \item \textbf{C. Les dysfonctionnements de l'intégration sous-régionale (limite de la solidarité / efficacité).}
    \item \emph{Preuve :} CEEAC longtemps paralysée (crise de financement, leadership), CEMAC : non-respect des critères de convergence, libre circulation des personnes et des biens incomplète, marché commun imparfait.
    \item \emph{Analyse :} le Cameroun finance une part importante des budgets CEMAC/CEEAC sans récolter pleinement les dividendes de la paix (insécurité persistante aux frontières Nord et Est). Asymétrie : le Cameroun est contributeur net, d'autres sont bénéficiaires nets de la sécurité.
\end{itemize}

\textbf{III. Perspectives : refonder le leadership par l'intégration profonde et la gouvernance démocratique (synthèse)}
\begin{itemize}
    \item \textbf{A. La paix intérieure, condition de la paix extérieure :} résolution inclusive de la crise NOSO (Grand Dialogue national 2019, statut spécial, DDR effectif) pour restaurer la crédibilité éthique.
    \item \textbf{B. Opérationnaliser la CEEAC/CEMAC :} mise en œuvre effective de la Force en attente (FOMAC), libre circulation réelle, marché commun de l'énergie et des transports.
    \item \textbf{C. La diplomatie économique comme levier de sécurité :} ZLECAf (Zone de libre-échange continentale africaine) : transformer l'interdépendance économique en garantie de paix.
\end{itemize}

\textbf{Conclusion :}
Le Cameroun est indéniablement l'\textbf{« État-charnière »} de l'Afrique centrale (géographie, démographie, économie, armée). Son action pour la paix (médiation, MNJTF, Casques bleus, Bakassi) est \textbf{réelle, coûteuse et conforme à ses principes}. Cependant, la \textbf{crise anglophone interne} et la \textbf{fragilité de l'intégration régionale} constituent des failles majeures. L'avenir de son leadership dépend de sa capacité à \textbf{résoudre son conflit interne par le dialogue} (cohérence principes/pratique) et à \textbf{faire de l'intégration économique (ZLECAf/CEMAC) le socle durable de la sécurité collective}.
\end{exoresolu}

\begin{methode}[Cartographier et argumenter l'action internationale (compétence TP)]
\textbf{Objectif :} Réaliser un croquis commenté (carte mentale ou carte géographique simplifiée) et l'exploiter à l'oral ou à l'écrit.
\begin{enumerate}
    \item \textbf{Choisir le fond de carte :} Afrique centrale (Cameroun centré) + monde (pour les partenariats globaux). Simplifier les contours.
    \item \textbf{Construire la légende (3 à 4 items max) :}
    \begin{itemize}
        \item \textbf{Figé ponctuel (ronds/symboles) :} missions de Casques bleus (taille $\propto$ effectifs), sièges d'institutions, sommets accueillis.
        \item \textbf{Figé linéaire (flèches) :} flux de coopération (épaisseur $\propto$ volume : Chine, France, UE, USA, Nigeria, CEMAC). Axes routiers/ferroviaires d'intégration (corridors).
        \item \textbf{Figé surfacique (zones colorées) :} espaces d'influence diplomatique (CEMAC, CEEAC, Golfe de Guinée, bassin du Lac Tchad, Commonwealth/Francophonie/OIF).
        \item \textbf{Signes spécifiques :} conflits gérés (Bakassi $\checkmark$, NOSO $\times$, RCA $\rightarrow$, Boko Haram $\rightarrow$).
    \end{itemize}
    \item \textbf{Respecter la sémiologie :} titre précis, auteur/source/date, orientation (Nord), échelle (qualitative), légende organisée (hiérarchie visuelle).
    \item \textbf{Argumenter à partir de la carte :} « Cette carte montre que... », « L'épaisseur de la flèche Chine illustre... », « L'enclavement du Nord justifie... ».
\end{enumerate}
\cle{Sémiologie graphique} -- \cle{Figé ponctuel/linéaire/surfacique} -- \cle{Proportionnalité} -- \cle{Hiérarchie visuelle}
\end{methode}

\begin{exoresolu}[TP : croquis commenté « Le Cameroun dans les relations internationales (2020-2024) »]
\textbf{Énoncé :} Réalisez un croquis schématique (description textuelle structurée) de l'action internationale du Cameroun sur la période récente. Légendez-le et rédigez un commentaire argumenté de 15 lignes montrant la diversification des partenariats et la centralité sécuritaire.
\tcblower
\textbf{Solution détaillée :}

\textbf{1. Description du croquis (à reproduire sur copie) :}
\begin{itemize}
    \item \textbf{Fond de carte :} Afrique centrale (Cameroun au centre, frontières avec 6 pays + mer). Monde en médaillon (coin supérieur droit) pour les partenaires globaux.
    \item \textbf{Zones colorées (surfacique) :}
    \begin{itemize}
        \item \textcolor{popblue}{\textbf{Bleu clair :}} espace CEMAC/CEEAC (intégration prioritaire) -- hachuré vertical.
        \item \textcolor{popgreen}{\textbf{Vert clair :}} espace bassin du Lac Tchad / MNJTF (sécurité urgente) -- hachuré horizontal.
        \item \textcolor{popgold}{\textbf{Jaune clair :}} Golfe de Guinée / zone maritime (économie bleue, sécurité maritime) -- pointillé.
    \end{itemize}
    \item \textbf{Flux (linéaire -- flèches proportionnelles) :}
    \begin{itemize}
        \item \textbf{Flèche ÉPAISSE rouge (bidirectionnelle) :} Cameroun $\leftrightarrow$ \textbf{Chine} (infrastructures de Kribi, autoroutes, santé, défense, FOCAC).
        \item \textbf{Flèche MOYENNE bleue (bidirectionnelle) :} Cameroun $\leftrightarrow$ \textbf{France / UE} (aide publique au développement, formation, défense, commerce, accords de partenariat).
        \item \textbf{Flèche MOYENNE verte (sortante) :} Cameroun $\rightarrow$ \textbf{RCA / Tchad / Nigeria} (aide humanitaire, médiation, sécurité MNJTF, énergie).
        \item \textbf{Flèches FINES orange (multidirectionnelles) :} Cameroun $\leftrightarrow$ \textbf{USA, Russie, Turquie, Maroc, pays du Golfe} (défense ciblée, agriculture, éducation, investissement).
    \end{itemize}
    \item \textbf{Points / symboles (ponctuel) :}
    \begin{itemize}
        \item \textbf{Étoile rouge (Yaoundé) :} capitale diplomatique (ambassades, ministères, institutions).
        \item \textbf{Casque bleu (Nord/Extrême-Nord/Est) :} bases BIR / MNJTF / MINUSCA (déploiement sortant).
        \item \textbf{Balance (frontière de Bakassi) :} symbole du règlement pacifique (CIJ).
        \item \textbf{Éclair rouge (NOSO) :} foyer de tension interne à résonance internationale.
    \end{itemize}
\end{itemize}

\textbf{2. Commentaire argumenté (modèle de réponse Bac) :}

\emph{Ce croquis révèle la double dimension de l'action internationale camerounaise : un \textbf{ancrage régional obligatoire} et une \textbf{diversification partenariale stratégique}.}

\emph{La \textbf{centralité sécuritaire} saute aux yeux : l'épaisseur des flèches vers le Nord (MNJTF/Boko Haram) et l'Est (RCA/MINUSCA), ainsi que les zones hachurées (Lac Tchad, CEEAC), prouvent que la survie de l'État et la stabilité sous-régionale sont les moteurs premiers. Le Cameroun assume son rôle de \textbf{fournisseur net de sécurité publique régionale} (principes : solidarité, règlement pacifique, intangibilité des frontières).}

\emph{Parallèlement, la \textbf{diversification des partenariats} (flèches Chine, France/UE, nouveaux partenaires) illustre le \textbf{non-alignement actif} : le Cameroun refuse la dépendance unique et joue la concurrence pour maximiser les gains (infrastructures vs gouvernance, défense vs développement). La flèche Chine (la plus épaisse) marque le virage infrastructurel et économique ; la flèche UE/France maintient le lien historique et de gouvernance.}

\emph{Toutefois, le \textbf{symbole d'alerte (NOSO)} et la \textbf{fragilité des zones d'intégration} (hachures discontinues) rappellent que la projection extérieure est menacée par la cohésion interne et l'inachèvement de l'intégration (libre circulation bloquée, insécurité persistante). Le défi est de transformer la centralité géographique en \textbf{puissance normative et économique} durable.}
\end{exoresolu}

\begin{attention}[Les 5 erreurs qui coûtent des points au Bac en ECM (diplomatie/RI)]
\begin{enumerate}
    \item \textbf{Confondre « non-ingérence » et « indifférence / isolationnisme ».} \textbf{Correction :} le Cameroun intervient (médiation, Casques bleus, aide humanitaire) \textbf{AVEC le consentement} de l'État hôte ou un mandat ONU/UA. C'est une ingérence \textbf{autorisée / sollicitée}, interdite seulement si unilatérale et coercitive.
    \item \textbf{Mélanger « bilatéral » et « Nord-Sud » / « multilatéral » et « Sud-Sud ».} \textbf{Correction :} ce sont deux grilles de lecture \textbf{indépendantes}. L'OHADA est multilatérale ET Sud-Sud. L'accord Cameroun-Chine est bilatéral ET Sud-Sud. L'UE-Cameroun est bilatéral (accord de partenariat) MAIS l'UE est une organisation multilatérale $\rightarrow$ préciser : « cadre bilatéral UE-pays » vs « cadre multilatéral ONU ».
    \item \textbf{Oublier la « primauté du droit international » comme principe camerounais.} \textbf{Correction :} citer \textbf{l'affaire Bakassi (CIJ 2002)} est LA preuve majeure. Ne pas la citer affaiblit l'argumentation sur le « règlement pacifique » et le « respect du droit ». Le Préambule de la Constitution y fait référence (adhésion aux principes de la Charte de l'ONU).
    \item \textbf{Traiter l'assistance humanitaire comme de la « coopération au développement » classique.} \textbf{Correction :} l'humanitaire = \textbf{urgence, vie, principes HNII (neutralité/impartialité)}. La coopération = \textbf{long terme, développement structurel, conditionnalité politique possible}. Ne pas dire « l'aide humanitaire vise le développement » : elle vise la \textbf{survie et la dignité immédiate}.
    \item \textbf{Rédiger une « liste de courses » au lieu d'une argumentation structurée (thèse/antithèse/synthèse).} \textbf{Correction :} pour une question « Montrez que... », « Appréciez... », il faut un \textbf{plan dialectique} (I. Réalités / II. Limites / III. Perspectives) avec des \textbf{arguments ARE (Affirmation -- Référence -- Explication)}. Une énumération d'exemples (Bakassi, MNJTF, RCA, Chine, France) sans articulation logique ne vaut que 1 à 2 points sur 6.
\end{enumerate}
\end{attention}

\newpage

\coursec{Exercices}

\courssub{Je vérifie mes connaissances}

\begin{exercice}[QCM : Notions fondamentales]
\emph{Choisissez la unique bonne réponse pour chaque question.}
\begin{enumerate}
    \item Les relations internationales désignent : \\
    A. Les échanges commerciaux entre deux entreprises privées. \\
    B. L'ensemble des interactions (politiques, économiques, culturelles, militaires) entre acteurs étatiques et non étatiques. \\
    C. Uniquement les réunions des chefs d'État à l'ONU. \\
    D. La politique intérieure d'un pays.

    \item La coopération \cle{bilatérale} se caractérise par : \\
    A. L'implication d'une organisation internationale comme l'ONU. \\
    B. Un accord entre deux États souverains. \\
    C. L'absence de traité écrit. \\
    D. La participation de tous les pays du continent.

    \item L'\cle{assistance humanitaire} se distingue de l'aide au développement par : \\
    A. Son caractère obligatoire pour les États donateurs. \\
    B. Son objectif d'urgence et de sauvegarde de vies humaines à court terme. \\
    C. Son financement exclusif par le secteur privé. \\
    D. Sa durée illimitée.

    \item Le principe de \cle{non-ingérence} dans la diplomatie camerounaise implique : \\
    A. De ne pas voter les résolutions de l'ONU. \\
    B. De refuser toute aide extérieure. \\
    C. De ne pas intervenir dans les affaires intérieures relevant de la souveraineté d'un autre État. \\
    D. De ne pas entretenir de relations diplomatiques.
\end{enumerate}
\end{exercice}

\begin{exercice}[Vrai / Faux justifié]
\emph{Indiquez si l'affirmation est vraie (V) ou fausse (F) et justifiez brièvement votre réponse en une phrase.}
\begin{enumerate}
    \item Le Cameroun pratique la politique de \cle{non-alignement} en refusant d'adhérer à toute organisation internationale.
    \item La \cle{coopération multilatérale} met nécessairement en jeu plus de deux États, souvent au sein d'une organisation internationale.
    \item Un \cle{réfugié} (au sens de la Convention de Genève de 1951) a franchi une frontière internationale, contrairement au \cle{déplacé interne}.
    \item Le règlement pacifique du différend de la \cle{péninsule de Bakassi} par la Cour internationale de Justice (CIJ) illustre le principe du respect du droit international.
\end{enumerate}
\end{exercice}

\begin{exercice}[Définitions et distinctions]
\emph{Définissez précisément les couples de notions suivants en soulignant la différence majeure :}
\begin{enumerate}
    \item \cle{Relations internationales} vs \cle{Politique étrangère}.
    \item \cle{Coopération bilatérale} vs \cle{Coopération multilatérale}.
    \item \cle{Aide humanitaire d'urgence} vs \cle{Aide au développement}.
    \item \cle{Diplomatie préventive} vs \cle{Diplomatie coercitive}.
\end{enumerate}
\end{exercice}

\begin{exercice}[Classement rapide : Aspects et Formes]
\emph{Complétez le tableau ci-dessous en classant chaque exemple selon son \textbf{aspect} (Économique, Culturel, Technique/Scientifique, Militaire/Sécuritaire, Humanitaire) et sa \textbf{forme} (Bilatérale ou Multilatérale).}

\begin{center}
\begin{tabularx}{\linewidth}{|l|X|X|}\hline
\textbf{Exemple de coopération} & \textbf{Aspect} & \textbf{Forme} \\ \hline
1. Construction du port en eau profonde de Kribi (avec Chine) & & \\ \hline
2. Participation des casques bleus camerounais à la MINUSCA (RCA) & & \\ \hline
3. Attribution de bourses d'études par la France (programme Eiffel) & & \\ \hline
4. Intégration du Cameroun à la \cle{ZLECAf} (Zone de Libre-Échange Continentale Africaine) & & \\ \hline
5. Gestion du camp de réfugiés de Minawao (HCR, ONU, ONG, État camerounais) & & \\ \hline
6. Accord de défense signé avec les États-Unis (formation, matériel) & & \\ \hline
7. Projet de recherche agricole avec l'IITA (Institut International d'Agriculture Tropicale) & & \\ \hline
8. Festival culturel « FESTAC » organisé avec le Nigeria & & \\ \hline
9. Appui budgétaire de l'Union Européenne (FED) & & \\ \hline
10. Mise en place de la force multinationale mixte (FMM) contre Boko Haram (LCBC) & & \\ \hline
\end{tabularx}
\end{center}
\end{exercice}

\courssub{Je m'entraîne}

\begin{exercice}[Question à réponse construite : Synthèse des notions]
\emph{Traitez le sujet suivant en une vingtaine de lignes structurées (introduction, développement en deux paragraphes, conclusion).}
\begin{quote}
« Définissez la notion de \cle{relations internationales}. Présentez ensuite \textbf{deux aspects} et \textbf{deux formes} de la coopération internationale en vous appuyant sur des exemples concrets tirés de l'action extérieure du Cameroun. »
\end{quote}
\end{exercice}

\begin{exercice}[Étude de document : L'assistance humanitaire à Minawao]
\emph{Document : Extrait de rapport HCR 2023 — Camp de Minawao (Région de l'Extrême-Nord).}
\begin{quote}
« Le camp de Minawao, ouvert en 2013, accueille environ 70 000 réfugiés nigérians fuyant les exactions de Boko Haram. L'assistance est coordonnée par le HCR avec le gouvernement camerounais (MINAS, MINSANTE) et des ONG partenaires (Plan International, IMC, CARE). Les défis majeurs restent l'accès à l'eau potable (ratio 15 L/personne/jour visé), la scolarisation des enfants (taux de fréquentation \textasciitilde 60\%), la protection contre les violences basées sur le genre (VBG) et la dépendance prolongée à l'aide alimentaire (PAM) faute d'activités génératrices de revenus suffisantes. La cohabitation avec les populations d'accueil génère parfois des tensions foncières et sur les ressources naturelles. »
\end{quote}
\emph{Questions :}
\begin{enumerate}
    \item À partir du texte, donnez une définition opérationnelle de l'\cle{assistance humanitaire}.
    \item Identifiez \textbf{trois catégories d'acteurs} intervenant à Minawao et précisez le rôle spécifique de chacun.
    \item Dégagez \textbf{deux défis majeurs} de l'action humanitaire dans ce contexte et proposez une piste de solution durable pour chacun.
\end{enumerate}
\end{exercice}

\begin{exercice}[Application : Principes directeurs de la diplomatie camerounaise]
\emph{La Constitution du Cameroun (Préambule) et la Charte de l'ONU fondent l'action extérieure.}
\begin{enumerate}
    \item Citez \textbf{cinq principes cardinaux} de la diplomatie camerounaise (ex : respect de la souveraineté, non-ingérence, règlement pacifique des différends, non-alignement, coopération solidaire...).
    \item Pour \textbf{chacun de ces cinq principes}, citez une action concrète ou une position officielle du Cameroun (depuis 1960) qui l'illustre.
    \item Expliquez en quoi le principe de \cle{non-alignement} ne signifie pas \cle{neutralité} ni \cle{isolement}, en vous appuyant sur l'appartenance du Cameroun à la Francophonie, au Commonwealth et à l'UA.
\end{enumerate}
\end{exercice}

\begin{exercice}[Travail pratique noté : Croquis « Le Cameroun et ses partenaires dans le monde »]
\emph{Consigne : À partir des données ci-dessous, organisez la \textbf{légende complète} d'un croquis planisphérique (type Bac) intitulé « Le Cameroun et ses partenaires dans le monde ». Vous ne dessinez pas la carte, vous structurez la légende en 3 parties (Flux économiques, Coopérations politiques/sécuritaires, Migrations/Réfugiés) avec symboles, couleurs et largeurs de flèches proportionnelles.}

\begin{quote}
\textbf{Données fournies :}
\begin{itemize}
    \item \textbf{Exportations (2022) :} UE (45\%), Chine (15\%), Inde (10\%), USA (5\%), Afrique (CEMAC/CEEAC 10\%).
    \item \textbf{Importations (2022) :} Chine (25\%), UE (20\%), Nigeria (15\%), Inde (5\%), Autres.
    \item \textbf{Investissements directs étrangers (IDE) stock :} France (historique), Chine (infrastructures : barrages, routes, port Kribi), USA (pétrole/gaz), Maroc (banque, télécoms).
    \item \textbf{Coopération militaire :} France (accords de défense, formation EMIA), USA (formation, renseignement, drones), Russie (équipements historiques), Israël (modernisation), Chine (fournitures).
    \item \textbf{Missions de paix (ONU/U.A.) :} Centrafrique (MINUSCA), Mali (ex-MINUSMA), Soudan, Haïti (police).
    \item \textbf{Flux de réfugiés (HCR 2023) :} Entrants : Nigeria (\textasciitilde 120 000, Extrême-Nord), RCA (\textasciitilde 350 000, Est/Adamaoua). Sortants : Camerounais déplacés internes (NO/SO \textasciitilde 600 000 ; Extrême-Nord \textasciitilde 300 000) et réfugiés au Nigeria/Tchad.
    \item \textbf{Diplomatie culturelle/éducative :} Bourses France, Chine, Russie, Turquie, Maroc, Allemagne, Canada ; Instituts culturels (Goethe, Institut français, Centre culturel chinois, américain).
\end{itemize}
\end{quote}
\emph{Structurez votre légende ainsi :}
\begin{enumerate}
    \item \textbf{Signes cartographiques} (Points, Flèches, Hachures, Couleurs) — Précisez le code visuel.
    \item \textbf{Organisation de la légende} (3 rubriques titrées, items hiérarchisés).
    \item \textbf{Commentaire synthétique} (3 lignes) à insérer en bas du croquis.
\end{enumerate}
\end{exercice}

\courssub{Je me prépare au Bac}

\begin{exercice}[Sujet d'argumentation type Bac : Bakassi et la diplomatie camerounaise]
\emph{Durée conseillée : 45 min. Réponse attendue : environ 20-25 lignes, structurée, argumentée, illustrée.}
\begin{quote}
\textbf{Sujet :} « Le règlement pacifique du conflit de la \cle{péninsule de Bakassi} illustre les principes de la diplomatie camerounaise. »
\end{quote}
\emph{Consignes :}
\begin{enumerate}
    \item Analysez le sujet (définition des termes, problématique, annonce du plan).
    \item Développez en \textbf{deux parties équilibrées} :
    \begin{itemize}
        \item Partie 1 : En quoi l'affaire Bakassi (procédure CIJ 1994-2002, Accord de Greentree 2006, retrait 2008) met en œuvre le \cle{règlement pacifique des différends}, le \cle{respect du droit international} et la \cle{primauté du dialogue}.
        \item Partie 2 : Montrez que ce cas révèle aussi les \textbf{limites et tensions} (coût humain/social pour les populations, résistance locale, gestion post-conflit, articulation avec le principe de \cle{non-ingérence} vis-à-vis des pressions extérieures).
    \end{itemize}
    \item Concluez en évaluant l'héritage de cette affaire pour la crédibilité diplomatique du Cameroun aujourd'hui.
\end{enumerate}
\end{exercice}

\begin{exercice}[Étude de dossier : Crise humanitaire et action diplomatique dans le Nord-Ouest/Sud-Ouest (NOSO)]
\emph{Documents fournis :}
\begin{quote}
\textbf{Document 1 (Extrait discours Chef de l'État, 2019) :} « Le Cameroun reste attaché au respect de son unité nationale, de son intégrité territoriale et de sa souveraineté. Toute initiative de dialogue se fera dans le respect de la Constitution et des lois de la République. Le gouvernement a mis en place le Comité national de désarmement, de démobilisation et de réintégration (CNDDR) et le Plan d'urgence humanitaire pour le NOSO. »

\textbf{Document 2 (Rapport OCHA 2023, synthèse) :} « Dans les régions du Nord-Ouest et du Sud-Ouest, 1,7 million de personnes ont besoin d'assistance humanitaire. 600 000 déplacés internes. Accès humanitaire entravé par l'insécurité, les restrictions administratives et les attaques contre le personnel humanitaire (principes d'humanité, neutralité, impartialité, indépendance mis à l'épreuve). Financement du plan de réponse humanitaire (HRP) à 45\% seulement. »

\textbf{Document 3 (Communiqué MINREX, 2022) :} « Le Cameroun rejette toute ingérence dans ses affaires intérieures sous couvert d'humanitaire. La coordination de l'aide relève de la responsabilité de l'État hôte. Le Cameroun coopère avec les agences onusiennes (OCHA, HCR, PAM, UNICEF) sur la base du consentement souverain. »
\end{quote}
\emph{Questions :}
\begin{enumerate}
    \item \textbf{Identifiez} (Doc 1 et 3) les \textbf{trois principes diplomatiques} invoqués par l'État camerounais face à la crise du NOSO.
    \item \textbf{Caractérisez} (Doc 2) la \cle{crise humanitaire} dans le NOSO (nature, ampleur, populations cibles).
    \item \textbf{Analysez} la tension entre le principe de \cle{souveraineté/non-ingérence} (Doc 1, 3) et les \textbf{principes humanitaires} (Doc 2 : humanité, neutralité, accès). Comment le Cameroun tente-t-il de les concilier (CNDDR, Plan d'urgence, coordination OCHA) ?
    \item \textbf{Proposez} deux mesures concrètes pour améliorer l'\cle{efficacité de l'assistance humanitaire} tout en respectant le cadre juridique national.
\end{enumerate}
\end{exercice}

\begin{exercice}[Mise en situation professionnelle : Mémo diplomatique]
\emph{Vous êtes Conseiller à la Direction de la Coopération Multilatérale au MINREX. Le Ministre vous demande une \textbf{note de synthèse (2 pages max)} pour préparer la participation du Cameroun au \textbf{Forum sur la Coopération Sino-Africaine (FOCAC) 2024}.}
\emph{Votre note doit comporter les rubriques suivantes, renseignées avec des faits précis :}
\begin{enumerate}
    \item \textbf{Contexte et enjeux :} Rappel des relations Cameroun-Chine (historique, volume échanges, grands projets : Kribi, barrages, routes, stade Olembé, hôpital de Gyobassi, Université de Yaoundé II).
    \item \textbf{Principes directeurs :} Comment les \cle{principes de la diplomatie camerounaise} (égalité souveraine, avantage mutuel, non-ingérence, win-win) cadrent ce partenariat ?
    \item \textbf{Formes et aspects de la coopération visés :} Identifiez 3 priorités sectorielles (ex : transformation locale matières premières, industrialisation, formation professionnelle/emploi jeunes, santé, numérique, sécurité) et précisez pour chacune : \cle{Forme} (bilatérale/multilatérale FOCAC) et \cle{Aspect} (économique, technique, social...).
    \item \textbf{Risques et veille :} Endettement, transfert de technologie, contenu local, impact environnemental/social. Proposez 2 clauses de sauvegarde à négocier.
    \item \textbf{Position de négociation :} Formulez 3 \cle{revendications prioritaires} argumentées pour le Cameroun (ex : rééchelonnement dette, usine de transformation bois/cacao/aluminium, bourses formation professionnelle, maintenance infrastructures).
\end{enumerate}
\end{exercice}

\newpage

\coursec{Corrigés des exercices}

\coursec{Les principes de la diplomatie camerounaise}

\courssub{Je vérifie mes connaissances}

\begin{corrige}[Exercice 1 — QCM : Notions fondamentales]
\begin{enumerate}
    \item \textbf{B.} Les relations internationales englobent l'ensemble des interactions (politiques, économiques, culturelles, militaires, juridiques) entre une pluralité d'acteurs : États (acteurs principaux), organisations internationales, ONG, entreprises multinationales, groupes de pression, etc. L'option A relève du commerce privé, C est trop restrictif, D concerne la politique intérieure.
    \item \textbf{B.} La coopération \cle{bilatérale} est un rapport direct entre \textbf{deux États souverains} (ex : accords Cameroun-France, Cameroun-Chine). L'option A et D décrivent la coopération multilatérale. L'option C est fausse : la bilatérale repose souvent sur des traités écrits.
    \item \textbf{B.} L'\cle{assistance humanitaire} vise à \textbf{sauver des vies et alléger des souffrances} dans l'urgence (catastrophe, conflit), selon les principes d'humanité, neutralité, impartialité, indépendance. Elle est de court terme. L'aide au développement est structurelle, long terme. A est faux (non contraignante juridiquement pour les donateurs), C faux (financement public dominant), D faux (durée limitée).
    \item \textbf{C.} Le principe de \cle{non-ingérence} (Art. 2.7 Charte ONU, Préambule Constitution Cameroun) interdit d'intervenir dans les \textbf{affaires intérieures} relevant de la \cle{souveraineté} nationale d'un autre État (choix politique, électoral, législatif). A, B, D sont des contre-sens : le Cameroun vote à l'ONU, accepte l'aide, entretient des relations diplomatiques.
\end{enumerate}
\end{corrige}

\begin{corrige}[Exercice 2 — Vrai / Faux justifié]
\begin{enumerate}
    \item \textbf{Faux.} Le \cle{non-alignement} (Mouvement des Non-Alignés, 1961) signifie le refus de s'aligner sur un bloc militaire (Est/Ouest pendant la Guerre froide), \textbf{pas} le refus des organisations internationales. Le Cameroun est membre fondateur de l'OUA/UA, de la Francophonie, du Commonwealth, de l'ONU, de l'OIC, etc.
    \item \textbf{Vrai.} La \cle{coopération multilatérale} met en relation \textbf{plus de deux États}, généralement au sein d'une \textbf{organisation internationale} (ONU, UA, CEMAC, OMS, OMC, FOCAC, etc.) ou de conférences internationales.
    \item \textbf{Vrai.} La \cle{Convention de Genève (1951)} définit le réfugié comme une personne craignant avec raison d'être persécutée \textbf{hors de son pays de nationalité} (donc franchissement d'une frontière). Le \cle{déplacé interne} (DP/IDP) a fui son domicile mais \textbf{n'a pas franchi de frontière internationale} (ex : déplacés NOSO, Extrême-Nord).
    \item \textbf{Vrai.} Le différend de la \cle{péninsule de Bakassi} (1994-2008) a été réglé par la \textbf{Cour internationale de Justice (CIJ)} (arrêt 2002), puis l'\textbf{Accord de Greentree (2006)} sous égide de l'ONU, et le retrait effectif (2008). Cela illustre le \textbf{respect du droit international}, le \textbf{règlement pacifique} (Art. 33 Charte ONU) et la \cle{primauté du dialogue}.
\end{enumerate}
\end{corrige}

\begin{corrige}[Exercice 3 — Définitions et distinctions]
\begin{enumerate}
    \item \textbf{Relations internationales (RI) :} Champ d'étude et ensemble des \textbf{interactions réelles} (coopération, conflit, échanges) entre acteurs étatiques et non étatiques à l'échelle mondiale. \textbf{Politique étrangère (PE) :} \textbf{Stratégie délibérée} d'un État pour défendre ses intérêts et valeurs à l'extérieur. \textbf{Différence majeure :} Les RI sont le \emph{théâtre} global (objectif, descriptif) ; la PE est l'\emph{action volontariste} d'un acteur spécifique (subjectif, prescriptif).
    \item \textbf{Coopération bilatérale :} Accord \textbf{deux États} (ex : Commission mixte Cameroun-Nigeria). \textbf{Coopération multilatérale :} Accord \textbf{plus de deux États}, souvent via une \textbf{OI} (ex : ZLECAf, FOCAC, projets Banque Mondiale). \textbf{Différence majeure :} Nombre de parties contractantes et cadre institutionnel (direct vs institutionnalisé).
    \item \textbf{Aide humanitaire d'urgence :} Action \textbf{court terme} (semaines/mois) pour \textbf{sauver des vies} (nourriture, eau, abri, santé, protection) en situation de crise (conflit, catastrophe). Principes : humanité, neutralité, impartialité, indépendance. \textbf{Aide au développement :} Action \textbf{long terme} (années/décennies) pour \textbf{réduire la pauvreté structurelle} et renforcer les capacités (infrastructures, éducation, gouvernance, économie). \textbf{Différence majeure :} \textbf{Temporalité} (urgence vs structurel) et \textbf{finalité} (survie vs autonomie).
    \item \textbf{Diplomatie préventive :} Action \textbf{en amont} pour \textbf{empêcher l'apparition ou l'escalade} de conflits (alerte précoce, médiation, déploiement d'observateurs, confiance mutuelle). \textbf{Diplomatie coercitive :} Usage de \textbf{menaces ou pressions} (sanctions économiques, embargo, ultimatum, usage limité de la force) pour \textbf{contraindre} un acteur à changer de comportement. \textbf{Différence majeure :} \textbf{Consentement et confiance} (préventive) vs \textbf{Contrainte et rapport de force} (coercitive).
\end{enumerate}
\end{corrige}

\begin{corrige}[Exercice 4 — Classement rapide : Aspects et Formes]
\begin{center}
\begin{tabularx}{\linewidth}{|l|X|X|}\hline
\textbf{Exemple de coopération} & \textbf{Aspect} & \textbf{Forme} \\ \hline
1. Construction du port en eau profonde de Kribi (avec Chine) & Économique / Technique & Bilatérale \\ \hline
2. Participation des casques bleus camerounais à la MINUSCA (RCA) & Militaire / Sécuritaire & Multilatérale (ONU) \\ \hline
3. Attribution de bourses d'études par la France (programme Eiffel) & Culturel / Éducatif & Bilatérale \\ \hline
4. Intégration du Cameroun à la ZLECAf & Économique / Commercial & Multilatérale (UA/Continentale) \\ \hline
5. Gestion du camp de réfugiés de Minawao (HCR, ONU, ONG, État) & Humanitaire & Multilatérale (multi-acteurs) \\ \hline
6. Accord de défense signé avec les États-Unis (formation, matériel) & Militaire / Sécuritaire & Bilatérale \\ \hline
7. Projet de recherche agricole avec l'IITA (CGIAR) & Technique / Scientifique & Multilatérale (Centre de recherche international) \\ \hline
8. Festival culturel « FESTAC » organisé avec le Nigeria & Culturel & Bilatérale \\ \hline
9. Appui budgétaire de l'Union Européenne (FED) & Économique / Financier & Multilatérale (Organisation régionale d'intégration) \\ \hline
10. Mise en place de la FMM contre Boko Haram (LCBC) & Militaire / Sécuritaire & Multilatérale (Organisation sous-régionale LCBC) \\ \hline
\end{tabularx}
\end{center}
\end{corrige}

\courssub{Je m'entraîne}

\begin{corrige}[Exercice 5 — Question à réponse construite : Synthèse des notions]
\textbf{Introduction}
Les \cle{relations internationales} désignent l'ensemble des interactions — politiques, économiques, juridiques, culturelles, militaires — qui s'établissent entre une pluralité d'acteurs (États, organisations internationales, ONG, entreprises, individus) à l'échelle planétaire. Elles constituent le cadre global dans lequel s'insère l'action extérieure de chaque État.

\textbf{Développement — Deux aspects de la coopération internationale (exemples camerounais)}
Premièrement, l'\textbf{aspect économique et commercial} est prépondérant. Il vise l'intégration aux marchés, l'attraction d'investissements et le développement d'infrastructures. Le Cameroun l'illustre par son adhésion à la \cle{ZLECAf} (marché unique africain), la construction du port de Kribi avec la Chine (IDE, infrastructures), ou les accords de partenariat économique (APE) avec l'Union Européenne.
Deuxièmement, l'\textbf{aspect militaire et sécuritaire} répond aux menaces transnationales. Le Cameroun y participe activement : contribution de contingents aux opérations de maintien de la paix de l'ONU (MINUSCA en RCA, ex-MINUSMA au Mali) et de l'UA ; coopération bilatérale de défense avec la France (formation EMIA, appui renseignement) et les USA (formation, drones) ; force multinationale mixte (FMM) de la Commission du Bassin du Lac Tchad (CBLT) contre Boko Haram.

\textbf{Développement — Deux formes de la coopération internationale (exemples camerounais)}
La \textbf{coopération bilatérale} lie directement deux États souverains. Exemples : l'accord de défense Cameroun-France (1974, renouvelé), les conventions de financement chinois pour les barrages (Mekin, Memve'ele, Nachtigal) et routes, le programme de bourses « Eiffel » français, la commission mixte Cameroun-Nigeria pour la frontière.
La \textbf{coopération multilatérale} engage plusieurs États au sein d'institutions. Exemples : l'appui budgétaire sectoriel de l'UE (FED) ; les programmes de la Banque Mondiale et du FMI (facilités élargies de crédit) ; l'action humanitaire coordonnée par OCHA/HCR/PAM à Minawao ; la participation au FOCAC (Forum sur la Coopération Sino-Africaine) et au sommet Turquie-Afrique.

\textbf{Conclusion}
La diplomatie camerounaise combine habilement ces aspects et formes pour diversifier ses partenariats (Nord-Sud, Sud-Sud), assurer sa sécurité, attirer des ressources pour son développement (SND30) et affirmer son rôle de puissance régionale stabilisatrice en Afrique centrale.
\end{corrige}

\begin{corrige}[Exercice 6 — Étude de document : L'assistance humanitaire à Minawao]
\begin{enumerate}
    \item \textbf{Définition opérationnelle :} L'assistance humanitaire est l'action coordonnée de \textbf{secours d'urgence} et de \textbf{protection} mise en œuvre par une pluralité d'acteurs (État hôte, agences onusiennes, ONG) au profit de populations en détresse (réfugiés, déplacés), visant à couvrir les \textbf{besoins vitaux immédiats} (eau, nourriture, abri, santé, éducation, protection) dans le respect des principes d'humanité, neutralité, impartialité et indépendance.
    \item \textbf{Trois catégories d'acteurs et rôles :}
    \begin{itemize}
        \item \textbf{L'État camerounais (MINAS, MINSANTE, MINAT, forces de défense) :} \cle{Responsable premier} (souveraineté), sécurité du camp et des abords, enregistrement (avec HCR), accès aux services publics de santé/éducation, coordination nationale, gestion foncière.
        \item \textbf{Les agences du système des Nations Unies (HCR chef de file, PAM, UNICEF, OMS, OCHA) :} \cle{Mandat de protection et d'assistance} (HCR), \cle{logistique et financement} (PAM : vivres ; UNICEF : éducation, eau, protection de l'enfance ; OMS : santé), \cle{coordination humanitaire} (OCHA : clusters, plan de réponse HRP).
        \item \textbf{Les ONG internationales et nationales (Plan International, IMC, CARE, etc.) :} \cle{Mise en œuvre opérationnelle de terrain} (distribution, construction latrines/écoles, gestion de cas VBG, appui psychosocial, activités génératrices de revenus AGR), \cle{proximité avec bénéficiaires}, \cle{plaidoyer}.
    \end{itemize}
    \item \textbf{Deux défis majeurs et pistes durables :}
    \begin{enumerate}
        \item \textbf{Dépendance prolongée à l'aide alimentaire (PAM) / insuffisance des AGR :} \emph{Piste :} Développement de \textbf{programmes de résilience et d'autonomisation économique} (accès à la terre pour agriculture/maraîchage, formation professionnelle certifiante, appui à l'entrepreneuriat jeune/femmes, liaison avec marchés locaux) pour passer de l'assistance à l'inclusion économique.
        \item \textbf{Tensions foncières et sur ressources (eau, bois) avec populations d'accueil :} \emph{Piste :} \cle{Approche « Nexus humanitaire-développement-paix »} : programmation \textbf{intégrée} bénéficiant aux deux communautés (forages partagés, écoles/santé communes, comités de gestion locale des ressources, dialogue social, appui aux services publics locaux) pour renforcer la \cle{cohésion sociale} et prévenir les conflits.
    \end{enumerate}
\end{enumerate}
\end{corrige}

\begin{corrige}[Exercice 7 — Application : Principes directeurs de la diplomatie camerounaise]
\begin{enumerate}
    \item \textbf{Cinq principes cardinaux (Préambule Constitution, Charte ONU, pratique) :}
    \begin{enumerate}
        \item Respect de la \cle{souveraineté} et de l'\cle{intégrité territoriale} (égalité souveraine des États).
        \item \cle{Non-ingérence} dans les affaires intérieures.
        \item \cle{Règlement pacifique des différends} (négociation, médiation, arbitrage, justice internationale).
        \item \cle{Non-alignement} (liberté d'appréciation, refus des blocs, diplomatie tous azimuts).
        \item \cle{Coopération solidaire} et \cle{avantage mutuel} (win-win, Sud-Sud, Nord-Sud, intégration régionale).
    \end{enumerate}
    \item \textbf{Illustrations concrètes (depuis 1960) :}
    \begin{enumerate}
        \item \textbf{Souveraineté/Intégrité :} Affaire \cle{Bakassi} (saisine CIJ 1994, exécution arrêt 2002, Accord Greentree 2008) ; défense frontières (opérations BIR contre Boko Haram, sécurisation frontière Est/RCA).
        \item \textbf{Non-ingérence :} Refus des conditionnalités politiques liées à l'aide (années 90) ; position sur crise NOSO (rejet ingérence sous couvert droits de l'homme/humanitaire, Communiqués MINREX) ; non-reconnaissance unilatérale d'entités sécessionnistes.
        \item \textbf{Règlement pacifique :} \cle{Bakassi} (CIJ) ; \cle{Conflit frontalier Lac Tchad} (CBLT, Commission mixte Nigeria-Cameroun) ; \cle{Médiation} régionale (RCA, Tchad, Guinée équatoriale) ; respect décisions justice internationale (CIJ, CCJA OHADA).
        \item \textbf{Non-alignement :} Adhésion simultanée \cle{Francophonie} (1970) et \cle{Commonwealth} (1995) ; membre \cle{UA}, \cle{ONU}, \cle{OCI}, \cle{G77+Chine}, \cle{Mouvement Non-Alignés} ; relations équilibrées France/USA/Chine/Russie/Turquie/Brésil/Inde/Israël/Maroc.
        \item \textbf{Coopération solidaire :} \cle{Casques bleus} (top contributeur africain ONU) ; \cle{Accueil réfugiés} (Nigeria, RCA, Tchad — loi 2005, camps Minawao, Gado, Borgop) ; \cle{Partenariats Sud-Sud} (FOCAC, Turquie-Afrique, Corée-Afrique, Brésil, Inde : transferts technologie, formation, infrastructures) ; \cle{Intégration CEMAC/CEEAC/ZLECAf}.
    \end{enumerate}
    \item \textbf{Non-alignement $\neq$ Neutralité $\neq$ Isolement :}
    La \cle{neutralité} est un statut juridique (ex : Suisse, Autriche) impliquant l'abstention dans tout conflit armé et la non-appartenance à des alliances militaires. L'\cle{isolement} est une rupture volontaire des relations.
    Le \cle{non-alignement} camerounais est une \textbf{liberté d'action} : le pays \textbf{choisit ses partenaires} au cas par cas selon ses intérêts, sans allégeance automatique. Preuve : membership actif dans la \cle{Francophonie} (liens linguistiques, culturels, institutionnels avec France/Québec/Afrique), le \cle{Commonwealth} (liens anglophones, commerce, éducation, gouvernance avec UK/Canada/Inde/Nigeria/Afrique du Sud), l'\cle{UA} (intégration continentale, paix/sécurité), l'\cle{OCI} (solidarité islamique), le \cle{G77+Chine} (négociation Nord-Sud). Cette multiplicité d'ancrages \textbf{renforce} son influence et sa marge de manœuvre, loin de tout isolement.
\end{enumerate}
\end{corrige}

\begin{corrige}[Exercice 8 — Travail pratique noté : Croquis « Le Cameroun et ses partenaires dans le monde »]
\begin{enumerate}
    \item \textbf{Signes cartographiques (Code visuel)}
    \begin{itemize}
        \item \textbf{Flèches de flux commerciaux (Export/Import) :} Largeur \textbf{proportionnelle au volume} (échelle : 1 mm = 5\% ou 10\% du total). \textbf{Couleur :} \textcolor{popblue}{Bleu} = Exportations camerounaises (sortantes) ; \textcolor{poporange}{Orange} = Importations (entrantes). Pointe de flèche vers le partenaire.
        \item \textbf{Flèches d'Investissements (IDE stock) :} \textbf{Trait discontinu épais} (\textcolor{popgreen}{Vert}) du pays investisseur vers le Cameroun (ou symbole \texteuro/\$ sur le Cameroun avec origine en légende). Taille proportionnelle au stock.
        \item \textbf{Flèches de Coopération militaire/sécuritaire :} \textbf{Trait pointillé} (\textcolor{poppurple}{Violet}) bidirectionnel ou vers le Cameroun (formation, équipement). Symbole ou casque.
        \item \textbf{Flèches Missions de paix (Casques bleus) :} \textbf{Flèche fine} (\textcolor{popblue}{Bleu ONU}) du Cameroun vers théâtres d'opérations (RCA, Mali, Soudan, Haïti). Symbole .
        \item \textbf{Flèches Migrations/Réfugiés :} \textbf{Flèches courbes} (\textcolor{popred}{Rouge} pour entrants réfugiés/IDP, \textcolor{poppink}{Rose} pour sortants réfugiés camerounais). Largeur proportionnelle aux effectifs (échelle : 1 mm = 50 000 pers.). Hachures \textcolor{popred!30}{rouge clair} sur zones d'origine (Extrême-Nord, Nord-Ouest, Sud-Ouest, Est, Adamaoua).
        \item \textbf{Points / Pictogrammes :} (bourses/coopération universitaire), (instituts culturels), (projets industriels/IDE), (capitales partenaires).
        \item \textbf{Fond de carte :} Pays partenaires colorés par \textbf{groupe d'affinité} (UE \textcolor{popblueL}{bleu clair}, Chine \textcolor{popredL}{rouge clair}, USA \textcolor{popblue}{bleu foncé}, Afrique \textcolor{popgreenL}{vert clair}, Russie \textcolor{poporangeL}{orange clair}, Turquie/Maroc/Monde arabe \textcolor{poppurpleL}{violet clair}).
    \end{itemize}
    \item \textbf{Organisation de la légende (3 rubriques titrées, items hiérarchisés)}
    \begin{center}
    \textbf{LÉGENDE : Le Cameroun et ses partenaires dans le monde}
    \end{center}
    \textbf{I. FLUX ÉCONOMIQUES (Commerce \& Investissements)}
    \begin{itemize}
        \item \textcolor{popblue}{$\blacktriangleright$} \textbf{Exportations (2022) :} Largeur prop. au \%. UE (45\%), Chine (15\%), Inde (10\%), USA (5\%), Afrique CEMAC/CEEAC (10\%).
        \item \textcolor{poporange}{$\blacktriangleright$} \textbf{Importations (2022) :} Largeur prop. au \%. Chine (25\%), UE (20\%), Nigeria (15\%), Inde (5\%).
        \item \textcolor{popgreen}{$\blacktriangledown$} \textbf{IDE Stock (Principaux pays) :} France (historique), Chine (Infra : Kribi, Barrages, Routes), USA (Pétrole/Gaz), Maroc (Banque/Telecoms). Symbole sur Cameroun + origine.
    \end{itemize}
    \textbf{II. COOPÉRATIONS POLITIQUES / SÉCURITAIRES}
    \begin{itemize}
        \item \textcolor{poppurple}{$\cdots\blacktriangleright$} \textbf{Accords de défense / Formation / Équipement :} France (EMIA), USA (Renseignement/Drones), Russie (Équipements), Israël (Modernisation), Chine (Fournitures).
        \item \textcolor{popblue}{$\blacktriangleright$} \textbf{Missions de Paix (ONU/UA) :} Cameroun $\to$ RCA (MINUSCA), Mali (ex-MINUSMA), Soudan, Haïti (Police). Effectifs $\sim$ 1000-2000 pers.
        \item \textbf{Diplomatie multilatérale :} Sièges UA (Addis), ONU (NY/Geneva), OIC, Francophonie, Commonwealth, CEMAC, CBLT, ZLECAf.
    \end{itemize}
    \textbf{III. MIGRATIONS / RÉFUGIÉS / DÉPLACÉS (HCR 2023)}
    \begin{itemize}
        \item \textcolor{popred}{$\blacktriangleright$} \textbf{Réfugiés entrants :} Nigeria $\to$ Extrême-Nord ($\sim$120k), RCA $\to$ Est/Adamaoua ($\sim$350k). Largeur prop.
        \item \textcolor{popred!50}{$\blacktriangleright$} \textbf{Déplacés Internes (IDP) :} NOSO ($\sim$600k), Extrême-Nord ($\sim$300k). Hachures zones concernées.
        \item \textcolor{poppink}{$\blacktriangleright$} \textbf{Réfugiés camerounais sortants :} Vers Nigeria, Tchad (effectifs moindres).
    \end{itemize}
    \item \textbf{Commentaire synthétique (3 lignes, bas du croquis)}
    \begin{quote}
    « Puissance régionale d'équilibre, le Cameroun diversifie ses partenariats stratégiques (Nord-Sud, Sud-Sud) pour son développement (SND30) et sa sécurité. Premier contributeur africain aux OMP, il assume un leadership stabilisateur tout en gérant une crise humanitaire interne (NOSO, Extrême-Nord) et externe (refugiés RCA/Nigeria), illustrant la complexité de sa diplomatie tous azimuts. »
    \end{quote}
\end{enumerate}
\end{corrige}

\courssub{Je me prépare au Bac}

\begin{corrige}[Exercice 9 — Sujet d'argumentation type Bac : Bakassi et la diplomatie camerounaise]
\textbf{Barème indicatif (sur 20) :} Analyse sujet (3) + Partie 1 (6) + Partie 2 (6) + Conclusion (3) + Qualité rédaction/structure (2).

\textbf{Analyse du sujet}
\textbf{Termes clés :} « Règlement pacifique » (Art. 33 Charte ONU : négociation, enquête, médiation, conciliation, arbitrage, règlement judiciaire) ; « Conflit de la péninsule de Bakassi » (différend frontalier Cameroun-Nigeria, zone pétrolifère/pêche, population nigériane majoritaire) ; « Principes diplomatie camerounaise » (souveraineté, droit international, non-violence, dialogue, non-ingérence).
\textbf{Problématique :} En quoi le règlement de l'affaire Bakassi, par la voie judiciaire (CIJ) puis diplomatique (Greentree), constitue-t-il l'archétype de l'application des principes constitutionnels camerounais, tout en révélant les tensions inhérentes à leur mise en œuvre concrète sur les populations ?
\textbf{Annonce du plan :} Nous montrerons d'abord que l'affaire Bakassi est une application exemplaire des principes de règlement pacifique, de primauté du droit et du dialogue (I), avant d'analyser les limites humaines, sociales et politiques qui en tempèrent le succès diplomatique (II).

\textbf{Partie 1 : Une illustration textbook des principes camerounais (Règlement pacifique, Droit international, Dialogue)}
Le Cameroun, fidèle au Préambule de sa Constitution (« règlement pacifique des différends », « respect du droit international »), saisit la \textbf{CIJ en 1994} plutôt que la force, malgré l'occupation militaire nigériane. C'est le choix de la \cle{voie juridique} sur la \cle{voie des armes}. L'\textbf{arrêt du 10 octobre 2002} confirme la souveraineté camerounaise sur Bakassi (traité 1913, effet d'uti possidetis juris). Le Cameroun \textbf{exige l'exécution intégrale}, refusant tout partage politique.
La \cle{primauté du dialogue} s'incarne dans l'\textbf{Accord de Greentree (12 juin 2006)}, signé sous l'égide du Secrétaire général de l'ONU (Kofi Annan), avec témoins (France, UK, Allemagne, USA). Cet accord organise le \textbf{transfert progressif, pacifique et ordonné} de l'administration et des forces de sécurité (retrait nigérian août 2006, retrait total août 2008), la protection des populations (droit de rester, nationalité, biens), et la délimitation terrestre et maritime. Le Cameroun fait preuve de \textbf{retenue stratégique}, de \textbf{patience diplomatique} et de \textbf{confiance dans le multilatéralisme} (ONU comme garant). C'est une victoire du \cle{droit} sur la \cle{force}.

\textbf{Partie 2 : Limites et tensions (Coût humain, Résistance locale, Post-conflit, Non-ingérence vs Pressions)}
Ce « succès » diplomatique masque des \textbf{tensions majeures}.
\textbf{1. Coût humain et social :} La population (majoritairement nigériane, pêcheurs Efik/Ibibio) se sent \textbf{abandonnée} par le Nigeria et \textbf{méfiante} envers l'administration camerounaise. Perte de repères, traumatismes, exactions ponctuelles (rapports ONG), difficulté d'accès à la nationalité camerounaise (délais, corruption), perte d'accès aux zones de pêche traditionnelles.
\textbf{2. Résistance locale et insécurité résiduelle :} Apparition de groupes armés (ex : \textit{Bakassi Freedom Fighters}, \textit{SCNC} radicaux) revendiquant l'autodétermination ou le rattachement au Nigeria, menant des attaques contre l'armée et l'administration camerounaises (ex : attaque péninsule 2010, enlèvements). Gestion sécuritaire lourde (militarisation) vs développement civil.
\textbf{3. Gestion post-conflit perfectible :} Promesses de développement (routes, électricité, hôpital, port artisanal) \textbf{peu tenues} ou lentes. Sentiment de \cle{territoire conquis mais pas intégré}. Défis de l'état civil, de l'éducation, de la santé.
\textbf{4. Non-ingérence vs Pressions extérieures :} Le Cameroun a dû \textbf{résister aux pressions} pour un « partage politique » ou un statut spécial (autonomie) réclamé par certains acteurs locaux/internationaux, au nom du principe d'\cle{intégrité territoriale} et de \cle{non-ingérence}. L'Accord de Greentree, bien que négocié, a été perçu par une frange de l'opinion comme une \cle{ingérence} de la « communauté internationale » (ONU, puissances occidentales) dans le règlement d'un différend bilatéral.

\textbf{Conclusion}
L'affaire Bakassi \textbf{consacre la crédibilité juridique et diplomatique} du Cameroun : un État qui gagne par le droit, respecte sa signature, et assume ses obligations internationales. Elle renforce son \textbf{leadership moral} en Afrique (modèle de règlement pacifique). Cependant, l'\textbf{héritage reste contrasté} : la \cle{souveraineté} formelle est acquise, mais la \cle{souveraineté effective} (administration, développement, adhésion des populations) reste un chantier ouvert. La crédibilité future du Cameroun dépendra de sa capacité à transformer cette « victoire de papier » en \cle{paix durable} et \cle{développement partagé} à Bakassi, prouvant que le droit international sert aussi les peuples.
\end{corrige}

\begin{corrige}[Exercice 10 — Étude de dossier : Crise humanitaire et action diplomatique dans le NOSO]
\begin{enumerate}
    \item \textbf{Trois principes diplomatiques invoqués (Doc 1 \& 3) :}
    \begin{enumerate}
        \item \cle{Respect de la souveraineté nationale} et de l'\cle{intégrité territoriale} (Doc 1 : « unité nationale, intégrité territoriale, souveraineté » ; Doc 3 : « rejette toute ingérence dans ses affaires intérieures »).
        \item \cle{Non-ingérence} dans les affaires intérieures (Doc 3 : « rejette toute ingérence sous couvert d'humanitaire » ; « coordination relève de la responsabilité de l'État hôte »).
        \item \cle{Primauté du cadre constitutionnel et légal national} pour le dialogue et la résolution (Doc 1 : « dans le respect de la Constitution et des lois de la République » ; mise en place CNDDR, Plan d'urgence par décret/loi).
    \end{enumerate}
    \item \textbf{Caractérisation de la crise humanitaire NOSO (Doc 2) :}
    \begin{itemize}
        \item \textbf{Nature :} Crise de \textbf{protection} complexe (conflit armé non international, violences contre civils, VBG, recrutement enfants, destruction infrastructures) + crise de \textbf{déplacement massif}.
        \item \textbf{Ampleur :} \textbf{1,7 million de personnes} en besoin d'assistance (sur ~4-5M d'habitants des 2 régions). \textbf{600 000 déplacés internes} (IDP). Taux de financement HRP \textbf{45\% seulement} (sous-financement critique).
        \item \textbf{Populations cibles :} IDP (familles entières, femmes, enfants), communautés d'accueil surchargées, populations restées sur place (enclaves, zones difficiles d'accès), personnel humanitaire ciblé.
        \item \textbf{Contraintes d'accès :} Insécurité (IED, embuscades), restrictions administratives (autorisations, checkpoints), attaques directes contre humanitaires (violation principes humanitaires).
    \end{itemize}
    \item \textbf{Analyse de la tension Souveraineté/Non-ingérence vs Principes humanitaires :}
    \begin{itemize}
        \item \textbf{Tension :} L'État (Doc 1, 3) affirme son \textbf{monopole de la coordination} et son \textbf{droit de regard/contrôle} (souveraineté, non-ingérence) : l'aide doit passer par ses structures (CNDDR, Plan d'urgence, MINAS), avec son accord. Les principes humanitaires (Doc 2 : \cle{humanité, neutralité, impartialité, indépendance}) exigent un \textbf{accès direct, sans entrave, aux populations} selon les seuls besoins, sans discrimination, et une autonomie de décision opérationnelle vis-à-vis des parties au conflit (État comme groupes armés).
        \item \textbf{Risques :} Instrumentalisation de l'aide (conditionnement accès zones, listes bénéficiaires), confusion civilo-militaire (convois escortés par armée = perte neutralité), exclusion de zones « non contrôlées » (violation impartialité), criminalisation de l'action humanitaire (arrestations staff).
        \item \textbf{Tentatives de conciliation (Cameroun) :}
        \begin{itemize}
            \item \cle{CNDDR} (Comité National DDR) : Cadre \textbf{étatique} pour le désarmement, la réinsertion socio-économique (aspect développement/paix), tentative d'offrir une sortie aux combattants.
            \item \cle{Plan d'urgence humanitaire NOSO (2018-2020+)} : \textbf{Financement 100\% national} (budget État), exécution par structures étatiques (MINAS, MINSANTE, MINEDUB) + ONG nationales. Réaffirme la \cle{maîtrise souveraine} de la réponse.
            \item \cle{Coordination avec OCHA/Clusters} : Participation aux réunions de coordination (HNO/HRP), partage d'informations, acceptation progressive d'ONG internationales dans certaines zones, sous \textbf{supervision étatique}. Mécanisme de « notification » mouvements humanitaires.
        \end{itemize}
        \textbf{Bilan :} Conciliation \textbf{inachevée et fragile}. L'accès reste entravé, le financement externe bas, la confiance faible. L'État privilégie la \cle{sécurisation d'abord} (approche militaire/DDR) avant l'humanitaire pur.
    \end{itemize}
    \item \textbf{Deux mesures concrètes pour améliorer l'efficacité (cadre juridique respecté) :}
    \begin{enumerate}
        \item \textbf{Mise en place d'un \cle{Mécanisme de déconfliction humanitaire civilo-militaire (CMCoord) formel et opérationnel :}} Accord écrit entre MINDEF/MINAT et OCHA/HCR définissant : zones/jours « fenêtres humanitaires » (cessez-le-feu localisés), procédure de notification simplifiée (48h $\to$ 24h), identification claire des convois (logos, fréquences radio), formation conjointe droit humanitaire/DIH pour militaires et humanitaires. \emph{Base légale :} DIH (Conventions Genève, Protocoles additionnels), Loi 2005 sur réfugiés, Constitution (protection civils).
        \item \textbf{\cle{Localisation accrue de l'aide (Localization Agenda) :}} Transfert direct de \textbf{financements (HRP) et de capacités} aux \textbf{ONG nationales et organisations communautaires de base (OCB)} qui ont l'accès, la confiance et la connaissance du terrain (langues, coutumes, réseaux). Appui technique (renforcement capacités gestion, reddition comptes, protection) par ONU/ONGI. \emph{Base légale :} Loi sur la liberté d'association (1990), Convention de partenariat État-ONG, Grand Bargain (engagement 2016 : 25\% financement direct aux acteurs locaux).
    \end{enumerate}
\end{enumerate}
\end{corrige}

\begin{corrige}[Exercice 11 — Mise en situation professionnelle : Mémo diplomatique FOCAC 2024]
\textbf{NOTE DE SYNTHÈSE — CONFIDENTIEL}
\textbf{À :} Monsieur le Ministre des Relations Extérieures (MINREX)
\textbf{De :} Conseiller, Direction de la Coopération Multilatérale
\textbf{Objet :} Préparation participation Cameroun au FOCAC 2024 (Pékin) — Position de négociation
\textbf{Date :} Septembre 2024

\vspace{0.5cm}
\noindent\textbf{1. CONTEXTE ET ENJEUX}
Les relations Cameroun-Chine, établies en \textbf{1971}, sont un pilier de la diplomatie « tous azimuts ». La Chine est le \textbf{1er partenaire commercial} (échanges $\sim$ 1500-2000 Mds FCFA/an), le \textbf{1er bailleur de fonds bilatéral} (prêts concessionnels, dons) et un acteur majeur des \textbf{infrastructures structurantes} (SND30) :
\begin{itemize}
    \item \textbf{Port de Kribi (Phase 1 \& 2) :} Hub régional, ZES.
    \item \textbf{Barrages :} Mekin, Memve'ele, Nachtigal (énergie industrielle).
    \item \textbf{Routes :} Yaoundé-Nsimalen, Kribi-Lolabé, sections Transcamerounaise.
    \item \textbf{Équipements sociaux :} Stade Olembé (CAN 2021), Hôpital de Gyobassi (référence sous-régionale), Université de Yaoundé II (Soa).
    \item \textbf{Numérique :} Huawei (backbone, 4G/5G, Safe City, e-gov).
\end{itemize}
\textbf{Enjeux 2024 :} Passer du « tout infrastructures » à la \textbf{transformation structurelle} (industrialisation, emploi jeunes, souveraineté alimentaire/sanitaire/numérique), gérer la \textbf{soutenabilité de la dette} (DSA FMI : risque modéré $\to$ élevé), exiger le \textbf{transfert de technologie} et le \textbf{contenu local}.

\vspace{0.3cm}
\noindent\textbf{2. PRINCIPES DIRECTEURS (Cadre diplomatie camerounaise)}
\begin{itemize}
    \item \cle{Égalité souveraine} : Partenariat d'égal à égal, non conditionné (politique/droits de l'homme), respect choix de développement.
    \item \cle{Non-ingérence} : La Chine ne s'immisce pas dans les affaires intérieures (NOSO, gouvernance) ; le Cameroun respecte la politique « Une Chine » (Taiwan, HK, Xinjiang).
    \item \cle{Avantage mutuel / Win-Win} : Infrastructures/ressources (bois, minerais, coton, pétrole) vs Marché chinois, IDE, expertise. Doit évoluer vers \cle{transformation locale} (valeur ajoutée au Cameroun).
    \item \cle{Respect du droit international} : Contractualisation transparente (marchés publics, conventions de prêt), respect normes environnementales/sociales (ESIA).
\end{itemize}

\vspace{0.3cm}
\noindent\textbf{3. FORMES ET ASPECTS DE LA COOPÉRATION VISÉS (3 Priorités sectorielles)}

\begin{center}
\begin{tabularx}{\linewidth}{|l|X|X|X|}\hline
\textbf{Priorité sectorielle} & \textbf{Forme} & \textbf{Aspect} & \textbf{Objectif opérationnel FOCAC 2024} \\ \hline
1. \textbf{Transformation locale matières premières} (Bois 2\&3 transfo, Cacao poudre/beurre, Coton fil/tissu, Aluminium/Alumine, Pétrole raffinage) & Bilatérale (Accords cadre FOCAC) + Multilatérale (FOCAC Plan d'action 2025-2027) & \textbf{Économique / Industriel / Commercial} & Négocier \textbf{usines clefs en main} (clé en main + transfert techno) dans ZES (Kribi, Edea, Ngaoundéré). Exiger quotas export brut $\downarrow$, export transformé $\uparrow$. \\ \hline
2. \textbf{Formation professionnelle / Emploi jeunes / Insertion} (Écoles métiers BTP, agro-industrie, numérique, maintenance industrielle, hôtellerie) & Bilatérale (Gouv-Gouv, Minfopra-Ministère Commerce Chine) + Multilatérale (FOCAC « Programme 1000 villages / Luban Workshop ») & \textbf{Social / Technique / Éducatif} & Créer \textbf{5-10 Ateliers Luban} (formation duale entreprises chinoises au Cameroun). \textbf{10 000 bourses/an} formation courte/certifiante (pas seulement académique). Intégration jeunes dans projets chinois (contenu local 30-50\%). \\ \hline
3. \textbf{Santé (Production vaccins/médicaments, Télémédecine, Hôpitaux de référence) \& Numérique (Data centers, Cybersécurité, E-gov, IA)} & Bilatérale (Ministères Santé/Posts \& Télécoms) + Multilatérale (FOCAC Santé / Numérique) & \textbf{Social / Technique / Scientifique} & Usine de remplissage/conditionnement vaccins (transfer tech). Extension \textbf{Safe City / Smart City} (Yaoundé, Douala, villes secondaires). Souveraineté données (Data center national). \\ \hline
\end{tabularx}
\end{center}

\vspace{0.3cm}
\noindent\textbf{4. RISQUES ET VEILLE — CLAUSES DE SAUVEGARDE À NÉGOCIER}
\begin{enumerate}
    \item \textbf{Risque dette / Soutenabilité :} Stock dette Chine $\sim$ 30-40\% dette ext. publique. Risque de « dette piège » (actifs stratégiques port/énergie en garantie).
    \item \textbf{Risque faible transfert technologie / Contenu local :} Main-d'œuvre/encadrement chinois majoritaires, sous-traitance locale limitée, opacité coûts.
    \item \textbf{Risque Environnemental/Social :} Déforestation (bois), pollution (mines, usines), déplacements populations, conditions travail.
\end{enumerate}
\textbf{2 Clauses de sauvegarde prioritaires :}
\begin{enumerate}
    \item \textbf{Clause « Contenu Local Obligatoire \& Transfert Technologique » :} \textbf{Minimum 40\% main-d'œuvre camerounaise} (dont 20\% qualifiés formés sur place), \textbf{sous-traitance 30\% PME locales}, \textbf{plan de transfert techno} contractualisé (brevets, savoir-faire, maintenance) avec pénalités financières si non-respect. Audit indépendant annuel.
    \item \textbf{Clause « Soutenabilité Dette \& Évaluation Ex-Ante » :} Tout nouveau prêt \textbf{conditionné à une Analyse de Soutenabilité Dette (DSA) conjointe FMI/Banque Mondiale} publiée. \textbf{Clause de rééchelonnement automatique} (maturité +5-10 ans, différé 5-7 ans, taux $\le$ 1-2\%) en cas de choc exogène (prix matières premières, pandémie, climatique). \textbf{Interdiction garanties sur actifs stratégiques souverains} (port, barrages, aéroports, fréquences).
\end{enumerate}

\vspace{0.3cm}
\noindent\textbf{5. POSITION DE NÉGOCIATION — 3 REVENDICATIONS PRIORITAIRES ARGUMENTÉES}
\begin{enumerate}
    \item \textbf{Rééchelonnement / Restructuration proactive de la dette existante (Stock $\sim$ 3000 Mds FCFA) :} \emph{Argument :} « Le Cameroun a honoré tous ses engagements (pas de défaut). La maturité courte (7-10 ans) et les taux (2-3\%) des prêts concessionnels récents pèsent sur le service de la dette (ratio service/recettes $>$ 30\%). Demande : \textbf{Allongement maturités à 20-25 ans, différé 8-10 ans, taux fixe 1\%, clause catastrophe naturelle.} Contrepartie : Accélération projets à fort retour économique (Kribi Phase 2, Aluminium, Raffinerie) + transparence contrats. »
    \item \textbf{Implantation d'unités de transformation industrielle intégrées (Bois, Cacao, Coton, Aluminium) dans les ZES :} \emph{Argument :} « Le Cameroun exporte 80\% de ses grumes (bois) et fèves de cacao brutes. La Chine est le 1er acheteur. L'industrialisation (SND30) exige la \textbf{transformation sur place}. Demande : \textbf{Financement/concession d'usines clés} (Sciage/Placage/Meubles 2\&3 transfo ; Torréfaction/Poudre/Beurre cacao ; Filature/Tissage ; Aluminerie) \textbf{clé en main avec transfert techno} dans ZES Kribi/Edea/Ngaoundéré. Accès préférentiel marché chinois pour produits finis (zéro droit ZLECAf + protocole FOCAC). »
    \item \textbf{Programme massif de formation professionnelle duale (Ateliers Luban) et maintenance infrastructures :} \emph{Argument :} « Le chômage jeunes (15-35 ans) est la 1ère menace sociale. Les infrastructures chinoises (routes, barrages, stades, hôpitaux, numérique) vieillissent sans maintenance locale. Demande : \textbf{10 Ateliers Luban opérationnels 2025} (BTP, Énergie, Numérique, Agro, Maintenance industrielle) co-gérés MINFOPRA/Entreprises chinoises. \textbf{Fonds de maintenance dédié} (1-2\% CAPEX/an) financé sur prêt/projet, géré par structure mixte. \textbf{5000 bourses/an} formation courte certifiante (pas diplômes longs seulement). »
\end{enumerate}
\vspace{0.5cm}
\noindent\textbf{Avis du Conseiller :} La délégation (MINREX, MINEPAT, MINFI, MINCOMMERCE, MININDUSTRIE, MINFOPRA, MINSANTE, MINPOSTEL) doit arriver avec \textbf{dossiers techniques ficelés} (pré-faisabilité usines, cahiers charges Luban, DSA actualisée) et une \textbf{équipe de négociation unique} (Lead MINREX) pour éviter la fragmentation. Le « Win-Win » 2024 se joue sur la \textbf{valeur ajoutée locale} et la \textbf{soutenabilité financière}, pas sur le volume de nouveaux prêts.
\end{corrige}

\newpage

\coursec{Fiche bilan}

\begin{popfigure}
\begin{center}\tcbox[colback=poporange,colframe=poporange,coltext=white,arc=9pt,boxsep=4pt,left=6pt,right=6pt]{\bfseries\small Diplomatie camerounaise}\end{center}
\vspace{-2pt}
\begin{tcbraster}[raster columns=3,raster equal height=rows,raster column skip=6pt,raster row skip=6pt,enhanced,blankest]
\begin{tcolorbox}[enhanced,colback=poporange,colframe=poporange,coltext=white,boxrule=0.9pt,arc=3pt,left=4pt,right=4pt,top=3pt,bottom=3pt,boxsep=1pt,halign=left]\bfseries\scriptsize Relations internationales (concept)\end{tcolorbox}
\begin{tcolorbox}[enhanced,colback=popgreen,colframe=popgreen,coltext=white,boxrule=0.9pt,arc=3pt,left=4pt,right=4pt,top=3pt,bottom=3pt,boxsep=1pt,halign=left]\bfseries\scriptsize Coopération : aspects (politique, éco., culturel)\end{tcolorbox}
\begin{tcolorbox}[enhanced,colback=poppurple,colframe=poppurple,coltext=white,boxrule=0.9pt,arc=3pt,left=4pt,right=4pt,top=3pt,bottom=3pt,boxsep=1pt,halign=left]\bfseries\scriptsize Coopération : formes (bilatérale, multilatérale)\end{tcolorbox}
\begin{tcolorbox}[enhanced,colback=poppink,colframe=poppink,coltext=white,boxrule=0.9pt,arc=3pt,left=4pt,right=4pt,top=3pt,bottom=3pt,boxsep=1pt,halign=left]\bfseries\scriptsize Assistance humanitaire (Dossier 1)\end{tcolorbox}
\begin{tcolorbox}[enhanced,colback=popgold,colframe=popgold,coltext=white,boxrule=0.9pt,arc=3pt,left=4pt,right=4pt,top=3pt,bottom=3pt,boxsep=1pt,halign=left]\bfseries\scriptsize Principes fondamentaux (non-ingérence, etc.)\end{tcolorbox}
\begin{tcolorbox}[enhanced,colback=popblue!70!popgreen,colframe=popblue!70!popgreen,coltext=white,boxrule=0.9pt,arc=3pt,left=4pt,right=4pt,top=3pt,bottom=3pt,boxsep=1pt,halign=left]\bfseries\scriptsize Acteurs et instruments (MINREX, ambassades)\end{tcolorbox}
\begin{tcolorbox}[enhanced,colback=popdark,colframe=popdark,coltext=white,boxrule=0.9pt,arc=3pt,left=4pt,right=4pt,top=3pt,bottom=3pt,boxsep=1pt,halign=left]\bfseries\scriptsize Enjeux actuels (paix, développement)\end{tcolorbox}
\end{tcbraster}

\legende{Carte mentale de synthèse : Leçon 01 -- Les principes de la diplomatie camerounaise. Chaque branche correspond à une partie du plan de la leçon.}
\end{popfigure}

\courssub{Définitions clés}

\begin{definition}[Les notions fondamentales de la leçon]
\begin{itemize}
  \item \cle{Relations internationales (RI)} : ensemble des interactions (politiques, économiques, juridiques, culturelles) entre les acteurs de la scène mondiale : États, organisations internationales (OI), organisations non gouvernementales (ONG) et entreprises transnationales.
  \item \cle{Coopération internationale} : action concertée entre États ou organisations pour atteindre des objectifs communs (développement, paix, santé, éducation, environnement).
  \item \cle{Diplomatie} : art de représenter un État, de négocier et de défendre ses intérêts auprès des autres États et des organisations internationales ; c'est aussi l'ensemble des agents (diplomates, ambassadeurs) et des procédures qui la mettent en œuvre.
  \item \cle{Assistance humanitaire} : aide d'urgence (vivres, soins, abris, protection) apportée aux populations en détresse (conflits, catastrophes naturelles, déplacements forcés), selon les principes d'\emph{humanité}, de \emph{neutralité}, d'\emph{impartialité} et d'\emph{indépendance}.
\end{itemize}
\end{definition}

\begin{attention}[Ne pas confondre]
La \textbf{politique étrangère} est la \emph{stratégie} définie par le chef de l'État (les objectifs, les orientations) ; la \textbf{diplomatie} est l'\emph{instrument} de mise en œuvre de cette stratégie (négociations, ambassades, traités). Au Bac, confondre les deux notions fait perdre des points de définition.
\end{attention}

\courssub{Principes de la diplomatie camerounaise (à connaître par cœur)}

\begin{center}
\begin{tabularx}{\linewidth}{|>{\bfseries}l|X|}\hline
\rowcolor{popblueL}\textbf{Principe} & \textbf{Contenu / Application} \\\hline
Non-ingérence & Respect de la souveraineté des États ; refus d'intervenir dans les affaires intérieures d'autrui.\\\hline
Bon voisinage & Recherche active de la paix, de la sécurité et de la coopération avec les pays frontaliers (Nigéria, Tchad, RCA, Congo, Gabon, Guinée équatoriale).\\\hline
Panafricanisme & Engagement pour l'unité africaine (UA, CEMAC, CEEAC) ; défense des intérêts du continent.\\\hline
Multilatéralisme & Privilégie le règlement des différends par le droit et les organisations internationales (ONU, OIF, Commonwealth, OCI).\\\hline
Coopération gagnant-gagnant & Partenariats économiques mutuellement avantageux (Sud-Sud, Nord-Sud : Chine, France, États-Unis, etc.).\\\hline
Respect du droit international & Adhésion aux chartes (ONU, UA) et aux traités ; règlement pacifique des conflits (ex. affaire de Bakassi tranchée par la CIJ en 2002, appliquée avec le Nigéria).\\\hline
Diplomatie économique & Promotion des investissements, des exportations et des transferts de technologie au service de la SND30 (Stratégie Nationale de Développement 2020-2030).\\\hline
\end{tabularx}
\end{center}

\begin{exemplebox}[Le règlement pacifique du différend de Bakassi]
Le Cameroun et le Nigéria s'opposaient sur la souveraineté de la péninsule de Bakassi. Plutôt que la guerre, le Cameroun a saisi la \textbf{Cour internationale de Justice} (CIJ), qui lui a donné raison en 2002 ; les \emph{Accords de Greentree} (2006) ont organisé le transfert pacifique. C'est l'illustration parfaite des principes de \cle{respect du droit international} et de \cle{bon voisinage}.
\end{exemplebox}

\courssub{Méthodes express}

\begin{methode}[Analyser une situation de coopération internationale]
  \begin{enumerate}
    \item Identifier les \textbf{acteurs} (États, OI, ONG, secteur privé).
    \item Qualifier la \textbf{forme} : bilatérale / multilatérale ; Nord-Sud / Sud-Sud / triangulaire.
    \item Repérer l'\textbf{aspect} : politique, économique, culturel, technique, militaire ou humanitaire.
    \item Évaluer la \textbf{plus-value} pour le Cameroun (alignement avec la SND30, transfert de compétences, financement).
  \end{enumerate}
\end{methode}

\begin{methode}[Justifier l'action diplomatique du Cameroun]
  \begin{enumerate}
    \item Citer le \textbf{principe} mobilisé (ex. non-ingérence, panafricanisme).
    \item Rappeler l'\textbf{instrument juridique} (Charte de l'ONU, Acte constitutif de l'UA, Traité de la CEMAC).
    \item Illustrer par un \textbf{fait concret} (médiation, sommet, accord commercial, règlement de Bakassi).
    \item Conclure sur la \textbf{cohérence} avec l'intérêt national et la SND30.
  \end{enumerate}
\end{methode}

\begin{aretenir}[Checklist avant le Bac]
  \begin{itemize}
    \item $\square$ Définir sans ambiguïté : relations internationales, coopération internationale, diplomatie, assistance humanitaire.
    \item $\square$ Citer et expliquer les \textbf{principes cardinaux} de la diplomatie camerounaise (non-ingérence, bon voisinage, panafricanisme, multilatéralisme, coopération gagnant-gagnant, respect du droit international, diplomatie économique).
    \item $\square$ Différencier \textbf{coopération bilatérale} et \textbf{multilatérale} avec des exemples camerounais.
    \item $\square$ Identifier les \textbf{aspects} (politique, économique, culturel, technique, militaire) d'un accord de coopération.
    \item $\square$ Résumer les \textbf{4 principes humanitaires} : humanité, neutralité, impartialité, indépendance.
    \item $\square$ Nommer les \textbf{principaux acteurs} de la diplomatie camerounaise : Président de la République, MINREX (Ministère des Relations Extérieures), ambassadeurs, consuls, missions permanentes.
    \item $\square$ Situer le Cameroun dans les \textbf{organisations clés} : ONU, UA, CEMAC, CEEAC, CIRGL, OIF, Commonwealth, OCI.
    \item $\square$ Expliquer le lien \textbf{diplomatie $\leftrightarrow$ SND30} (diplomatie économique, attractivité, paix).
    \item $\square$ Analyser un \textbf{texte officiel} (discours, communiqué, traité) en extrayant le principe et l'intérêt national.
    \item $\square$ Argumenter à l'écrit : structure \emph{Thèse / Arguments (principes + faits) / Conclusion}.
    \item $\square$ Éviter le piège : confondre \emph{politique étrangère} (stratégie) et \emph{diplomatie} (instrument de mise en œuvre).
  \end{itemize}
\end{aretenir}

\findecours

\end{document}
